% Leçon 36 — Expression écrite : caractères complexes liés à la gestion des réseaux sociaux — Chinois Tle A — Cours signé OnBuch+
% Programme officiel MINESEC. Compiler avec : tectonic ZH36-ecrit-caracteres-complexes-lies-a-la-gestion-des-r.tex
\newcommand{\DOCMATIERE}{Chinois}
\newcommand{\DOCNIVEAU}{Tle A}
\newcommand{\DOCMODULE}{Module 5 : Mass médias et télécommunication}
\newcommand{\DOCLECON}{ZH36}
\newcommand{\DOCTITRE}{Leçon 36 — Expression écrite : caractères complexes liés à la gestion des réseaux sociaux}
% OnBuch+ — préambule « cours signés » ENRICHI (compilé avec tectonic / XeLaTeX)
% Système visuel « Pop » d'OnBuch+ : fond crème (jamais blanc), bordures encre
% épaisses, ombres « dures » décalées, titres Archivo Black en capitales.
% Version MATHÉMATIQUES : graphes (pgfplots/tikz), boîtes pédagogiques
% (définition, propriété, méthode, attention, le savais-tu, exercice résolu,
% exercice, TP, bilan), page de garde et signature OnBuch+ ancrée en bas.
%
% Macros attendues dans main.tex AVANT \input{preamble} :
%   \DOCMATIERE  \DOCNIVEAU  \DOCMODULE  \DOCLECON (numéro)  \DOCTITRE
\documentclass[11pt]{article}

\usepackage{fontspec}
\usepackage[french]{babel} % français = langue principale ; le chinois via \zh{..} / textebox (xeCJK)
\usepackage{xeCJK}
\usepackage{pifont} % \ding{51} \ding{55} utilisés par les modèles
\usepackage[a4paper,margin=2cm,top=3.4cm,bottom=2.8cm,headheight=1.4cm,footskip=1.3cm]{geometry}
\usepackage{amsmath,amssymb,amsfonts}
\usepackage{mathtools} % \xmapsto, \coloneqq... souvent utilisés par les modèles
\usepackage{cancel} % simplifications barrées (bilans de forces)
% \abs{z} (module, valeur absolue) et \norm{u} (norme), utilisés par les modèles
\providecommand{\abs}{}\let\abs\relax\DeclarePairedDelimiter{\abs}{\lvert}{\rvert}
\providecommand{\norm}{}\let\norm\relax\DeclarePairedDelimiter{\norm}{\lVert}{\rVert}
\usepackage[table]{xcolor} % [table] : fournit aussi \rowcolors (alternance de lignes)
\usepackage{pagecolor}
\usepackage{tikz}
\usetikzlibrary{arrows.meta,positioning,calc,shapes.geometric,shapes.misc,shapes.arrows,decorations.pathmorphing,decorations.pathreplacing,decorations.markings,patterns,shadows,mindmap,trees,fit,backgrounds,matrix,intersections,angles,quotes}
\usepackage{pgfplots}
\usepackage[version=4]{mhchem} % équations nucléaires \ce{^{235}_{92}U}
\usepackage[european,siunitx]{circuitikz} % schémas électriques
\pgfplotsset{compat=1.18}
\usepgfplotslibrary{fillbetween,groupplots} % aires entre courbes (calcul intégral)
\usepackage{enumitem}
\usepackage{fancyhdr}
% [most] charge toutes les bibliothèques tcolorbox (listings, minted, raster,
% documentation...) et gonfle inutilement la mémoire TeX sur un document long
% (~300 boîtes) : on ne charge que ce qui est réellement utilisé.
\usepackage[skins,breakable]{tcolorbox}
\usepackage{graphicx}
\usepackage{colortbl}
\usepackage{booktabs}
\usepackage{tabularx}
\usepackage{diagbox} % case d'en-tête diagonale des tableaux à double entrée
% Colonnes extensibles centrée / gauche / droite (C, L, R), souvent utilisées par les modèles
\newcolumntype{C}{>{\centering\arraybackslash}X}
\newcolumntype{L}{>{\raggedright\arraybackslash}X}
\newcolumntype{R}{>{\raggedleft\arraybackslash}X}
\usepackage{array}
\usepackage{multicol}
\usepackage{multirow}
\usepackage{siunitx}
% parse-numbers=false : accepte aussi \SI{2,5 \times 10^{-3}}{...}
\sisetup{locale=FR,output-decimal-marker={,},group-minimum-digits=4,parse-numbers=false}
\DeclareSIUnit\ohm{\ensuremath{\symup{\Omega}}} % U+2126 absent de la police
\DeclareSIUnit\division{div} % réglages d'oscilloscope (ms/div, V/div)
\DeclareSIUnit\tr{tr} % tours (vitesse de rotation, ex. \SI{1500}{\tr\per\minute})
\DeclareSIUnit\molar{M} % concentration molaire (ex. \SI{3}{\molar}, \SI{1}{\micro\molar})
\DeclareSIUnit\an{an} % année (le modèle confond parfois avec \year, un compteur TeX) — ex. \SI{5}{\kilo\meter\per\an}
\DeclareSIUnit\FCFA{FCFA} % franc CFA (ex. \SI{500}{\FCFA\per\kilo\gram})
\DeclareSIUnit\degreeBrix{\textdegree Brix} % teneur en sucre d'un sirop/jus (réfractométrie)
\DeclareSIUnit\month{mois} % le modèle utilise parfois \month, un compteur TeX, comme unité de durée
\usepackage{qrcode}
% \degree : commande fréquemment utilisée par les modèles pour un angle ou une
% température en dehors de \SI{}{} (non fournie par les paquets chargés).
\providecommand{\degree}{^{\circ}}
% \qed (fin de démonstration), utilisé par les modèles sans amsthm
\providecommand{\qed}{\hfill$\square$}
% \barre : double barre des valeurs interdites dans un tableau de variations (tkz-tab)
\providecommand{\barre}{\ensuremath{\|}}
% environnement proof (amsthm non chargé) utilisé par les modèles
\ifdefined\proof\else\newenvironment{proof}[1][Démonstration]{\par\noindent\textit{#1.}\ }{\qed\par}\fi
\usepackage{lastpage}
\usepackage{hyperref}
\hypersetup{hidelinks}

% ============================================================
% Palette « Pop » OnBuch+ — mêmes hex que la classe P (plus_ui.dart)
% ============================================================
\definecolor{poporange}{HTML}{F59321}
\definecolor{popdark}{HTML}{E07A0C}
\definecolor{popcream}{HTML}{FFF6E8}
\definecolor{popink}{HTML}{1C1714}
\definecolor{poppurple}{HTML}{7A5AE0}
\definecolor{popgreen}{HTML}{2E9E6A}
\definecolor{poppink}{HTML}{E0457B}
\definecolor{popblue}{HTML}{2F7DE1}
\definecolor{popgold}{HTML}{C98A12}
\definecolor{popmuted}{HTML}{8A7F76}
\definecolor{popline}{HTML}{EADFD2}
\definecolor{popgray}{HTML}{8A7F76}
\colorlet{popgrey}{popgray}
% Alias tolérants (le modèle nomme parfois les couleurs différemment)
\colorlet{popwhite}{white}
\colorlet{popblack}{popink}
\colorlet{popred}{poppink}
\colorlet{popyellow}{popgold}
% Teintes claires (fonds de tableaux/encadrés) — le modèle nomme parfois
% les couleurs avec un suffixe L (ex. popblueL) sans les avoir définies.
\colorlet{poporangeL}{poporange!20}
\colorlet{popdarkL}{popdark!20}
\colorlet{poppurpleL}{poppurple!20}
\colorlet{popgreenL}{popgreen!20}
\colorlet{poppinkL}{poppink!20}
\colorlet{popblueL}{popblue!20}
\colorlet{popgoldL}{popgold!20}
\colorlet{popredL}{popred!20}
\colorlet{popgrayL}{popgray!20}
% Teintes claires pour fonds de tableaux / zones
\colorlet{popblueL}{popblue!10!white}
\colorlet{poporangeL}{poporange!14!white}
\colorlet{popgreenL}{popgreen!12!white}
\colorlet{poppurpleL}{poppurple!12!white}
\colorlet{poppinkL}{poppink!12!white}
\colorlet{popgoldL}{popgold!14!white}

\pagecolor{popcream}

% ============================================================
% Typographie
% ============================================================
\defaultfontfeatures{Ligatures=TeX}
\setmainfont{PlusJakartaSans-Medium.ttf}[Path=../../../fonts/,BoldFont=PlusJakartaSans-Bold.ttf]
% Symboles, API et emojis absents de la police principale : repli sur DejaVu Sans (caractères actifs)
\newfontface\symfont{DejaVuSans.ttf}[Path=../../../fonts/]
\makeatletter
\newcommand\symactive[1]{\catcode#1=13 \begingroup\lccode`\~=#1 \lowercase{\endgroup\def~}{{\symfont\char#1}}}
\newcommand\symblank[1]{\catcode#1=13 \begingroup\lccode`\~=#1 \lowercase{\endgroup\def~}{}}
\newcommand\symhyph[1]{\catcode#1=13 \begingroup\lccode`\~=#1 \lowercase{\endgroup\def~}{\mbox{-}}}
\newcount\symc
\newcommand\symrange[2]{\symc=#1 \@whilenum\symc<#2 \do{\expandafter\symactive\expandafter{\the\symc}\advance\symc by 1 }}
\newcommand\symblankrange[2]{\symc=#1 \@whilenum\symc<#2 \do{\expandafter\symblank\expandafter{\the\symc}\advance\symc by 1 }}
\makeatother
\symrange{"0250}{"02FF}\symactive{"2713}\symactive{"2714}\symactive{"2717}\symactive{"2718}\symactive{"2610}\symactive{"2611}\symactive{"2612}
\symrange{"2600}{"27BF}
\catcode"2705=13 \begingroup\lccode`\~="2705 \lowercase{\endgroup\def~}{{\symfont\char"2713}}\symblank{"26BD}\symblank{"26A1}
\symrange{"2460}{"257F}\symactive{"223C}\symactive{"2248}\symactive{"2260}
\symactive{"01F8}\symactive{"01F9}\symactive{"1E3E}\symactive{"1E3F}
\symhyph{"2011}\symhyph{"2E1A}\symblankrange{"1F300}{"1FAFF}\symblank{"FE0F}
% Caractères chinois : WenQuanYi Zen Hei (sous-ensemble GB2312 + pinyin), gras/italique simulés
\setCJKmainfont{WenQuanYiZenHei-subset.ttf}[Path=../../../fonts/,AutoFakeBold=3,AutoFakeSlant=0.2]
\xeCJKsetup{CJKspace=false}
% Pinyin : voyelles à caron (ǎ ǐ ǒ ǔ ǖ ǘ ǚ ǜ) absentes de la police principale -> caractères actifs
\newfontface\pyfont{WenQuanYiZenHei-subset.ttf}[Path=../../../fonts/]
\newcommand\pyactive[1]{\catcode#1=13 \begingroup\lccode`\~=#1 \lowercase{\endgroup\def~}{{\pyfont\char#1}}}
\pyactive{"01CD}\pyactive{"01CE}\pyactive{"01CF}\pyactive{"01D0}\pyactive{"01D1}\pyactive{"01D2}\pyactive{"01D3}\pyactive{"01D4}\pyactive{"01D5}\pyactive{"01D6}\pyactive{"01D7}\pyactive{"01D8}\pyactive{"01D9}\pyactive{"01DA}\pyactive{"01DB}\pyactive{"01DC}
% Maths en sans-serif (Fira Math) avec chiffres et lettres droites de Plus
% Jakarta, pour que les formules se fondent dans le texte.
\usepackage[math-style=ISO,bold-style=ISO]{unicode-math}
\setmathfont{FiraMath-Regular.otf}[Path=../../../fonts/]
\setmathfont{PlusJakartaSans-Medium.ttf}[Path=../../../fonts/,range={up/num,up/{Latin,latin},\mathup}]
\usepackage{icomma} % virgule décimale sans espace en mode math
% Caractères mathématiques Unicode absents des polices de texte (Plus Jakarta,
% Archivo Black) : rendus par leur équivalent mathématique, en texte comme en
% formule — sinon ils s'affichent en carré vide (ex. « Arithmétique dans ℤ »).
\usepackage{newunicodechar}
\newunicodechar{ℝ}{\ensuremath{\mathbb{R}}}
\newunicodechar{ℤ}{\ensuremath{\mathbb{Z}}}
\newunicodechar{ℂ}{\ensuremath{\mathbb{C}}}
\newunicodechar{ℕ}{\ensuremath{\mathbb{N}}}
\newunicodechar{ℚ}{\ensuremath{\mathbb{Q}}}
\newunicodechar{𝔻}{\ensuremath{\mathbb{D}}}
\newunicodechar{α}{\ensuremath{\alpha}}
\newunicodechar{β}{\ensuremath{\beta}}
\newunicodechar{γ}{\ensuremath{\gamma}}
\newunicodechar{δ}{\ensuremath{\delta}}
\newunicodechar{ε}{\ensuremath{\varepsilon}}
\newunicodechar{θ}{\ensuremath{\theta}}
\newunicodechar{λ}{\ensuremath{\lambda}}
\newunicodechar{μ}{\ensuremath{\mu}}
\newunicodechar{σ}{\ensuremath{\sigma}}
\newunicodechar{τ}{\ensuremath{\tau}}
\newunicodechar{φ}{\ensuremath{\varphi}}
\newunicodechar{ψ}{\ensuremath{\psi}}
\newunicodechar{ω}{\ensuremath{\omega}}
\newunicodechar{Σ}{\ensuremath{\Sigma}}
% majuscules produites par \MakeUppercase dans les titres (α-aminés -> Α)
\newunicodechar{Α}{\ensuremath{\alpha}}
\newunicodechar{Β}{\ensuremath{\beta}}
\newunicodechar{Γ}{\ensuremath{\Gamma}}
\newunicodechar{Δ}{\ensuremath{\Delta}}
\newunicodechar{Ε}{\ensuremath{\varepsilon}}
\newunicodechar{Θ}{\ensuremath{\Theta}}
\newunicodechar{Λ}{\ensuremath{\Lambda}}
\newunicodechar{Μ}{\ensuremath{\mu}}
\newunicodechar{Τ}{\ensuremath{\tau}}
\newunicodechar{Φ}{\ensuremath{\Phi}}
\newunicodechar{Ψ}{\ensuremath{\Psi}}
\newunicodechar{Ω}{\ensuremath{\Omega}}
\newunicodechar{↦}{\ensuremath{\mapsto}}
\newunicodechar{∈}{\ensuremath{\in}}
\newunicodechar{∉}{\ensuremath{\notin}}
\newunicodechar{∧}{\ensuremath{\wedge}}
\newunicodechar{⇔}{\ensuremath{\Leftrightarrow}}
\newunicodechar{⟺}{\ensuremath{\Longleftrightarrow}}
\newunicodechar{⇒}{\ensuremath{\Rightarrow}}
\newunicodechar{⟹}{\ensuremath{\Longrightarrow}}
\newunicodechar{∘}{\ensuremath{\circ}}
\newunicodechar{∪}{\ensuremath{\cup}}
\newunicodechar{∩}{\ensuremath{\cap}}
\newunicodechar{≡}{\ensuremath{\equiv}}
\newunicodechar{⊂}{\ensuremath{\subset}}
\newunicodechar{⊥}{\ensuremath{\perp}}
\newunicodechar{∃}{\ensuremath{\exists}}
\newunicodechar{∀}{\ensuremath{\forall}}
\newunicodechar{∛}{\ensuremath{\sqrt[3]{\,}}}
\newunicodechar{∜}{\ensuremath{\sqrt[4]{\,}}}
\newunicodechar{⃗}{\ensuremath{{}^{\to}}}
\newunicodechar{ⁿ}{\ifmmode^{n}\else\textsuperscript{\ensuremath{n}}\fi}
\newunicodechar{ˣ}{\ifmmode^{x}\else\textsuperscript{\ensuremath{x}}\fi}
\newunicodechar{ᵇ}{\ifmmode^{b}\else\textsuperscript{\ensuremath{b}}\fi}
\newunicodechar{ʳ}{\ifmmode^{r}\else\textsuperscript{\ensuremath{r}}\fi}
\newunicodechar{ᵘ}{\ifmmode^{u}\else\textsuperscript{\ensuremath{u}}\fi}
\newunicodechar{ʸ}{\ifmmode^{y}\else\textsuperscript{\ensuremath{y}}\fi}
\newunicodechar{ᵃ}{\ifmmode^{a}\else\textsuperscript{\ensuremath{a}}\fi}
\newunicodechar{ᵉ}{\ifmmode^{e}\else\textsuperscript{\ensuremath{e}}\fi}
\newunicodechar{ᵏ}{\ifmmode^{k}\else\textsuperscript{\ensuremath{k}}\fi}
\newunicodechar{ᵖ}{\ifmmode^{p}\else\textsuperscript{\ensuremath{p}}\fi}
\newunicodechar{ᴺ}{\ifmmode^{N}\else\textsuperscript{\ensuremath{N}}\fi}
\newunicodechar{ᶜ}{\ifmmode^{c}\else\textsuperscript{\ensuremath{c}}\fi}
\newunicodechar{ʰ}{\ifmmode^{h}\else\textsuperscript{\ensuremath{h}}\fi}
\newunicodechar{ᵠ}{\ifmmode^{q}\else\textsuperscript{\ensuremath{q}}\fi}
\newunicodechar{ₙ}{\ifmmode_{n}\else\textsubscript{\ensuremath{n}}\fi}
\newunicodechar{ₐ}{\ifmmode_{a}\else\textsubscript{\ensuremath{a}}\fi}
\newunicodechar{ᵢ}{\ifmmode_{i}\else\textsubscript{\ensuremath{i}}\fi}
\newunicodechar{ᵣ}{\ifmmode_{r}\else\textsubscript{\ensuremath{r}}\fi}
\newunicodechar{ₖ}{\ifmmode_{k}\else\textsubscript{\ensuremath{k}}\fi}
\newunicodechar{ᵦ}{\ifmmode_{\beta}\else\textsubscript{\ensuremath{\beta}}\fi}
\newunicodechar{ₓ}{\ifmmode_{x}\else\textsubscript{\ensuremath{x}}\fi}
\newfontfamily\popdisplay{ArchivoBlack-Regular.ttf}[Path=../../../fonts/]
\newfontfamily\poplogo{SpaceGrotesk-Bold.ttf}[Path=../../../fonts/]
\newfontfamily\popsemi{PlusJakartaSans-SemiBold.ttf}[Path=../../../fonts/]
\newfontfamily\popxbold{PlusJakartaSans-ExtraBold.ttf}[Path=../../../fonts/]
\newfontfamily\popxboldls{PlusJakartaSans-ExtraBold.ttf}[Path=../../../fonts/,LetterSpace=3.0]

\setlist[enumerate]{leftmargin=1.7em,itemsep=5pt,topsep=4pt}
\setlist[itemize]{leftmargin=1.7em,itemsep=4pt,topsep=4pt,label={\color{poporange}\textbullet}}
\setlength{\parindent}{0pt}
\setlength{\parskip}{5pt}
\renewcommand{\arraystretch}{1.25}

% pgfplots : style Pop par défaut
\pgfplotsset{
  every axis/.append style={axis line style={popink,line width=1pt},tick style={popink},
    grid style={popline,line width=0.5pt},label style={font=\small},tick label style={font=\footnotesize}},
  popcurve/.style={poporange,line width=1.8pt,no markers,smooth},
}
% Même style disponible dans un \draw TikZ simple (hors pgfplots)
\tikzset{popcurve/.style={poporange,line width=1.8pt}}
\tikzset{popgrid/.style={popline,very thin}}
\colorlet{popcurve}{poporange} % le nom du style est parfois utilisé comme couleur

% ============================================================
% Pill / squiggle
% ============================================================
\newcommand{\poppill}[3]{%
  \tikz[baseline=-0.55ex]\node[draw=popink,line width=1.2pt,rounded corners=2.6pt,%
    fill=#2,inner xsep=6.5pt,inner ysep=3pt]{%
    \popxboldls\fontsize{7}{7}\selectfont\color{#3}\MakeUppercase{#1}};%
}
\newcommand{\popsquiggle}[1]{%
  \tikz[baseline=-0.4ex]{
    \draw[#1,line width=2.6pt,line cap=round]
      plot [smooth] coordinates {(0,0.09) (0.34,0.01) (0.68,0.21) (1.02,0.01) (1.36,0.10)};
  }%
}
% Mot-clé surligné (marqueur orange sous le texte)
\newcommand{\cle}[1]{{\popxbold\color{popdark}#1}}

% ============================================================
% En-tête / pied
% ============================================================
\pagestyle{fancy}
\fancyhf{}
\renewcommand{\headrulewidth}{0pt}
\renewcommand{\footrulewidth}{0pt}
\lhead{\raisebox{-0.15ex}{\poplogo\fontsize{13}{13}\selectfont\color{popink}OnBuch\color{poporange}+}}
\rhead{\poppill{\DOCMATIERE\ \textbullet\ \DOCNIVEAU\ \textbullet\ Leçon \DOCLECON}{white}{popink}}
\lfoot{\popxbold\fontsize{7.6}{7.6}\selectfont\color{popmuted}\MakeUppercase{Cours signé OnBuch+}\ \textbullet\ {\popsemi\DOCTITRE}}
\rfoot{\tikz[baseline=-0.6ex]\node[rounded corners=8.5pt,fill=popink,minimum height=17pt,minimum width=17pt,inner xsep=5pt,inner sep=0]{\popxbold\fontsize{8}{8}\selectfont\color{popcream}\thepage\,/\,\pageref*{LastPage}};}

% ============================================================
% Titres
% ============================================================
\newcommand{\chaptitre}[2]{%
  \begin{tcolorbox}[enhanced,colback=poporange,colframe=popink,boxrule=2.6pt,arc=11pt,
      drop shadow=popink,left=15pt,right=15pt,top=12pt,bottom=11pt,before skip=3pt,after skip=18pt]
    \poppill{#2}{popink}{popcream}\par\vspace{9pt}
    {\raggedright\hyphenpenalty=10000\exhyphenpenalty=10000\popdisplay\fontsize{22}{24}\selectfont\color{popcream}\MakeUppercase{#1}\par}
  \end{tcolorbox}
}
% Section numérotée : pastille orange avec le numéro + titre Archivo
\newcounter{popsec}
\newcommand{\coursec}[1]{%
  \stepcounter{popsec}\setcounter{popsub}{0}%
  \par\vspace{16pt}\noindent
  \begin{minipage}{\linewidth}
  \tikz[baseline=-0.7ex]\node[draw=popink,line width=1.6pt,fill=poporange,rounded corners=4pt,minimum size=22pt,inner sep=0]{\popdisplay\fontsize{12}{12}\selectfont\color{popcream}\thepopsec};%
  \hspace{8pt}{\popdisplay\fontsize{15}{17}\selectfont\color{popink}\MakeUppercase{#1}}\par
  \vspace{4pt}\noindent{\color{poporange}\rule{36pt}{3.6pt}}
  \end{minipage}\par\nopagebreak\vspace{8pt}
}
\newcounter{popsub}
\newcommand{\courssub}[1]{%
  \stepcounter{popsub}%
  \par\vspace{9pt}\noindent{\popxbold\fontsize{11.5}{13}\selectfont\color{popink}{\color{poporange}\thepopsec.\thepopsub}\hspace{6pt}#1}\par\nopagebreak\vspace{4pt}
}

% ============================================================
% Boîtes pédagogiques — même recette : fond blanc, bordure épaisse colorée,
% ombre dure encre, badge plein. Titre optionnel [#1].
% ============================================================
\tcbset{popbase/.style={breakable,enhanced,colback=white,boxrule=2.3pt,arc=9pt,
  shadow={3pt}{-3pt}{0pt}{popink},left=14pt,right=14pt,top=11pt,bottom=11pt,before skip=11pt,after skip=15pt,
  fontupper=\popsemi\fontsize{10.6}{14.5}\selectfont\color{popink},
  fontlower=\popsemi\fontsize{10.6}{14.5}\selectfont\color{popink}}}
% Coche dessinée (le glyphe ✓ n'existe pas dans les polices du document)
\newcommand{\popcheck}[1]{\tikz[baseline=-0.5ex]\draw[#1,line width=1.6pt,line cap=round,line join=round](-0.13,0.0)--(-0.04,-0.09)--(0.13,0.1);}
\providecommand{\coche}{\popcheck{popgreen}}
\providecommand{\resultat}[1]{\par\smallskip\noindent\fcolorbox{popgreen}{popgreenL}{\parbox{\dimexpr\linewidth-2\fboxsep-2\fboxrule\relax}{\textbf{Résultat.} #1}}\par\smallskip}
\newcommand{\popboxhead}[3]{\poppill{#1}{#2}{white}\ifx\relax#3\relax\else\hspace{6pt}{\popxbold\fontsize{10.6}{12}\selectfont\color{popink}#3}\fi\par\vspace{7pt}}

\newtcolorbox{aretenir}[1][]{popbase,colframe=popgreen,before upper={\popboxhead{À retenir}{popgreen}{#1}}}
% Texte en chinois (document de lecture, dialogue, extrait) : cadre
% d'étude. Un titre optionnel (source, auteur) s'affiche dans l'en-tête.
\newtcolorbox{textebox}[1][]{popbase,colframe=popblue,before upper={\popboxhead{文本 · Texte}{popblue}{#1}},after upper={}}
% Marques ✓ / ✗ (forme correcte / fautive) souvent utilisées par les modèles
\providecommand{\cross}{\textcolor{poppink}{\ensuremath{\times}}}
\providecommand{\xmark}{\cross}\providecommand{\crossmark}{\cross}\providecommand{\cmark}{\ensuremath{\checkmark}}
% Ligne de réponse / justification à compléter (inventée par les modèles)
\providecommand{\justification}[1]{\par\noindent\textit{Justification~:} \dotfill\par}
% \textipa (package tipa non chargé) : la police ne couvre pas l'API -> simple passage
\providecommand{\textipa}[1]{#1}
% Soulignement (ulem non chargé) : \uline, \ul
\providecommand{\uline}[1]{\underline{#1}}\providecommand{\ul}[1]{\underline{#1}}
% \glossaire{\item ...} inventée par les modèles : liste de glossaire
\providecommand{\glossaire}[1]{\begin{itemize}#1\end{itemize}}
% \alert (beamer) inventée par les modèles : texte mis en évidence
\providecommand{\alert}[1]{\textbf{\textcolor{poppink}{#1}}}
% variantes ASCII d'ordinaux écrites en macro
\providecommand{\iere}{\textsuperscript{re}}\providecommand{\ieme}{\textsuperscript{e}}\providecommand{\ere}{\textsuperscript{re}}\providecommand{\eme}{\textsuperscript{e}}\providecommand{\er}{\textsuperscript{er}}\providecommand{\ie}{\textsuperscript{e}}
% Ordinaux français écrits en macro par les modèles : 1\ière, 3\ième
\newcommand{\ière}{\textsuperscript{re}}\newcommand{\ième}{\textsuperscript{e}}
% \exemple (inventée par les modèles) : étiquette « Exemple : »
\providecommand{\exemple}{\textbf{Exemple~:}\ }
% \square utilisable aussi hors mode math (cases à cocher)
\AtBeginDocument{\let\squareMath\square\renewcommand{\square}{\ensuremath{\squareMath}}}
% Mot ou phrase en chinois à l'intérieur d'un paragraphe français
\newcommand{\zh}[1]{\ifmmode\text{#1}\else{#1}\fi}
\let\eng\zh \let\esp\zh \let\all\zh % alias tolérés
% flèches d'intonation inventées par les modèles
\providecommand{\textsearrow}{}\renewcommand{\textsearrow}{\ensuremath{\searrow}}\providecommand{\textnearrow}{}\renewcommand{\textnearrow}{\ensuremath{\nearrow}}
% \fr{..} : mot français (inventé par les modèles)
\providecommand{\fr}[1]{\textit{#1}}
% zhbox : encadré de texte chinois inventé par les modèles
\newenvironment{zhbox}{\begin{quote}}{\end{quote}}
% \courssubsub : titre de niveau 3 inventé par les modèles
\providecommand{\courssubsub}[1]{\par\medskip\noindent\textbf{#1}\par\smallskip}
% \zhbf : texte chinois en gras inventé par les modèles
\providecommand{\zhbf}[1]{\textbf{#1}}
% \vs : « versus » inventé par les modèles
\providecommand{\vs}{\textit{vs}\ }
% \blank : case à compléter inventée par les modèles
\providecommand{\blank}[1][]{\underline{\hspace{1.6em}}}
\newtcolorbox{exemplebox}[1][]{popbase,colframe=poporange,before upper={\popboxhead{Exemple}{poporange}{#1}}}
\newtcolorbox{definition}[1][]{popbase,colframe=popblue,before upper={\popboxhead{Définition}{popblue}{#1}}}
\newtcolorbox{propriete}[1][]{popbase,colframe=poppurple,before upper={\popboxhead{Propriété}{poppurple}{#1}}}
\newtcolorbox{methode}[1][]{popbase,colframe=popgold,before upper={\popboxhead{Méthode}{popgold}{#1}}}
\newtcolorbox{attention}[1][]{popbase,colframe=poppink,before upper={\popboxhead{Attention}{poppink}{#1}}}
\newtcolorbox{experience}[1][]{popbase,colframe=popblue,colback=popblueL,before upper={\popboxhead{Expérience}{popblue}{#1}}}
% « Le savais-tu ? » : carte violette pleine (contexte, culture, Cameroun)
\newtcolorbox{savaistu}[1][]{popbase,colback=poppurple,colframe=popink,
  fontupper=\popsemi\fontsize{10.4}{14.2}\selectfont\color{white},
  before upper={\poppill{Le savais-tu\,?}{white}{poppurple}\ifx\relax#1\relax\else\hspace{6pt}{\popxbold\color{white}#1}\fi\par\vspace{7pt}}}
% Exercice résolu : énoncé en haut, \tcblower puis la solution sur fond crème
\newcounter{popexo}\newcounter{popexr}
\newtcolorbox{exoresolu}[1][]{popbase,colframe=poporange,colbacklower=popcream,
  segmentation style={popink,line width=1.2pt,dashed},
  before upper={\stepcounter{popexr}\popboxhead{Exercice résolu \thepopexr}{poporange}{#1}},
  before lower={\poppill{Solution}{popink}{popcream}\par\vspace{6pt}}}
% Exercice (non résolu en place) — la correction est dans la partie Corrigés
\newtcolorbox{exercice}[1][]{popbase,colframe=popink,
  before upper={\stepcounter{popexo}\popboxhead{Exercice \thepopexo}{popink}{#1}}}
% Corrigé
\newtcolorbox{corrige}[1][]{popbase,colframe=popgreen,colback=popgreenL,
  before upper={\popboxhead{Corrigé}{popgreen}{#1}}}
% Objectifs / prérequis / bilan
\newtcolorbox{objectifs}{popbase,colframe=popink,colback=poporangeL,
  before upper={\popboxhead{Objectifs de la leçon}{popink}{}}}
\newtcolorbox{prerequis}{popbase,colframe=popink,colback=popblueL,
  before upper={\popboxhead{Prérequis}{popblue}{}}}
\newtcolorbox{bilan}{popbase,colframe=popink,colback=white,boxrule=2.8pt,
  before upper={\popboxhead{Fiche bilan}{poporange}{L'essentiel à connaître pour l'examen}}}
% Figure encadrée avec légende
\newtcolorbox{popfigure}[1][]{enhanced,breakable=false,colback=white,colframe=popline,boxrule=1.4pt,arc=8pt,
  left=10pt,right=10pt,top=10pt,bottom=8pt,before skip=10pt,after skip=12pt,halign=center,
  lower separated=false,halign lower=center,
  fontlower=\popsemi\fontsize{9}{11.5}\selectfont\color{popmuted},
  #1}
\newcommand{\legende}[1]{\par\vspace{6pt}{\popsemi\fontsize{9}{11.5}\selectfont\color{popmuted}#1}}

% ============================================================
% Signature OnBuch+
% ============================================================
\newcommand{\poplogoinline}[1]{{\poplogo\fontsize{#1}{#1}\selectfont\color{popink}OnBuch\color{poporange}+}}
\newcommand{\signatureonbuch}{%
  \begin{tcolorbox}[enhanced,colback=popink,colframe=popink,boxrule=2.2pt,arc=10pt,
      drop shadow=poporange,left=14pt,right=14pt,top=10pt,bottom=10pt,before skip=0pt,after skip=0pt]
    \tikz[baseline=-0.7ex]\node[circle,fill=popgreen,draw=popcream,line width=1.3pt,minimum size=22pt,inner sep=0]{\popcheck{white}};%
    \hspace{9pt}\begin{minipage}[c]{0.62\linewidth}
      {\popxbold\fontsize{12}{13}\selectfont\color{popcream}Signé\ \textbullet\ {\poplogo OnBuch\color{poporange}+}}\par\vspace{2pt}
      {\raggedright\popsemi\fontsize{8}{10.5}\selectfont\color{popline}\MakeUppercase{Équipe pédagogique OnBuch+}\\\MakeUppercase{Conforme au programme officiel MINESEC}\par}
    \end{minipage}\hfill
    {\poplogo\fontsize{22}{22}\selectfont\color{popcream}OnBuch\color{poporange}+}
  \end{tcolorbox}%
}

% ============================================================
% Page de garde
% #1 titre · #2 matière · #3 niveau · #4 module · #5 n° leçon · #6 durée ·
% #7 description / objectifs (texte ou itemize)
% ============================================================
\newcommand{\pagegarde}[7]{%
  \thispagestyle{empty}
  \newgeometry{a4paper,margin=1.7cm,top=1.6cm,bottom=1.4cm}%
  \begin{tikzpicture}[remember picture,overlay]
    \fill[poporange!18] ([xshift=-3.2cm,yshift=-3.6cm]current page.north east) circle (4.6cm);
    \fill[poppurple!14] ([xshift=2.4cm,yshift=7.6cm]current page.south west) circle (3.8cm);
  \end{tikzpicture}%
  {\poplogo\fontsize{30}{30}\selectfont\color{popink}OnBuch\color{poporange}+}\hfill
  \raisebox{0.6ex}{\poppill{Cours signé \textbullet\ Premium}{popink}{popcream}}\par
  \vspace{1pt}\popsquiggle{poppurple}\par
  \vspace{8pt}
  \begin{tcolorbox}[enhanced,colback=poporange,colframe=popink,boxrule=2.8pt,arc=13pt,
      drop shadow=popink,left=18pt,right=18pt,top=12pt,bottom=13pt,after skip=10pt]
    \poppill{#2 \textbullet\ #3}{popink}{popcream}\hspace{5pt}\poppill{#4}{white}{popink}\par\vspace{8pt}
    {\popdisplay\fontsize{14}{14}\selectfont\color{popink}LEÇON #5}\par\vspace{4pt}
    {\raggedright\hyphenpenalty=10000\exhyphenpenalty=10000\popdisplay\fontsize{25}{27}\selectfont\color{popcream}\MakeUppercase{#1}\par}
    \vspace{8pt}
    {\popxbold\fontsize{9.5}{9.5}\selectfont\color{popink}\MakeUppercase{Durée conseillée : #6}}
  \end{tcolorbox}
  \begin{tcolorbox}[enhanced,colback=white,colframe=popink,boxrule=2.2pt,arc=10pt,
      drop shadow=popink,left=16pt,right=16pt,top=10pt,bottom=10pt,after skip=8pt]
    \poppill{Dans cette leçon}{poppurple}{white}\par\vspace{5pt}
    \popsemi\fontsize{9.3}{12.6}\selectfont\color{popink}#7
  \end{tcolorbox}
  \noindent\begin{minipage}[t]{0.487\linewidth}
    \begin{tcolorbox}[enhanced,colback=popgreen,colframe=popink,boxrule=2.2pt,arc=10pt,
        drop shadow=popink,left=11pt,right=11pt,top=10pt,bottom=10pt,height=3.1cm,valign=center]
      \tcbox[colback=white,colframe=popink,boxrule=1.3pt,arc=5pt,boxsep=0pt,nobeforeafter,
        left=3pt,right=3pt,top=3pt,bottom=3pt,valign=center]{\qrcode[height=1.9cm]{https://play.google.com/store/apps/details?id=com.luvvix.onbuch&hl=fr}}%
      \hspace{8pt}\begin{minipage}[c]{0.5\linewidth}
        {\popdisplay\fontsize{11}{12.5}\selectfont\color{white}\MakeUppercase{Télécharge OnBuch}}\par\vspace{4pt}
        {\raggedright\popsemi\fontsize{8.4}{11}\selectfont\color{white}Gratuit \textbullet\ Google Play\\Tuteur IA, annales, concours, résultats\par}
      \end{minipage}
    \end{tcolorbox}
  \end{minipage}\hfill
  \begin{minipage}[t]{0.487\linewidth}
    \begin{tcolorbox}[enhanced,colback=poppurple,colframe=popink,boxrule=2.2pt,arc=10pt,
        drop shadow=popink,left=11pt,right=11pt,top=10pt,bottom=10pt,height=3.1cm,valign=center]
      \tikz[baseline=-0.7ex]\node[circle,fill=white,draw=popink,line width=1.3pt,minimum size=1.6cm,inner sep=0]{\popdisplay\fontsize{24}{24}\selectfont\color{poppurple}+};%
      \hspace{8pt}\begin{minipage}[c]{0.6\linewidth}
        {\popdisplay\fontsize{11}{12.5}\selectfont\color{white}\MakeUppercase{Passe à OnBuch+}}\par\vspace{4pt}
        {\raggedright\popsemi\fontsize{8.4}{11}\selectfont\color{white}Tous les cours signés, Tuteur IA illimité, fiches PDF — onglet OnBuch+ de l'app\par}
      \end{minipage}
    \end{tcolorbox}
  \end{minipage}
  \vspace{8pt}
  \popcontenu
  \vfill
  \signatureonbuch
  \restoregeometry
  \newpage
}

% Rangée « Ce cours contient » (page de garde)
\newcommand{\poptile}[3]{% #1 couleur #2 gros texte #3 légende
  \begin{tcolorbox}[enhanced,colback=#1,colframe=popink,boxrule=2pt,arc=9pt,shadow={2.5pt}{-2.5pt}{0pt}{popink},
      left=6pt,right=6pt,top=9pt,bottom=9pt,halign=center,valign=center,height=1.75cm,width=\linewidth,nobeforeafter]
    {\popdisplay\fontsize{12.5}{13}\selectfont\color{white}\MakeUppercase{#2}}\par\vspace{3pt}
    {\centering\popsemi\fontsize{7.8}{9.5}\selectfont\color{white}#3\par}
  \end{tcolorbox}}
\newcommand{\popcontenu}{%
  \par\noindent{\popxboldls\fontsize{7.5}{7.5}\selectfont\color{popmuted}CE COURS SIGNÉ CONTIENT}\par\vspace{7pt}
  \noindent\begin{minipage}[t]{0.235\linewidth}\poptile{popblue}{Cours}{complet, illustré, pas à pas}\end{minipage}\hfill
  \begin{minipage}[t]{0.235\linewidth}\poptile{poporange}{Méthodes}{et exercices résolus}\end{minipage}\hfill
  \begin{minipage}[t]{0.235\linewidth}\poptile{poppink}{Exercices}{type examen + corrigés}\end{minipage}\hfill
  \begin{minipage}[t]{0.235\linewidth}\poptile{popgreen}{Fiche bilan}{l'essentiel à retenir}\end{minipage}\par}

% Sommaire
\newtcolorbox{sommaire}{enhanced,colback=white,colframe=popink,boxrule=2.2pt,arc=10pt,
  drop shadow=popink,left=18pt,right=18pt,top=13pt,bottom=13pt,before skip=6pt,after skip=20pt,
  before upper={\poppill{Sommaire}{poppurple}{white}\par\vspace{9pt}\popsemi\fontsize{10.6}{14.5}\selectfont\color{popink}}}

% Signature de fin de cours — poussée tout en bas de la dernière page
\newcommand{\findecours}{%
  \par\vfill\nopagebreak
  \begin{center}\popsquiggle{poporange}\end{center}\vspace{4pt}
  \signatureonbuch
}
\AtBeginDocument{\def\llbracket{\mathopen{[\mkern-3.5mu[}}\def\rrbracket{\mathclose{]\mkern-3.5mu]}}} % intervalles d'entiers (glyphe ⟦ absent de la police)
% Glyphes absents de Fira Math : redéfinis à partir de glyphes disponibles
\AtBeginDocument{%
  \def\setminus{\mathbin{\backslash}}\let\smallsetminus\setminus
  \def\vdots{\mathord{\vbox{\baselineskip3.2pt\lineskiplimit0pt\kern4pt\hbox{.}\hbox{.}\hbox{.}}}}%
  \def\ddots{\mathinner{\mkern1mu\raise6.5pt\hbox{.}\mkern2mu\raise3.6pt\hbox{.}\mkern2mu\raise0.7pt\hbox{.}\mkern1mu}}%
  \def\triangle{\mathord{\scalebox{0.9}{$\symup{\Delta}$}}}%
  \def\nabla{\mathord{\raisebox{\depth}{\scalebox{1}[-1]{$\symup{\Delta}$}}}}%
  \def\checkmark{\popcheck{popdark}}%
  \def\star{\mathbin{\ast}}\def\bigstar{\mathord{\ast}}%
  \def\diamond{\mathbin{\scalebox{0.6}{$\Diamond$}}}%
  \def\bigcup{\mathop{\mathchoice{\raisebox{-0.3em}{\scalebox{1.7}{$\cup$}}}{\scalebox{1.25}{$\cup$}}{\cup}{\cup}}\displaylimits}%
  \def\bigcap{\mathop{\mathchoice{\raisebox{-0.3em}{\scalebox{1.7}{$\cap$}}}{\scalebox{1.25}{$\cap$}}{\cap}{\cap}}\displaylimits}%
  \def\lesseqqgtr{\mathrel{\lessgtr}}\def\bigoplus{\mathop{\oplus}\displaylimits}%
}
\providecommand{\ssi}{si et seulement si} % abréviation fréquente
\AtBeginDocument{\providecommand{\wideparen}{\overparen}} % arc de cercle

\begin{document}

\pagegarde{Leçon 36 — Expression écrite : caractères complexes liés à la gestion des réseaux sociaux}{Chinois}{Tle A}{Module 5 : Mass médias et télécommunication}{ZH36}{2 h}{Cette leçon d'expression écrite entraîne les élèves à produire un texte simple en chinois sur la gestion des réseaux sociaux : écriture rigoureuse des caractères complexes du thème (traits, radicaux, caractères proches), utilisation de connecteurs, et traduction dans les deux sens (thème et version).
\vspace{4pt}
\begin{multicols}{2}\small
\begin{itemize}
\item À la fin de cette leçon, je sais écrire correctement les caractères complexes liés aux réseaux sociaux (网, 络, 微, 信, 博, 频, 赞, 注, 删, 屏...)
\item je sais identifier les radicaux et respecter l'ordre des traits des caractères du thème
\item je sais distinguer les caractères graphiquement proches (微/徽, 博/搏, 络/路, 辨/辩)
\item je sais rédiger un texte simple et structuré (introduction, développement, conclusion) sur mes habitudes sur les réseaux sociaux
\item je sais utiliser les connecteurs utiles : 首先, 然后, 但是, 所以, 因为...所以...
\item je sais traduire des phrases simples du français vers le chinois (thème)
\end{itemize}\end{multicols}}

\begin{sommaire}
\begin{multicols}{2}
\begin{enumerate}
\item Lexique du thème : les réseaux sociaux
\item Écriture des caractères : radicaux et ordre des traits
\item Caractères proches : ne plus se tromper
\item Modèle de texte commenté et connecteurs
\item Thème et version : traduire dans les deux sens
\item Grille d'évaluation et entraînement à l'examen
\item Activité d'intégration
\item Méthodes et exercices résolus
\item Exercices
\item Corrigés des exercices
\item Fiche bilan
\end{enumerate}
\end{multicols}
\end{sommaire}

\begin{objectifs}
\begin{itemize}
\item À la fin de cette leçon, je sais écrire correctement les caractères complexes liés aux réseaux sociaux (网, 络, 微, 信, 博, 频, 赞, 注, 删, 屏...)
\item je sais identifier les radicaux et respecter l'ordre des traits des caractères du thème
\item je sais distinguer les caractères graphiquement proches (微/徽, 博/搏, 络/路, 辨/辩)
\item je sais rédiger un texte simple et structuré (introduction, développement, conclusion) sur mes habitudes sur les réseaux sociaux
\item je sais utiliser les connecteurs utiles : 首先, 然后, 但是, 所以, 因为...所以...
\item je sais traduire des phrases simples du français vers le chinois (thème)
\item je sais traduire un court texte du chinois vers le français (version)
\item je sais m'auto-évaluer à l'aide d'une grille (exactitude des caractères, grammaire, connecteurs, respect du sujet)
\end{itemize}
\end{objectifs}

\begin{prerequis}
\begin{itemize}
\item Vocabulaire de base des médias et d'Internet (手机, 电脑, 上网, 发消息) vu en Première
\item Ordre général des traits et notion de radical (règles d'écriture : haut-bas, gauche-droite)
\item Structures de base : 是, 有, 用 + nom, 喜欢 + verbe
\item Connecteurs simples : 和, 也, 因为...所以...
\item Nombres, expressions de fréquence (每天, 经常, 有时候)
\end{itemize}
\end{prerequis}

\newpage

\coursec{Lexique du thème : les réseaux sociaux}

À Douala comme à Yaoundé, la plupart des élèves de Terminale consultent WhatsApp, Facebook ou TikTok plusieurs fois par jour : au marché Mokolo, dans les taxis de Biyem-Assi, à la sortie du lycée, les écrans brillent partout. En Chine, les jeunes utilisent WeChat \zh{微信} et Weibo \zh{微博} pour discuter, publier des photos et suivre l'actualité. Un ami chinois te demande par message : \zh{你经常用什么社交媒体？} \emph{(Nǐ jīngcháng yòng shénme shèjiāo méitǐ ?)} « Quels réseaux sociaux utilises-tu souvent ? » Sauras-tu lui répondre par écrit, avec des caractères correctement tracés, et lui décrire tes habitudes en ligne ? C'est tout l'enjeu de cette leçon. Avant d'écrire, il faut d'abord posséder le \cle{lexique} exact du thème : c'est l'objet de cette première section.

\courssub{Observer d'abord : un message authentique}

\begin{textebox}[Un message d'un ami chinois]
\zh{我每天用微信跟朋友聊天，也经常在微博上发照片。}\\
\emph{Wǒ měitiān yòng Wēixìn gēn péngyou liáotiān, yě jīngcháng zài Wēibó shàng fā zhàopiàn.}\\
« Je discute chaque jour avec mes amis sur WeChat, et je publie souvent des photos sur Weibo. »\\[4pt]
\zh{我的账号有很多粉丝，他们都给我的视频点赞、评论。}\\
\emph{Wǒ de zhànghào yǒu hěn duō fěnsī, tāmen dōu gěi wǒ de shìpín diǎnzàn, pínglùn.}\\
« Mon compte a beaucoup d'abonnés ; ils likent et commentent tous mes vidéos. »\\[4pt]
\zh{网络很方便，但是也很危险：你要保护好自己的隐私和密码。}\\
\emph{Wǎngluò hěn fāngbiàn, dànshì yě hěn wēixiǎn : nǐ yào bǎohù hǎo zìjǐ de yǐnsī hé mìmǎ.}\\
« Internet est très pratique, mais aussi très dangereux : tu dois bien protéger ta vie privée et ton mot de passe. »
\end{textebox}

\textbf{Questions de compréhension}
\begin{enumerate}
\item Quelles sont les deux plateformes mentionnées dans le message ?
\item Que font les \zh{粉丝} du narrateur ?
\item Quel contraste est exprimé par \zh{但是} (dànshì) ?
\end{enumerate}

Ce court message contient presque tout le lexique de la leçon. Observons-le mot à mot, puis organisons-le.

\courssub{Les mots-clés : plateformes et contenus}

\begin{definition}[社交媒体 shèjiāo méitǐ]
\cle{社交媒体} (shèjiāo méitǐ) désigne les \textbf{médias sociaux / réseaux sociaux} : les plateformes d'échange en ligne où l'on publie, commente et partage. Le mot se construit ainsi : \zh{社交} (shèjiāo, « relations sociales ») + \zh{媒体} (méitǐ, « médias »). On le retrouve dans \zh{社交媒体平台} (shèjiāo méitǐ píngtái, « plateforme de médias sociaux »).
\end{definition}

\begin{center}
\begin{tabularx}{\linewidth}{|X|X|X|X|}
\hline
\rowcolor{popblueL} Caractères & Pinyin & Nature & Traduction \\
\hline
网络 & wǎngluò & nom (cl. 个) & réseau, Internet \\
\hline
社交媒体 & shèjiāo méitǐ & nom (cl. 个) & réseaux sociaux, médias sociaux \\
\hline
微信 & Wēixìn & nom propre & WeChat (litt. « micro-message ») \\
\hline
微博 & Wēibó & nom propre & Weibo (litt. « micro-blog ») \\
\hline
视频 & shìpín & nom (cl. 个, 条) & vidéo \\
\hline
照片 & zhàopiàn & nom (cl. 张) & photo \\
\hline
消息 & xiāoxi & nom (cl. 条) & message, nouvelle \\
\hline
朋友圈 & péngyouquān & nom (cl. 个) & « Moments » (fil d'actualité WeChat) \\
\hline
\end{tabularx}
\end{center}

\begin{exemplebox}[Les mots-clés en contexte]
\zh{现在的年轻人都喜欢上网。} \emph{(Xiànzài de niánqīngrén dōu xǐhuan shàngwǎng.)} « Les jeunes d'aujourd'hui aiment tous se connecter à Internet. »\\
\zh{这个视频很有意思，你看看吧！} \emph{(Zhège shìpín hěn yǒuyìsi, nǐ kànkan ba !)} « Cette vidéo est très intéressante, regarde-la ! »\\
\zh{她在朋友圈发了一张照片。} \emph{(Tā zài péngyouquān fā le yì zhāng zhàopiàn.)} « Elle a publié une photo dans ses Moments. »
\end{exemplebox}

\courssub{Les verbes d'action : publier, liker, partager}

Ces verbes sont le cœur de la leçon : ce sont eux qui permettent de décrire des habitudes en ligne. Plusieurs sont des \cle{verbes séparables} (离合词) : verbe + objet figé (ex. \zh{点赞} = 点 + 赞), ce qui permet d'insérer un classificateur ou un complément entre les deux caractères.

\begin{center}
\begin{tabularx}{\linewidth}{|X|X|X|X|}
\hline
\rowcolor{popblueL} Caractères & Pinyin & Nature & Traduction \\
\hline
发 & fā & verbe & envoyer, publier, poster \\
\hline
点赞 & diǎnzàn & verbe séparable (点+赞) & liker (litt. « cliquer-louer ») \\
\hline
评论 & pínglùn & verbe / nom (cl. 条) & commenter ; commentaire \\
\hline
关注 & guānzhù & verbe séparable (关+注) & suivre, s'abonner à \\
\hline
分享 & fēnxiǎng & verbe & partager \\
\hline
删除 & shānchú & verbe & supprimer, effacer \\
\hline
屏蔽 & píngbì & verbe & bloquer, masquer \\
\hline
加 & jiā & verbe & ajouter \\
\hline
刷 & shuā & verbe & faire défiler (écran), balayer \\
\hline
\end{tabularx}
\end{center}

\begin{exemplebox}[Exemples traduits]
\zh{请关注我的账号。} \emph{(Qǐng guānzhù wǒ de zhànghào.)} « Veuillez suivre mon compte. »\\
\zh{他删掉了那条评论。} \emph{(Tā shānchú diào le nà tiáo pínglùn.)} « Il a supprimé ce commentaire. »\\
\zh{我把这个视频分享给我的网友了。} \emph{(Wǒ bǎ zhège shìpín fēnxiǎng gěi wǒ de wǎngyǒu le.)} « J'ai partagé cette vidéo avec mes amis en ligne. »\\
\zh{别给他点赞，先屏蔽他吧！} \emph{(Bié gěi tā diǎnzàn, xiān píngbì tā ba !)} « Ne like pas son contenu, bloque-le d'abord ! »\\
\zh{我每天刷半小时视频。} \emph{(Wǒ měitiān shuā bàn xiǎoshí shìpín.)} « Je fais défiler des vidéos pendant une demi-heure chaque jour. »
\end{exemplebox}

\begin{attention}[发 fā : un verbe à multiples visages]
\zh{发} (fā) a de nombreux emplois : \zh{发消息} (envoyer un message), \zh{发照片} (publier une photo), \zh{发烧} (fāshāo, avoir de la fièvre), \zh{发工资} (fā gōngzī, verser le salaire). Le sens dépend \textbf{toujours du complément} qui suit. Erreur fréquente des francophones : traduire mécaniquement \zh{发} par « envoyer ». Dans \zh{发微博}, il faut comprendre « publier sur Weibo », pas « envoyer Weibo ». Vérifie toujours le nom placé après le verbe avant de traduire.
\end{attention}

\courssub{Les personnes et les notions associées}

\begin{center}
\begin{tabularx}{\linewidth}{|X|X|X|X|}
\hline
\rowcolor{popblueL} Caractères & Pinyin & Nature & Traduction \\
\hline
网友 & wǎngyǒu & nom (cl. 个, 位) & internaute, ami en ligne \\
\hline
粉丝 & fěnsī & nom (cl. 个, 位) & fans, abonnés \\
\hline
账号 & zhànghào & nom (cl. 个) & compte (d'utilisateur) \\
\hline
密码 & mìmǎ & nom (cl. 个, 组) & mot de passe \\
\hline
隐私 & yǐnsī & nom (cl. 个) & vie privée, données personnelles \\
\hline
好友 & hǎoyǒu & nom (cl. 个, 位) & ami (contact sur un réseau) \\
\hline
\end{tabularx}
\end{center}

\begin{savaistu}[粉丝 : des vermicelles devenus des fans]
\zh{粉丝} (fěnsī) signifie à l'origine « vermicelles de haricot » (appelés \emph{fěnsī} ou \emph{fěntiáo}), un aliment très courant en Chine ! Mais comme sa prononciation \emph{fěnsī} ressemble à l'anglais \emph{fans}, les Chinois l'ont choisi comme emprunt phonétique pour désigner les admirateurs et les abonnés. C'est un exemple célèbre de mot emprunté par le son : le sens premier (vermicelles) n'a rien à voir avec le nouveau sens (fans). De même, \zh{微博} (Wēibó) imite le son de l'anglais \emph{micro-blog}.
\end{savaistu}

\courssub{Adjectifs et jugements : pratique ou dangereux ?}

\begin{center}
\begin{tabularx}{\linewidth}{|X|X|X|X|}
\hline
\rowcolor{popblueL} Caractères & Pinyin & Nature & Traduction \\
\hline
方便 & fāngbiàn & adjectif & pratique, commode \\
\hline
有用 & yǒuyòng & adjectif & utile \\
\hline
危险 & wēixiǎn & adjectif & dangereux \\
\hline
浪费时间 & làngfèi shíjiān & locution verbale & perdre du temps, une perte de temps \\
\hline
有意思 & yǒuyìsi & adjectif & intéressant, amusant \\
\hline
上瘾 & shàngyǐn & verbe / adjectif & addictif, rendre accro \\
\hline
\end{tabularx}
\end{center}

\zh{微信很方便，可是刷视频太浪费时间了，容易上瘾。} \emph{(Wēixìn hěn fāngbiàn, kěshì shuā shìpín tài làngfèi shíjiān le, róngyì shàngyǐn.)} « WeChat est très pratique, mais faire défiler des vidéos fait vraiment perdre trop de temps, c'est facile d'en devenir accro. »

\courssub{Point de grammaire : les verbes séparables (离合词)}

\begin{propriete}[Structure et emploi des verbes séparables]
Certains verbes chinois sont formés de \textbf{Verbe + Objet} (ex. \zh{点赞} 点+赞, \zh{上网} 上+网, \zh{关注} 关+注, \zh{睡觉} 睡+觉). Ils se comportent comme un seul verbe mais peuvent être « séparés » :
\begin{enumerate}
\item \textbf{Insertion d'un classificateur / quantificateur :} \zh{点个赞} (diǎn ge zàn, « mettre un like »), \zh{上了两个小时网} (shàng le liǎng ge xiǎoshí wǎng, « être resté en ligne deux heures »).
\item \textbf{Insertion de \zh{了} (aspect accompli) entre les deux :} \zh{点了赞} (diǎn le zàn), \zh{关了注} (guān le zhù).
\item \textbf{Reduplication du premier caractère (tentative) :} \zh{点点赞} (diǎn dian zàn, « liker un peu »).
\end{enumerate}
\emph{Attention :} \zh{评论}, \zh{删除}, \zh{屏蔽}, \zh{分享}, \zh{加} ne sont \textbf{pas} séparables (ce sont des verbes simples ou V+V).
\end{propriete}

\courssub{Carte mentale du lexique}

\begin{popfigure}
\begin{tikzpicture}[mindmap, grow cyclic, every node/.style={concept, align=center, font=\small}, concept color=popblue!40, level 1/.append style={sibling angle=90, level distance=3.5cm}]
\node[concept color=popblue!60, align=center]{\zh{社交媒体}\\ shèjiāo méitǐ}
  child[concept color=popgreen!45] { node[align=center]{Plateformes\\ \zh{微信} \zh{微博}\\ \zh{朋友圈} \zh{网络}} }
  child[concept color=poporange!50] { node[align=center]{Actions\\ \zh{发} \zh{点赞} \zh{评论}\\ \zh{关注} \zh{分享}\\ \zh{删除} \zh{屏蔽}\\ \zh{加} \zh{刷}} }
  child[concept color=poppurple!45] { node[align=center]{Personnes\\ \zh{网友} \zh{粉丝}\\ \zh{好友}} }
  child[concept color=poppink!50] { node[align=center]{Notions / Jugements\\ \zh{账号} \zh{密码} \zh{隐私}\\ \zh{方便} \zh{危险}\\ \zh{浪费时间} \zh{上瘾}} };
\end{tikzpicture}
\legende{Carte mentale du lexique des réseaux sociaux : quatre familles de mots autour de \zh{社交媒体}.}
\end{popfigure}

\courssub{Les collocations fréquentes : verbe + complément}

En chinois comme en français, certains verbes « appellent » certains compléments. Ces associations fixes (\cle{collocations}) sont indispensables pour écrire des phrases naturelles. On indique le classificateur le plus courant entre parenthèses.

\begin{center}
\begin{tabularx}{\linewidth}{|X|X|X|}
\hline
\rowcolor{popblueL} Verbe + complément & Pinyin & Traduction \\
\hline
发 + 消息 (条) & fā xiāoxi & envoyer un message \\
\hline
发 + 照片 (张) & fā zhàopiàn & publier une photo \\
\hline
发 + 朋友圈 (条/个) & fā péngyouquān & publier un « Moment » \\
\hline
点 + 赞 (个) & diǎn zàn & liker \\
\hline
删除 + 照片 (张) / 评论 (条) & shānchú zhàopiàn / pínglùn & supprimer une photo / un commentaire \\
\hline
刷 + 视频 (个/条) & shuā shìpín & faire défiler des vidéos \\
\hline
加 + 好友 (个) & jiā hǎoyǒu & ajouter un ami (contact) \\
\hline
上 + 网 & shàng wǎng & se connecter à Internet \\
\hline
关注 + 账号 (个) & guānzhù zhànghào & suivre un compte \\
\hline
保护 + 隐私 & bǎohù yǐnsī & protéger sa vie privée \\
\hline
修改 + 密码 & xiūgǎi mìmǎ & modifier son mot de passe \\
\hline
\end{tabularx}
\end{center}

Remarque la structure de \zh{点赞} et \zh{上网} : ce sont des verbes dits « séparables » (verbe + objet figé). On peut insérer un élément entre les deux caractères : \zh{点个赞} (diǎn ge zàn, « mettre un petit like »), \zh{上了两个小时网} (shàng le liǎng ge xiǎoshí wǎng, « être resté deux heures en ligne »).

\begin{exoresolu}[Construire des phrases avec les collocations]
Énoncé : traduis en chinois les deux phrases suivantes en utilisant les collocations du tableau. (a) « J'ai ajouté un nouvel ami sur WeChat. » (b) « Elle passe trop de temps à faire défiler des vidéos. »
\tcblower
Solution détaillée :\\
(a) Le verbe « ajouter un ami » se dit \zh{加好友}. « Sur WeChat » est un complément de lieu introduit par \zh{在}, placé \textbf{avant} le verbe (ordre chinois : Sujet + 在 + lieu + Verbe). L'accompli se marque avec \zh{了}. On utilise le classificateur \zh{个} pour \zh{好友}. On obtient : \zh{我在微信上加了一个新好友。} \emph{(Wǒ zài Wēixìn shàng jiā le yí ge xīn hǎoyǒu.)}\\
(b) « Faire défiler des vidéos » = \zh{刷视频}. « Passer du temps à » se rend par \zh{花} (huā) + durée + \zh{在} + Verbe (ou directement Verbe). Structure : Sujet + 花 + durée + 刷视频. On obtient : \zh{她每天花很多时间刷视频。} \emph{(Tā měitiān huā hěn duō shíjiān shuā shìpín.)} « Elle passe beaucoup de temps chaque jour à faire défiler des vidéos. »
\end{exoresolu}

\begin{attention}[Pièges fréquents des francophones]
\begin{itemize}
\item Ne traduis pas « aimer une publication » par \zh{喜欢} (xǐhuan, sentiment) : pour « liker » (action technique), le mot exact est \zh{点赞}. \zh{我喜欢这张照片} = « J'aime cette photo » (sentiment) ; \zh{我给这张照片点了赞} = « J'ai liké cette photo » (action).
\item « Suivre » quelqu'un sur un réseau = \zh{关注}, jamais \zh{跟} (gēn, « suivre physiquement / avec »). \zh{我关注了他} = « Je me suis abonné à son compte ».
\item Attention à l'homophonie / paronymie : \zh{密码} (mìmǎ, mot de passe, secret) ne se confond pas avec \zh{号码} (hàomǎ, numéro de téléphone, numéro de série). Retiens : \zh{账号} (compte) contient \zh{号}, \zh{密码} (code secret) commence par \zh{密} « secret ».
\item \zh{发} vs \zh{发送} (fāsòng) : \zh{发送} est plus formel, « envoyer/transmettre » (ex. email, données). \zh{发} suffit pour les réseaux sociaux.
\item Classificateurs : n'oublie pas \zh{条} pour \zh{消息}, \zh{评论}, \zh{视频} (longs) et \zh{张} pour \zh{照片}.
\end{itemize}
\end{attention}

\begin{exercice}[À toi de jouer]
\begin{enumerate}
\item Complète avec le verbe qui convient (\zh{发}, \zh{点赞}, \zh{关注}, \zh{删除}, \zh{屏蔽}, \zh{加}) :\\
\zh{请\_\_\_\_我的微博账号。} \\
\zh{这条评论不好，我要\_\_\_\_它。} \\
\zh{我每天都给朋友的照片\_\_\_\_。} \\
\zh{这个陌生人总发广告，我把他\_\_\_\_了。}
\item Traduis en français : \zh{网络很方便，但是你要保护好自己的隐私和密码。}
\item Traduis en chinois : « J'ai partagé une photo avec mes amis en ligne. » (Utilise la structure \zh{把}).
\item Traduis en chinois : « Il a liké ma vidéo et a écrit un commentaire. » (Utilise \zh{了} pour les deux actions successives).
\end{enumerate}
\end{exercice}

\begin{corrige}[Exercice « À toi de jouer »]
1. \zh{请关注我的微博账号。} (Qǐng guānzhù wǒ de Wēibó zhànghào.) — Veuillez suivre mon compte Weibo.\\
\zh{这条评论不好，我要删除它。} (Zhè tiáo pínglùn bù hǎo, wǒ yào shānchú tā.) — Ce commentaire n'est pas bien, je vais le supprimer.\\
\zh{我每天都给朋友的照片点赞。} (Wǒ měitiān dōu gěi péngyou de zhàopiàn diǎnzàn.) — Chaque jour je like les photos de mes amis.\\
\zh{这个陌生人总发广告，我把他屏蔽了。} (Zhège mòshēngrén zǒng fā guǎnggào, wǒ bǎ tā píngbì le.) — Cet inconnu poste tout le temps de la pub, je l'ai bloqué.\\[6pt]
2. « Internet est très pratique, mais tu dois bien protéger ta vie privée et ton mot de passe. »\\[6pt]
3. \zh{我把一张照片分享给我的网友了。} \emph{(Wǒ bǎ yì zhāng zhàopiàn fēnxiǎng gěi wǒ de wǎngyǒu le.)} — Structure \zh{把} : Sujet + 把 + Objet (quantifié) + Verbe + Complément de résultat/direction.\\
\emph{Variante sans \zh{把} :} \zh{我分享了一张照片给我的网友。}\\[6pt]
4. \zh{他给我的视频点了赞，还评论了。} \emph{(Tā gěi wǒ de shìpín diǎn le zàn, hái pínglùn le.)} — Note la séparation de \zh{点赞} pour insérer \zh{了} : \zh{点 + 了 + 赞}. \zh{还} (hái) marque l'action additionnelle.
\end{corrige}

\begin{aretenir}[L'essentiel du lexique]
\begin{itemize}
\item \textbf{Plateformes :} \zh{网络}, \zh{社交媒体}, \zh{微信}, \zh{微博}, \zh{朋友圈}.
\item \textbf{Actions (verbes) :} \zh{发} (polyvalent), \zh{点赞}, \zh{评论}, \zh{关注}, \zh{分享}, \zh{删除}, \zh{屏蔽}, \zh{加}, \zh{刷}.
\item \textbf{Verbes séparables clés :} \zh{点赞} (点+赞), \zh{上网} (上+网), \zh{关注} (关+注) — permettent insertion de \zh{了}, classificateurs, reduplication.
\item \textbf{Personnes et notions :} \zh{网友}, \zh{粉丝} (emprunt phonétique \emph{fans}), \zh{账号}, \zh{密码}, \zh{隐私}, \zh{好友}.
\item \textbf{Jugements :} \zh{方便}, \zh{有用}, \zh{危险}, \zh{浪费时间}, \zh{上瘾}.
\item \textbf{Collocations incontournables :} \zh{发消息/照片/朋友圈}, \zh{点赞}, \zh{刷视频}, \zh{加好友}, \zh{上网}, \zh{关注账号}, \zh{保护隐私}, \zh{修改密码}.
\item \textbf{Classificateurs :} \zh{条} (message, commentaire, vidéo), \zh{张} (photo), \zh{个} (compte, ami, fan, vidéo courte).
\end{itemize}
\end{aretenir}

\coursec{Écriture des caractères : radicaux et ordre des traits}

Dans la section précédente, nous avons découvert le lexique des réseaux sociaux : \zh{网络} (wǎngluò, « réseau »), \zh{微博} (wēibó, « microblog »), \zh{评论} (pínglùn, « commentaire »), \zh{删除} (shānchú, « supprimer »), \zh{视频} (shìpín, « vidéo »). Vous savez maintenant les lire et les comprendre. Mais à l'examen, on vous demandera aussi de les \emph{écrire} correctement, trait par trait. Un caractère mal tracé — un trait oublié, un radical confondu — peut changer complètement le sens. Cette section vous donne les clés pour écrire ces caractères sans faute : les radicaux, l'ordre des traits et la décomposition en blocs.

\courssub{Observer d'abord : les radicaux racontent une histoire}

\begin{textebox}[Petit texte d'observation]
我每天都在网上看微博，写评论。\\
Wǒ měitiān dōu zài wǎngshang kàn wēibó, xiě pínglùn.\\
« Chaque jour, je lis des microblogs sur Internet et j'écris des commentaires. »\\[4pt]
昨天我删了一条不好的评论。\\
Zuótiān wǒ shān le yì tiáo bù hǎo de pínglùn.\\
« Hier, j'ai supprimé un mauvais commentaire. »\\[4pt]
网络把喀麦隆和中国连在一起。\\
Wǎngluò bǎ Kāmàilóng hé Zhōngguó lián zài yìqǐ.\\
« Internet relie le Cameroun et la Chine. »
\end{textebox}

Regardez bien les caractères que votre œil repère : \zh{网}, \zh{络}, \zh{微}, \zh{博}, \zh{评}, \zh{论}, \zh{删}. Chacun contient une petite pièce récurrente, le \cle{radical} (ou « clé », \zh{部首} bùshǒu), qui indique la famille de sens du caractère. Le reste du caractère donne souvent le son. Comprendre cette logique, c'est apprendre deux fois plus vite.

\begin{definition}[Radical (部首)]
Le \cle{radical} est la partie d'un caractère qui classe son sens. Il sert aussi à chercher le caractère dans un dictionnaire. Exemple : dans \zh{评} (píng, « commenter »), le radical est \zh{讠}, la « parole » : commenter, c'est d'abord parler.
\end{definition}

\courssub{Les radicaux clés du thème « réseaux sociaux »}

\begin{center}
\begin{tabularx}{\linewidth}{|X|X|X|X|}
\hline
\rowcolor{popblueL} Radical & Sens & Caractères du thème & Exemple traduit \\
\hline
纟 (3 traits) & soie, fil & 络 (luò), 红 (hóng) & 网络 wǎngluò « réseau » ; 网红 wǎnghóng « célébrité du web » \\
\hline
讠 (2 traits) & parole & 评 (píng), 论 (lùn), 语 (yǔ) & 评论 pínglùn « commentaire » ; 汉语 Hànyǔ « chinois » \\
\hline
彳 (3 traits) & pas, marcher & 微 (wēi), 很 (hěn) & 微博 wēibó « microblog » ; 很好 hěn hǎo « très bien » \\
\hline
十 (2 traits) & dix & 博 (bó) & 微博 wēibó ; 博士 bóshì « docteur » \\
\hline
页 (6 traits) & tête, page & 频 (pín), 题 (tí) & 视频 shìpín « vidéo » ; 问题 wèntí « question » \\
\hline
刂 (2 traits) & couteau & 删 (shān) & 删除 shānchú « supprimer » \\
\hline
\end{tabularx}
\end{center}

\begin{exemplebox}[Le sens caché dans le caractère]
\zh{网络的“络”有丝旁（纟）。}\\
\emph{Wǎngluò de “luò” yǒu sīpáng (纟).}\\
« Le 络 de 网络 porte la clé de la soie (纟). »\\
Image à retenir : le réseau est comme un \textbf{fil de soie} qui relie les gens entre eux, de Yaoundé à Pékin. De même, \zh{删} contient le couteau \zh{刂} : supprimer, c'est « couper » le message.
\end{exemplebox}

\begin{savaistu}[讠, la parole simplifiée]
Le radical \zh{讠} est la forme simplifiée de \zh{言} (yán, « parole »), qui existe encore comme caractère complet. Presque tous les caractères liés à la communication le contiennent : \zh{说} (shuō, « parler »), \zh{话} (huà, « parole »), \zh{评} (píng, « commenter »), \zh{论} (lùn, « discuter »), \zh{语} (yǔ, « langue »), \zh{读} (dú, « lire »), \zh{谢} (xiè, « remercier »). Quand vous voyez \zh{讠}, pensez : « ici, on parle ! »
\end{savaistu}

\courssub{Rappel : les quatre règles d'or de l'ordre des traits}

\begin{propriete}[Règles générales du tracé]
\begin{enumerate}
\item \textbf{De haut en bas} : \zh{言} se trace du trait supérieur vers le bas.
\item \textbf{De gauche à droite} : dans \zh{络}, \zh{评}, \zh{删}, on trace toujours la partie \textbf{gauche} avant la partie \textbf{droite}.
\item \textbf{L'horizontal avant le vertical} : dans \zh{十}, le trait horizontal — passe avant le vertical 丨.
\item \textbf{L'extérieur avant l'intérieur, et l'on ferme en dernier} : dans \zh{网}, on trace le cadre extérieur, puis l'intérieur ; le trait qui fermerait un enclos complet (comme dans \zh{国}) vient toujours en dernier.
\end{enumerate}
\end{propriete}

\begin{attention}[Erreur fréquente des francophones]
En français, on écrit les lettres dans l'ordre où elles sonnent. En chinois, on ne « suit pas le son » : on suit la \textbf{géométrie}. Ne commencez jamais un caractère gauche-droite par sa partie droite, même si c'est elle qui donne le son. \zh{评} commence toujours par \zh{讠}, jamais par \zh{平}.
\end{attention}

\courssub{Étude détaillée : 网 (6 traits)}

\zh{网} (wǎng, « filet, réseau ») est un caractère à structure semi-englobante : son radical est \zh{网} lui-même (dont la variante \zh{罒} apparaît en haut d'autres caractères). Ordre des traits : (1) le vertical gauche 丨, (2) le trait horizontal-vertical-crochet qui forme le haut et la droite, (3) le petit oblique intérieur gauche 丿, (4) le point intérieur 丶, (5) le second oblique intérieur, (6) le second point. Règle appliquée : \textbf{extérieur avant intérieur}.

\begin{popfigure}
\begin{center}
\begin{tikzpicture}[scale=0.9]
% Grille 米字格
\draw[popline, thin] (0,0) rectangle (3,3);
\draw[popline, dashed] (1.5,0) -- (1.5,3);
\draw[popline, dashed] (0,1.5) -- (3,1.5);
\draw[popline, dashed] (0,0) -- (3,3);
\draw[popline, dashed] (0,3) -- (3,0);
% Tracé stylisé de 网
\draw[popblue, very thick] (0.6,0.5) -- (0.6,2.6);
\draw[popblue, very thick] (0.6,2.6) -- (2.4,2.6) -- (2.4,0.5);
\draw[poporange, very thick] (0.95,2.0) -- (1.35,1.35);
\draw[poporange, very thick] (1.0,1.3) -- (1.4,1.55);
\draw[popgreen, very thick] (2.05,2.0) -- (1.65,1.35);
\draw[popgreen, very thick] (2.0,1.3) -- (1.6,1.55);
\node[popblue, font=\small] at (0.35,2.7) {1};
\node[popblue, font=\small] at (1.5,2.85) {2};
\node[poporange, font=\small] at (0.85,2.25) {3};
\node[poporange, font=\small] at (0.8,1.15) {4};
\node[popgreen, font=\small] at (2.2,2.25) {5};
\node[popgreen, font=\small] at (2.25,1.15) {6};
\end{tikzpicture}
\end{center}
\legende{Grille 米字格 (« grille en riz ») : tracé de 网 en 6 traits. Bleu : le cadre extérieur d'abord ; orange et vert : l'intérieur ensuite.}
\end{popfigure}

La grille \zh{米字格} (mǐzìgé, « grille du caractère riz ») est votre meilleure amie : ses lignes médianes et diagonales vous aident à \emph{centrer} chaque trait. Entraînez-vous toujours sur du papier quadrillé.

\courssub{Étude détaillée : 络 (9 traits)}

\zh{络} (luò) est un caractère \textbf{gauche-droite} : à gauche \zh{纟} (3 traits : deux pliés obliques puis un relevé horizontal), à droite \zh{各} (6 traits : l'oblique 丿, le plié horizontal-oblique, le petit oblique, puis la « bouche » \zh{口} en 3 traits). On trace \textbf{d'abord tout le 纟, ensuite tout le 各}, de haut en bas.

\begin{propriete}[Caractères gauche-droite]
Dans un caractère de structure gauche-droite (\zh{络}, \zh{评}, \zh{论}, \zh{删}, \zh{微}), on trace \textbf{toujours} la partie gauche en entier avant de commencer la partie droite. La partie gauche, souvent le radical, est étroite ; la partie droite est plus large. Proportion conseillée : environ un tiers à gauche, deux tiers à droite.
\end{propriete}

\courssub{Étude détaillée : 微 (13 traits)}

\zh{微} (wēi, « petit, micro- ») est le caractère le plus difficile de cette leçon. Il se décompose en \textbf{trois blocs}, tracés de gauche à droite :
\begin{itemize}
\item Bloc 1 (gauche) : \zh{彳} (3 traits) — deux obliques puis un vertical.
\item Bloc 2 (milieu) : \zh{山} + \zh{一} + \zh{几} (7 traits) — la « montagne » d'abord, de haut en bas.
\item Bloc 3 (droite) : \zh{攵} (4 traits) — le radical « frapper » : oblique, horizontal, puis le croisement final oblique-pointé.
\end{itemize}

\begin{popfigure}
\begin{center}
\begin{tikzpicture}[scale=1.0]
\node[draw=popblue, fill=popblueL, rounded corners, minimum width=1.1cm, minimum height=1.1cm, font=\Large] (a) at (0,0) {彳};
\node[draw=poporange, fill=poporangeL, rounded corners, minimum width=2.2cm, minimum height=1.1cm, font=\Large] (b) at (2.2,0) {山 一 几};
\node[draw=popgreen, fill=popgreenL, rounded corners, minimum width=1.1cm, minimum height=1.1cm, font=\Large] (c) at (4.4,0) {攵};
\draw[-{Stealth}, very thick, popdark] (a) -- (b) node[midway, above, font=\small] {1};
\draw[-{Stealth}, very thick, popdark] (b) -- (c) node[midway, above, font=\small] {2};
\node[font=\small, align=center, popmuted] at (0,-1.0) {3 traits\\(gauche)};
\node[font=\small, align=center, popmuted] at (2.2,-1.0) {7 traits\\(milieu)};
\node[font=\small, align=center, popmuted] at (4.4,-1.0) {3 traits\\(droite)};
\node[font=\large] at (2.2,1.2) {微 wēi = 13 traits};
\end{tikzpicture}
\end{center}
\legende{Décomposition de 微 en trois blocs colorés, tracés de gauche à droite : 彳, puis 山一几, puis 攵.}
\end{popfigure}

\courssub{Étude détaillée : 博 (12 traits) et 删 (7 traits)}

\zh{博} (bó, « vaste, érudit ») : à gauche \zh{十} (2 traits, horizontal avant vertical), à droite \zh{甫} + \zh{寸} (10 traits), de haut en bas. \textbf{Attention au point} \zh{丶} en haut à droite de \zh{甫} : c'est le trait le plus souvent oublié par les élèves ! Sans ce point, le caractère est faux.

\zh{删} (shān, « supprimer ») : à gauche \zh{册} (5 traits, le « registre » de lamelles de bambou), à droite le couteau \zh{刂} (2 traits : vertical court, puis long vertical-crochet). Règle : le couteau à droite se trace \textbf{en dernier}, comme si l'on « coupait » le caractère à la fin.

\begin{methode}[Mémoriser un caractère complexe en 4 étapes]
\begin{enumerate}
\item \textbf{Identifier le radical} : quelle famille de sens ? (\zh{评} → \zh{讠} parole).
\item \textbf{Décomposer en blocs} : gauche / milieu / droite, ou extérieur / intérieur (\zh{微} = 彳 + 山一几 + 攵).
\item \textbf{Écrire 5 fois en comptant les traits} à voix haute : « un, deux, trois… treize ! »
\item \textbf{Réécrire de mémoire}, puis vérifier : le radical est-il à sa place ? Aucun trait oublié (le point de \zh{博}) ?
\end{enumerate}
\end{methode}

\begin{attention}[Ne pas confondre 讠 et 氵]
\zh{讠} (parole, \textbf{2 traits} : point + plié horizontal-vertical-relevé) et \zh{氵} (eau, \textbf{3 traits} : deux points + relevé) se ressemblent pour un œil pressé. Mais \zh{评论} (pínglùn, « commenter ») n'a rien à voir avec \zh{汗} (hàn, « sueur ») ! Test rapide : si le mot parle de communication, c'est \zh{讠} ; s'il parle de liquide, c'est \zh{氵}. Comparez : \zh{语} (yǔ, langue) / \zh{江} (jiāng, fleuve) ; \zh{说} (shuō, parler) / \zh{洗} (xǐ, laver).
\end{attention}

\begin{center}
\begin{tabularx}{\linewidth}{|X|X|X|X|}
\hline
\rowcolor{popblueL} Caractère & Traits & Structure & Piège à éviter \\
\hline
网 wǎng & 6 & semi-englobant & tracer le cadre extérieur d'abord, l'intérieur ensuite \\
\hline
络 luò & 9 & gauche-droite & 纟 en 3 traits, pas en 2 ; 各 à droite ensuite \\
\hline
微 wēi & 13 & gauche-milieu-droite & ne pas oublier le 一 entre 山 et 几 \\
\hline
博 bó & 12 & gauche-droite & le point 丶 en haut à droite de 甫 \\
\hline
删 shān & 7 & gauche-droite & le couteau 刂 se trace en dernier \\
\hline
频 pín & 13 & gauche-droite & 页 à droite : la « tête » avec ses deux traits finaux \\
\hline
\end{tabularx}
\end{center}

\begin{exoresolu}[Décomposer et compter les traits]
Énoncé : décomposez le caractère \zh{评} (píng, « commenter ») en blocs, donnez son radical et comptez ses traits. Puis écrivez la phrase \zh{我写评论。} en respectant l'ordre gauche-droite.
\tcblower
Solution détaillée :
\begin{enumerate}
\item \zh{评} est un caractère gauche-droite : à gauche le radical \zh{讠} (parole, 2 traits : le point 丶 puis le plié horizontal-vertical-relevé), à droite \zh{平} (píng, « plat », 5 traits : horizontal, point, oblique, horizontal, vertical).
\item Nombre total : 2 + 5 = \textbf{7 traits}. Ordre : tout le \zh{讠} d'abord, puis tout le \zh{平}, de haut en bas.
\item Phrase : \zh{我写评论。} \emph{Wǒ xiě pínglùn.} « J'écris un commentaire. » Pour \zh{写} (5 traits), on trace le couvre-chef \zh{冖} (2 traits) avant la partie inférieure \zh{与} (3 traits) ; pour \zh{论} (6 traits), \zh{讠} avant \zh{仑}.
\end{enumerate}
Vérification : le radical \zh{讠} annonce bien le sens « parole », et \zh{平} donne le son píng. Caractère juste, sens juste.
\end{exoresolu}

\begin{exercice}[À vous de tracer]
\begin{enumerate}
\item Décomposez \zh{语} (yǔ) en blocs et comptez ses traits.
\item Écrivez \zh{微} cinq fois sur une grille 米字格 en comptant les traits à voix haute.
\item Recopiez la phrase \zh{我删了一条评论。} (Wǒ shān le yì tiáo pínglùn. « J'ai supprimé un commentaire. ») en respectant l'ordre des traits, puis traduisez-la.
\item Trouvez l'intrus et justifiez : \zh{评} — \zh{论} — \zh{河} — \zh{语}.
\end{enumerate}
\end{exercice}

\begin{corrige}[Exercice « À vous de tracer »]
\begin{enumerate}
\item \zh{语} = \zh{讠} (2 traits) + \zh{吾} (7 traits : \zh{五} 4 traits + \zh{口} 3 traits) = \textbf{9 traits}. Structure gauche-droite : \zh{讠} d'abord.
\item Vérification au comptage : 彳 (3) + 山一几 (7) + 攵 (3) = 13 traits, tracés de gauche à droite.
\item Traduction : « J'ai supprimé un commentaire. » Dans \zh{删}, le couteau \zh{刂} se trace en dernier ; dans \zh{条} (7 traits), le haut \zh{夂} avant le bas \zh{木}.
\item L'intrus est \zh{河} (hé, « rivière ») : son radical est \zh{氵} (eau), alors que \zh{评}, \zh{论}, \zh{语} portent tous \zh{讠} (parole), la famille de la communication.
\end{enumerate}
\end{corrige}

\begin{aretenir}[L'essentiel de la section]
\begin{itemize}
\item Le radical donne la famille de sens : \zh{纟} soie (\zh{络}), \zh{讠} parole (\zh{评}, \zh{论}, \zh{语}), \zh{彳} pas (\zh{微}), \zh{页} tête (\zh{频}), \zh{刂} couteau (\zh{删}).
\item Quatre règles d'or : haut avant bas, gauche avant droite, horizontal avant vertical, extérieur avant intérieur (et l'on ferme en dernier).
\item Décomposer en blocs avant d'écrire : \zh{微} = 彳 + 山一几 + 攵 (13 traits) ; \zh{博} = 十 + 甫 + 寸 (12 traits, sans oublier le point de \zh{甫}).
\item Méthode en 4 étapes : identifier le radical, décomposer, écrire 5 fois en comptant, réécrire de mémoire.
\end{itemize}
\end{aretenir}

\coursec{Caractères proches : ne plus se tromper}

Dans la section précédente, nous avons appris à écrire les caractères du thème en respectant les radicaux et l'ordre des traits. Mais connaître l'ordre des traits ne suffit pas : certains caractères se ressemblent tellement qu'un seul trait ou un seul composant les différencie. En chinois, comme en français où « côte » et « cote » n'ont pas le même sens, confondre deux caractères proches change complètement le sens de la phrase — et à l'examen, c'est une faute. Cette section vous donne une méthode rigoureuse pour ne plus jamais confondre les paires de caractères du thème des réseaux sociaux.

\courssub{Le principe général : le radical sémantique et le composant phonétique}

\begin{propriete}[Radical sémantique + composant phonétique]
La grande majorité des caractères chinois sont des caractères \cle{phonéto-sémantiques} : ils sont formés de deux parties.
\begin{itemize}
\item Le \cle{radical sémantique} indique la \textbf{famille de sens} : \zh{讠} (parole), \zh{扌} (main), \zh{贝} (argent), \zh{足} (pied), \zh{虫} (insecte/animal), \zh{巾} (tissu), \zh{纟} (soie, fil).
\item Le \cle{composant phonétique} suggère la \textbf{prononciation}, souvent de façon approximative.
\end{itemize}
Conséquence pratique : deux caractères qui partagent le même composant phonétique se prononcent souvent pareil ou presque (\zh{辨} et \zh{辩} se lisent tous deux \emph{biàn}) ; deux caractères qui se ressemblent visuellement peuvent se prononcer différemment si leurs radicaux diffèrent (\zh{络} \emph{luò} et \zh{路} \emph{lù} partagent le composant \zh{各}, mais pas le même radical).
\end{propriete}

\begin{attention}[La ressemblance visuelle ne garantit ni le son ni le sens]
Les francophones ont tendance à lire les caractères « globalement », comme des dessins. Erreur : \zh{络} et \zh{路} contiennent tous deux \zh{各}, mais l'un a le radical \zh{纟} (soie, relier) et l'autre \zh{足} (pied, route). Leurs prononciations diffèrent (\emph{luò} / \emph{lù}) et leurs sens n'ont rien à voir. Il faut donc toujours décomposer le caractère : \textbf{quelle partie donne le sens ? quelle partie donne le son ?}
\end{attention}

\courssub{Six paires de caractères proches du thème des réseaux sociaux}

\textbf{Paire 1 : 微 (wēi) ≠ 徽 (huī)}\par
\zh{微} signifie « petit, micro » : c'est le premier caractère de \zh{微信} (\emph{Wēixìn}, WeChat) et de \zh{微博} (\emph{Wēibó}, Weibo). \zh{徽} signifie « emblème, insigne » : \zh{徽章} (\emph{huīzhāng}, insigne). La différence se situe dans la partie inférieure : \zh{微} se termine à droite par \zh{攵} (radical « frapper »), tandis que \zh{徽} contient \zh{糸} (soie) en bas.

\begin{attention}[Erreur fréquente : écrire 徽信 au lieu de 微信]
Beaucoup d'élèves francophones écrivent \zh{徽信} au lieu de \zh{微信}. Phrase-mémo : pensez « \zh{微} = petit, comme un micro-message ». WeChat, c'est le « micro-message » : \zh{微信}. Le caractère \zh{徽} (emblème) ne sert pratiquement jamais au lycée ; si vous hésitez, c'est presque toujours \zh{微} qu'il faut écrire.
\end{attention}

\textbf{Paire 2 : 博 (bó) ≠ 搏 (bó)}\par
Même prononciation, sens très différents. \zh{博} signifie « vaste, étendu » : \zh{微博} (\emph{Wēibó}), littéralement « micro-vaste », la plateforme de microblogging. \zh{搏}, avec le radical de la main \zh{扌}, signifie « lutter, combattre » : \zh{搏斗} (\emph{bódòu}, lutter).

\begin{exemplebox}[Deux phrases à ne pas confondre]
\zh{我在微博上关注他。}\\
\emph{Wǒ zài Wēibó shang guānzhù tā.} « Je le suis sur Weibo. »\\[4pt]
\zh{他在搏斗。}\\
\emph{Tā zài bódòu.} « Il est en train de lutter. »
\end{exemplebox}

Phrase-mémo : « Weibo, c'est vaste (\zh{博}) comme le monde ; lutter (\zh{搏}), c'est avec la main (\zh{扌}). »

\textbf{Paire 3 : 络 (luò) ≠ 路 (lù)}\par
\zh{络}, avec le radical de la soie \zh{纟}, évoque les fils qui relient : \zh{网络} (\emph{wǎngluò}, réseau, Internet). \zh{路}, avec le radical du pied \zh{足}, signifie « route, chemin » : \zh{马路} (\emph{mǎlù}, route). Phrase-mémo : « le réseau (\zh{络}) est fait de fils (\zh{纟}) ; la route (\zh{路}) se parcourt avec les pieds (\zh{足}). »

\textbf{Paire 4 : 辨 (biàn) ≠ 辩 (biàn)}\par
Ces deux caractères se prononcent exactement pareil. \zh{辨}, avec \zh{刀} (couteau) au centre, signifie « distinguer, discerner » : \zh{辨别是非} (\emph{biànbié shìfēi}, distinguer le vrai du faux) — expression essentielle pour parler des fausses informations en ligne. \zh{辩}, avec \zh{讠} (parole) au centre, signifie « débattre, argumenter » : \zh{辩论} (\emph{biànlùn}, débat). Phrase-mémo : « débattre (\zh{辩}) se fait avec la parole (\zh{讠}) ; distinguer (\zh{辨}) se fait comme avec un couteau qui sépare. »

\textbf{Paire 5 : 账 (zhàng) ≠ 帐 (zhàng)}\par
\zh{账}, avec le radical \zh{贝} (coquillage, ancienne monnaie → argent), signifie « compte » : \zh{账号} (\emph{zhànghào}, numéro de compte, identifiant). \zh{帐}, avec le radical \zh{巾} (tissu), signifie « tente, moustiquaire » : \zh{帐篷} (\emph{zhàngpeng}, tente). Phrase-mémo : « un compte (\zh{账}), c'est de l'argent (\zh{贝}) ; une tente (\zh{帐}), c'est du tissu (\zh{巾}). »

\textbf{Paire 6 : 密 (mì) ≠ 蜜 (mì)}\par
\zh{密}, avec \zh{山} (montagne) en bas, signifie « secret, dense » : \zh{密码} (\emph{mìmǎ}, mot de passe). \zh{蜜}, avec \zh{虫} (insecte) en bas, signifie « miel » : \zh{蜂蜜} (\emph{fēngmì}, miel). Phrase-mémo : « le miel (\zh{蜜}) vient des abeilles (\zh{虫}) ; le mot de passe (\zh{密}) est caché comme derrière une montagne (\zh{山}). »

\courssub{Tableau récapitulatif des six paires}

\begin{center}
\begin{tabularx}{\linewidth}{|X|X|X|X|}
\hline
\rowcolor{popblueL} Caractère & Pinyin et sens & Mot du thème & Phrase-mémo \\
\hline
微 & wēi, petit & 微信 WeChat & micro-message \\
\hline
徽 & huī, emblème & 徽章 insigne & emblème avec soie 糸 \\
\hline
博 & bó, vaste & 微博 Weibo & vaste comme le monde \\
\hline
搏 & bó, lutter & 搏斗 lutter & lutter avec la main 扌 \\
\hline
络 & luò, relier & 网络 réseau & réseau de fils 纟 \\
\hline
路 & lù, route & 马路 route & route pour les pieds 足 \\
\hline
辨 & biàn, distinguer & 辨别是非 distinguer le vrai du faux & couteau qui sépare \\
\hline
辩 & biàn, débattre & 辩论 débat & débat avec la parole 讠 \\
\hline
账 & zhàng, compte & 账号 identifiant & compte et argent 贝 \\
\hline
帐 & zhàng, tente & 帐篷 tente & tente en tissu 巾 \\
\hline
密 & mì, secret & 密码 mot de passe & secret derrière la montagne 山 \\
\hline
蜜 & mì, miel & 蜂蜜 miel & miel des abeilles 虫 \\
\hline
\end{tabularx}
\end{center}

\begin{popfigure}
\begin{tikzpicture}[
  node distance=0.5cm,
  carac/.style={rectangle, draw=popblue, fill=popblueL, rounded corners, minimum width=2.6cm, minimum height=1cm, align=center, font=\large},
  memo/.style={rectangle, draw=poporange, fill=poporangeL, rounded corners, minimum width=3.6cm, minimum height=1cm, align=center, font=\small}
]
\node[carac, align=center] (a1){微 wēi\\ petit};
\node[carac, right=of a1, align=center] (b1){徽 huī\\ emblème};
\node[memo, right=of b1, align=center] (m1){micro-message\\ = 微信};
\node[carac, below=of a1, align=center] (a2){博 bó\\ vaste};
\node[carac, right=of a2, align=center] (b2){搏 bó\\ lutter};
\node[memo, right=of b2] (m2) {lutter = main 扌};
\node[carac, below=of a2, align=center] (a3){账 zhàng\\ compte};
\node[carac, right=of a3, align=center] (b3){帐 zhàng\\ tente};
\node[memo, right=of b3] (m3) {compte = argent 贝};
\draw[-{Stealth}, popdark] (a1) -- (b1);
\draw[-{Stealth}, popdark] (b1) -- (m1);
\draw[-{Stealth}, popdark] (a2) -- (b2);
\draw[-{Stealth}, popdark] (b2) -- (m2);
\draw[-{Stealth}, popdark] (a3) -- (b3);
\draw[-{Stealth}, popdark] (b3) -- (m3);
\end{tikzpicture}
\legende{Tableau comparatif : caractère A, caractère B, phrase-mémo.}
\end{popfigure}

\courssub{Méthode : comment distinguer deux caractères proches}

\begin{methode}[Distinguer deux caractères proches en quatre étapes]
\begin{enumerate}
\item \textbf{Décomposer} chaque caractère : repérer le radical sémantique et le composant phonétique.
\item \textbf{Associer le radical au sens} : \zh{贝} = argent → \zh{账} (compte) ; \zh{讠} = parole → \zh{辩} (débattre) ; \zh{扌} = main → \zh{搏} (lutter).
\item \textbf{Créer une phrase-mémo complète} qui utilise le bon caractère dans un contexte réel : \zh{我用账号登录微博。} (\emph{Wǒ yòng zhànghào dēnglù Wēibó.} « Je me connecte à Weibo avec mon identifiant. »)
\item \textbf{Vérifier la prononciation} : si les deux caractères se prononcent différemment (\zh{络} \emph{luò} / \zh{路} \emph{lù}), le son est un indice supplémentaire pour ne pas les confondre.
\end{enumerate}
\end{methode}

\begin{popfigure}
\begin{tikzpicture}[
  node distance=0.6cm and 1.2cm,
  quest/.style={rectangle, draw=poppurple, fill=poppurpleL, rounded corners, align=center, minimum width=4.2cm, minimum height=0.9cm},
  res/.style={rectangle, draw=popgreen, fill=popgreenL, rounded corners, align=center, minimum width=3.4cm, minimum height=0.9cm}
]
\node[quest, align=center] (q){Quel est le sens\\ recherché ?};
\node[quest, below left=of q] (q1) {lié à la parole ?};
\node[quest, below right=of q] (q2) {lié à la main ?};
\node[res, below=of q1, align=center] (r1){radical 讠\\ → 辩 (débattre)};
\node[res, below=of q2, align=center] (r2){radical 扌\\ → 搏 (lutter)};
\draw[-{Stealth}, popdark] (q) -- (q1);
\draw[-{Stealth}, popdark] (q) -- (q2);
\draw[-{Stealth}, popdark] (q1) -- (r1) node[midway, left, font=\small] {oui};
\draw[-{Stealth}, popdark] (q2) -- (r2) node[midway, right, font=\small] {oui};
\end{tikzpicture}
\legende{Arbre de décision : le sens guide le choix du radical.}
\end{popfigure}

\courssub{Exercice résolu et entraînement}

\begin{exoresolu}[Choisir le bon caractère]
Énoncé. Complétez chaque phrase avec le bon caractère de la paire indiquée, puis traduisez.\\
1. 我的微\_\_\_密码很复杂。（博 / 搏 ; 密 / 蜜）\\
2. 他用\_\_\_号登录网\_\_\_。（账 / 帐 ; 络 / 路）\\
3. 我们要学会\_\_\_别是非。（辨 / 辩）
\tcblower
Solution détaillée.\\
1. \zh{我的微博密码很复杂。} \emph{Wǒ de Wēibó mìmǎ hěn fùzá.} « Mon mot de passe Weibo est très complexe. » \zh{博} = vaste (pas de main \zh{扌}, donc pas \zh{搏}) ; \zh{密} = secret avec \zh{山} (pas d'insecte \zh{虫}, donc pas \zh{蜜}).\\
2. \zh{他用账号登录网络。} \emph{Tā yòng zhànghào dēnglù wǎngluò.} « Il se connecte au réseau avec son identifiant. » \zh{账} avec \zh{贝} (argent → compte) ; \zh{络} avec \zh{纟} (fils → réseau).\\
3. \zh{我们要学会辨别是非。} \emph{Wǒmen yào xuéhuì biànbié shìfēi.} « Nous devons apprendre à distinguer le vrai du faux. » \zh{辨} (distinguer) et non \zh{辩} (débattre) : il s'agit de séparer le vrai du faux, pas de discuter.
\end{exoresolu}

\begin{exercice}[À vous de jouer]
Traduisez en chinois en choisissant le bon caractère à chaque fois :\\
1. « J'ai oublié le mot de passe de mon compte WeChat. »\\
2. « Sur Internet, il faut distinguer le vrai du faux. »\\
3. « Il n'est pas en train de lutter, il débat avec ses amis. »
\end{exercice}

\begin{corrige}[Exercice : traduction]
1. \zh{我忘了我的微信账号的密码。} \emph{Wǒ wàng le wǒ de Wēixìn zhànghào de mìmǎ.} — \zh{微} (pas \zh{徽}), \zh{账} (pas \zh{帐}), \zh{密} (pas \zh{蜜}).\\
2. \zh{在网络上，我们要辨别是非。} \emph{Zài wǎngluò shang, wǒmen yào biànbié shìfēi.} — \zh{络} (pas \zh{路}), \zh{辨} (pas \zh{辩}).\\
3. \zh{他不是在搏斗，他是在和朋友们辩论。} \emph{Tā bú shì zài bódòu, tā shì zài hé péngyoumen biànlùn.} — \zh{搏} (main, lutter) et \zh{辩} (parole, débattre). Remarquez la négation \zh{不是} : on nie une identification (« ce n'est pas qu'il lutte »), pas une simple action.
\end{corrige}

\begin{savaistu}[À l'examen, la précision prime sur la vitesse]
À l'épreuve de chinois, un caractère mal tracé — trait manquant, radical confondu — est compté \textbf{faux}, même si l'intention est claire. Écrire \zh{徽信} au lieu de \zh{微信} ou \zh{蜜码} au lieu de \zh{密码} coûte des points. Mieux vaut écrire moins de caractères, mais les écrire exactement. Entraînez-vous à recopier chaque paire de cette leçon cinq fois en prononçant le pinyin à voix haute : la mémoire de la main et celle de l'oreille se renforcent mutuellement.
\end{savaistu}

\begin{aretenir}[L'essentiel de la section]
\begin{itemize}
\item Tout caractère proche se distingue par son \cle{radical sémantique} : \zh{讠} parole, \zh{扌} main, \zh{贝} argent, \zh{纟} fil/soie, \zh{足} pied, \zh{虫} insecte, \zh{巾} tissu, \zh{山} montagne.
\item Les six paires à maîtriser : \zh{微}/\zh{徽}, \zh{博}/\zh{搏}, \zh{络}/\zh{路}, \zh{辨}/\zh{辩}, \zh{账}/\zh{帐}, \zh{密}/\zh{蜜}.
\item Stratégie : mémoriser chaque paire avec une \textbf{phrase complète} en contexte, jamais le caractère isolé.
\item La ressemblance visuelle ne garantit ni le sens ni la prononciation : toujours décomposer avant d'écrire.
\end{itemize}
\end{aretenir}

\coursec{Modèle de texte commenté et connecteurs}

Dans la section précédente, nous avons appris à distinguer les caractères proches du thème des réseaux sociaux (\zh{聊} et \zh{柳}, \zh{网} et \zh{冈}, \zh{费} et \zh{佛}…). Vous savez désormais écrire chaque caractère sans le confondre avec son voisin. Il reste une dernière étape pour réussir l'expression écrite au baccalauréat : \textbf{assembler ces caractères en un texte bien construit}. Un bon texte chinois n'est pas une suite de phrases jetées au hasard : il suit une architecture claire — introduction, développement, opinion, conclusion — et ses phrases sont reliées par des \cle{connecteurs}. C'est exactement ce que nous allons apprendre ici, à partir d'un texte modèle intitulé \zh{我和社交媒体} (« Les réseaux sociaux et moi »).

\courssub{Texte modèle : 我和社交媒体}

Lisez d'abord le texte en entier, à voix haute, sans regarder la traduction. Essayez de repérer les mots que vous connaissez déjà et de deviner le sens des phrases nouvelles grâce au contexte.

\begin{textebox}[我和社交媒体 — Wǒ hé shèjiāo méitǐ]
现在，很多年轻人都喜欢用社交媒体。\\
Xiànzài, hěn duō niánqīngrén dōu xǐhuan yòng shèjiāo méitǐ.\\
« Maintenant, beaucoup de jeunes aiment utiliser les réseaux sociaux. »\\[4pt]
我每天用微信跟朋友聊天，也经常刷视频。\\
Wǒ měitiān yòng Wēixìn gēn péngyou liáotiān, yě jīngcháng shuā shìpín.\\
« Chaque jour, j'utilise WeChat pour discuter avec mes amis, et je fais aussi souvent défiler des vidéos. »\\[4pt]
除了聊天以外，我还用微信付钱、看新闻。\\
Chúle liáotiān yǐwài, wǒ hái yòng Wēixìn fù qián, kàn xīnwén.\\
« En plus de discuter, j'utilise aussi WeChat pour payer et lire les informations. »\\[4pt]
社交媒体很方便，但是也很浪费时间。\\
Shèjiāo méitǐ hěn fāngbiàn, dànshì yě hěn làngfèi shíjiān.\\
« Les réseaux sociaux sont très pratiques, mais ils font aussi perdre beaucoup de temps. »\\[4pt]
虽然有时候我很累，但是我还是要先写作业。\\
Suīrán yǒu shíhou wǒ hěn lèi, dànshì wǒ háishì yào xiān xiě zuòyè.\\
« Bien que je sois parfois fatigué, je dois quand même d'abord faire mes devoirs. »\\[4pt]
我觉得我们应该合理使用网络。\\
Wǒ juéde wǒmen yīnggāi hélǐ shǐyòng wǎngluò.\\
« Je pense que nous devons utiliser Internet de manière raisonnable. »
\end{textebox}

\textbf{Glossaire du texte}

\begin{center}
\begin{tabularx}{\linewidth}{|X|X|X|X|}
\hline
\rowcolor{popblueL} Caractères & Pinyin & Nature & Traduction \\
\hline
社交媒体 & shèjiāo méitǐ & nom & réseaux sociaux (médias sociaux) \\
\hline
微信 & Wēixìn & nom propre & WeChat \\
\hline
聊天 & liáotiān & verbe & bavarder, discuter \\
\hline
刷视频 & shuā shìpín & verbe + compl. & faire défiler des vidéos \\
\hline
付钱 & fù qián & verbe + compl. & payer \\
\hline
新闻 & xīnwén & nom & informations, nouvelles \\
\hline
方便 & fāngbiàn & adjectif & pratique, commode \\
\hline
浪费时间 & làngfèi shíjiān & verbe + compl. & perdre du temps \\
\hline
合理 & hélǐ & adjectif/adverbe & raisonnable(ment) \\
\hline
使用 & shǐyòng & verbe & utiliser \\
\hline
\end{tabularx}
\end{center}

\textbf{Questions de compréhension}
\begin{enumerate}
\item Que font les jeunes selon la première phrase ?
\item Quelles sont les trois choses que l'auteur fait avec WeChat ?
\item Quel est l'inconvénient des réseaux sociaux cité dans le texte ?
\item Quelle est la conclusion de l'auteur ?
\end{enumerate}

\courssub{La structure du texte : 开头、正文、看法、结尾}

Relisons le texte en observant \emph{le rôle de chaque phrase}. Vous remarquez qu'il suit toujours le même plan en quatre étapes, typique de la rédaction chinoise au lycée :

\begin{itemize}
\item \cle{开头} (kāitóu, « l'ouverture ») : on pose le sujet avec une phrase générale, souvent introduite par \zh{现在} ou \zh{很多人都…} : \zh{现在，很多年轻人都喜欢用社交媒体。}
\item \cle{正文} (zhèngwén, « le corps ») : on donne les faits, ses habitudes, avec \zh{我每天…}、\zh{我经常…} : \zh{我每天用微信跟朋友聊天，也经常刷视频。}
\item \cle{看法} (kànfǎ, « le point de vue ») : on donne son opinion avec \zh{我觉得…} ou \zh{我认为…} : \zh{我觉得我们应该合理使用网络。}
\item \cle{结尾} (jiéwěi, « la fin ») : on conclut avec un conseil ou une phrase toute faite : \zh{所以，我们应该合理使用网络。}
\end{itemize}

\begin{popfigure}
\begin{tikzpicture}[
  box/.style={draw=popblue, fill=popblueL, rounded corners, thick, minimum width=2.4cm, minimum height=1.1cm, align=center, font=\small},
  arr/.style={-{Stealth[length=3mm]}, thick, poporange},
  lab/.style={font=\footnotesize\itshape, popdark, midway, below}
]
\node[box, align=center] (a) at (0,0){开头\\ \footnotesize introduction};
\node[box, align=center] (b) at (3.9,0){正文\\ \footnotesize faits, habitudes};
\node[box, align=center] (c) at (7.8,0){看法\\ \footnotesize opinion};
\node[box, align=center] (d) at (11.7,0){结尾\\ \footnotesize conclusion};
\draw[arr] (a) -- (b) node[lab] {首先 / 每天};
\draw[arr] (b) -- (c) node[lab] {但是 / 我觉得};
\draw[arr] (c) -- (d) node[lab] {所以 / 应该};
\end{tikzpicture}
\legende{Architecture d'un texte simple en chinois : quatre blocs reliés par des connecteurs.}
\end{popfigure}

\begin{aretenir}[La règle d'or de la rédaction chinoise]
Une phrase = une idée. Un connecteur = un pont entre deux idées. Un texte de 8 à 10 lignes = 4 blocs (开头 → 正文 → 看法 → 结尾). Si vous respectez ce plan, votre copie est déjà à moitié réussie.
\end{aretenir}

\courssub{Les connecteurs utiles}

Observons d'abord les paires de connecteurs du texte modèle, puis énonçons les règles.

\begin{center}
\begin{tabularx}{\linewidth}{|X|X|X|X|}
\hline
\rowcolor{popblueL} Chinois & Pinyin & Français & Exemple \\
\hline
首先…然后… & shǒuxiān… ránhòu… & d'abord… ensuite… & 首先写作业，然后玩手机。 \\
\hline
但是 & dànshì & mais & 很方便，但是很浪费时间。 \\
\hline
因为…所以… & yīnwèi… suǒyǐ… & parce que… (donc)… & 因为方便，所以大家都用。 \\
\hline
虽然…但是… & suīrán… dànshì… & bien que… (mais)… & 虽然很累，但是我要写作业。 \\
\hline
除了…以外，还… & chúle… yǐwài, hái… & en plus de…, aussi… & 除了聊天以外，我还付钱。 \\
\hline
\end{tabularx}
\end{center}

\begin{propriete}[因为…所以… : le couple « parce que… donc… » est correct en chinois]
En français, on ne dit jamais « parce que… donc… » dans la même phrase : on choisit l'un ou l'autre. En chinois, c'est le contraire : les deux mots s'emploient \textbf{ensemble} très naturellement.\\
Structure : 因为 + cause，所以 + résultat。\\
Exemple : \zh{因为社交媒体很方便，所以大家都喜欢用。} \emph{Yīnwèi shèjiāo méitǐ hěn fāngbiàn, suǒyǐ dàjiā dōu xǐhuan yòng.} « Parce que les réseaux sociaux sont pratiques, tout le monde aime les utiliser. »\\
Remarque : dans la conversation rapide, on peut garder seulement \zh{所以} : \zh{我很累，所以我要睡觉。} \emph{Wǒ hěn lèi, suǒyǐ wǒ yào shuìjiào.} « Je suis fatigué, donc je vais dormir. » À l'écrit d'examen, préférez toujours le couple complet.
\end{propriete}

\begin{attention}[虽然…但是… : ne traduisez pas mot à mot vers le français]
En chinois, \zh{虽然} et \zh{但是} s'emploient \textbf{ensemble} dans la même phrase : \zh{虽然很累，但是我要写作业。} \emph{Suīrán hěn lèi, dànshì wǒ yào xiě zuòyè.} En français, on traduit par « Bien que je sois fatigué, je dois faire mes devoirs » — jamais « bien que… mais… ». À l'inverse, au thème (français → chinois), si le français dit « bien que », vous devez mettre \zh{但是} dans la seconde partie : c'est une erreur classique de l'oublier. À la place de \zh{但是}, on peut aussi utiliser \zh{可是} (kěshì) ou \zh{不过} (búguò), plus légers.
\end{attention}

\begin{exemplebox}[除了…以外，还… : « en plus de…, aussi… »]
\zh{除了聊天以外，我还用微信付钱。} \emph{Chúle liáotiān yǐwài, wǒ hái yòng Wēixìn fù qián.} « En plus de discuter, j'utilise aussi WeChat pour payer. »\\
\zh{除了英语以外，我还学中文。} \emph{Chúle Yīngyǔ yǐwài, wǒ hái xué Zhōngwén.} « En plus de l'anglais, j'apprends aussi le chinois. »\\
Attention à l'ordre : le bloc \zh{除了…以外} se place \textbf{avant le sujet} ou juste après lui, et \zh{还} se place toujours \textbf{avant le verbe principal}. Ne dites jamais \zh{我用微信还付钱} : \zh{还} doit précéder \zh{用}.
\end{exemplebox}

\courssub{Les expressions d'opinion et les phrases de conclusion}

Pour exprimer votre point de vue (看法), trois outils suffisent au niveau Terminale :

\begin{itemize}
\item \zh{我觉得…} \emph{wǒ juéde…} : « je pense que… » (le plus courant, oral et écrit) ;
\item \zh{我认为…} \emph{wǒ rènwéi…} : « je considère que… » (plus soutenu, idéal à l'écrit) ;
\item \zh{对我来说，…} \emph{duì wǒ lái shuō, …} : « pour moi, … » (pour introduire un point de vue personnel).
\end{itemize}

Exemple : \zh{对我来说，社交媒体是朋友，也是老师。} \emph{Duì wǒ lái shuō, shèjiāo méitǐ shì péngyou, yě shì lǎoshī.} « Pour moi, les réseaux sociaux sont un ami et aussi un professeur. »

Pour la conclusion (结尾), apprenez par cœur ces phrases toutes faites, réutilisables dans presque toutes les rédactions sur le thème :

\begin{itemize}
\item \zh{我们应该合理使用网络。} \emph{Wǒmen yīnggāi hélǐ shǐyòng wǎngluò.} « Nous devons utiliser Internet de façon raisonnable. »
\item \zh{所以，我们不能花太多时间玩手机。} \emph{Suǒyǐ, wǒmen bù néng huā tài duō shíjiān wán shǒujī.} « Donc, nous ne devons pas passer trop de temps sur nos téléphones. »
\item \zh{网络很有用，但是学习更重要。} \emph{Wǎngluò hěn yǒuyòng, dànshì xuéxí gèng zhòngyào.} « Internet est utile, mais les études sont plus importantes. »
\end{itemize}

Notez l'emploi de \zh{花} (huā, « dépenser, passer ») dans \zh{花时间} : c'est le verbe correct pour « passer du temps », et non \zh{用}, qui signifie « utiliser ».

\courssub{Méthode de rédaction et entraînement}

\begin{methode}[Rédiger un texte « Les réseaux sociaux et moi » en 4 étapes]
\begin{enumerate}
\item \textbf{Faire un plan en 3 idées} : mes habitudes (正文), le pour et le contre (avantages / inconvénients), mon avis (看法).
\item \textbf{Une phrase = une idée} : phrases courtes, structure Sujet + Verbe + Complément ; ne cherchez pas à impressionner, cherchez à être correct.
\item \textbf{Relier avec un connecteur} : au moins trois connecteurs différents dans le texte (\zh{但是}、\zh{因为…所以…}、\zh{除了…以外，还…}).
\item \textbf{Relire en vérifiant chaque caractère} : surtout les caractères proches vus dans la section précédente (\zh{聊}、\zh{网}、\zh{费}), les tons du pinyin au brouillon, et la présence de \zh{但是} après \zh{虽然}.
\end{enumerate}
\end{methode}

\begin{exoresolu}[Thème : traduire en chinois]
Énoncé : Traduisez en chinois : « Bien que les réseaux sociaux soient pratiques, ils font perdre du temps, donc nous devons les utiliser de façon raisonnable. »
\tcblower
Solution détaillée :\\
1. « Bien que… » impose le couple \zh{虽然…但是…} : ne pas oublier \zh{但是}.\\
2. « pratique » = \zh{方便} ; « perdre du temps » = \zh{浪费时间} ; « de façon raisonnable » = \zh{合理} (placé avant le verbe \zh{使用}).\\
3. « donc » = \zh{所以}, placé devant la dernière proposition.\\[3pt]
Traduction : \zh{虽然社交媒体很方便，但是也很浪费时间，所以我们应该合理使用。}\\
\emph{Suīrán shèjiāo méitǐ hěn fāngbiàn, dànshì yě hěn làngfèi shíjiān, suǒyǐ wǒmen yīnggāi hélǐ shǐyòng.}
\end{exoresolu}

\begin{exercice}[À vous]
1. Version : traduisez en français : \zh{除了刷视频以外，我还用微信跟家人聊天。}\\
2. Thème : traduisez en chinois : « Parce que WeChat est très pratique, beaucoup de jeunes l'utilisent chaque jour. »\\
3. Rédaction : écrivez 6 à 8 phrases sur « 我和社交媒体 » en suivant le plan 开头 → 正文 → 看法 → 结尾, avec au moins trois connecteurs différents.
\end{exercice}

\begin{corrige}[Exercice « À vous »]
1. « En plus de faire défiler des vidéos, j'utilise aussi WeChat pour discuter avec ma famille. »\\
2. \zh{因为微信很方便，所以很多年轻人每天都用。} \emph{Yīnwèi Wēixìn hěn fāngbiàn, suǒyǐ hěn duō niánqīngrén měitiān dōu yòng.}\\
3. Modèle attendu : une phrase d'ouverture avec \zh{现在}, deux ou trois phrases d'habitudes avec \zh{每天} / \zh{经常}, un contraste avec \zh{但是}, une opinion avec \zh{我觉得}, une conclusion avec \zh{我们应该合理使用网络。} Vérifier : présence des couples \zh{因为…所以…} et \zh{虽然…但是…} complets, caractères corrects.
\end{corrige}

\begin{attention}[Les trois pièges de la copie d'examen]
\begin{itemize}
\item Écrire \zh{虽然…} sans \zh{但是…} : la phrase reste « ouverte » et coûte des points.
\item Empiler les idées dans une seule longue phrase : préférez deux phrases courtes reliées par un connecteur.
\item Traduire le français mot à mot : « pour moi » se dit \zh{对我来说} et non une construction inventée ; utilisez uniquement les formules apprises.
\end{itemize}
\end{attention}

Vous disposez maintenant de tout le nécessaire pour construire un texte : un plan en quatre blocs, cinq connecteurs maîtrisés, trois expressions d'opinion et des phrases de conclusion toutes faites. La section suivante vous entraînera à traduire dans les deux sens, du français vers le chinois et du chinois vers le français.

\coursec{Thème et version : traduire dans les deux sens}

Dans la section précédente, nous avons étudié un modèle de texte commenté sur les réseaux sociaux et appris à enchaîner nos idées grâce aux connecteurs \zh{首先} (d'abord), \zh{然后} (ensuite), \zh{但是} (mais) et \zh{所以} (donc). Il est maintenant temps de franchir une étape décisive : passer d'une langue à l'autre. Au baccalauréat, l'épreuve de chinois comporte toujours un \cle{thème} (traduire du français vers le chinois) et une \cle{version} (traduire du chinois vers le français). Ces deux exercices ne se réussissent pas « au feeling » : ils exigent une méthode rigoureuse. C'est l'objet de cette section.

\courssub{Le thème : du français vers le chinois}

\textbf{Observer d'abord}\par

Lisons attentivement les deux phrases suivantes :

\zh{我经常在社交媒体上发照片。} \emph{(Wǒ jīngcháng zài shèjiāo méitǐ shàng fā zhàopiàn.)} « Je publie souvent des photos sur les réseaux sociaux. »

\zh{我每天用微信。} \emph{(Wǒ měitiān yòng Wēixìn.)} « J'utilise WeChat tous les jours. »

Comparons avec le français : en français, « souvent », « tous les jours », « sur les réseaux sociaux » se placent \emph{après} le verbe. En chinois, c'est l'inverse : \zh{经常} (souvent), \zh{每天} (tous les jours) et \zh{在社交媒体上} (sur les réseaux sociaux) se placent \emph{avant} le verbe \zh{发} (publier) et \zh{用} (utiliser). C'est la première grande difficulté du thème pour un francophone.

\begin{propriete}[La place des compléments avant le verbe]
En chinois, les compléments circonstanciels de \cle{temps}, de \cle{lieu} (introduits par \zh{在}) et de \cle{manière ou d'instrument} (introduits par \zh{用}, \zh{跟}) se placent \textbf{avant le verbe}, et non après comme en français.\\
Structure : \textbf{Sujet + [temps] + [lieu avec 在] + [avec 跟 / instrument 用] + Verbe + Complément d'objet}\\
\zh{我每天用微信。} \emph{(Wǒ měitiān yòng Wēixìn.)} « J'utilise WeChat tous les jours. » — et non *\zh{我用微信每天} (faute typique du francophone !)
\end{propriete}

\begin{popfigure}
\begin{tikzpicture}[node distance=0.45cm and 0.5cm]
\node[draw=popblue, fill=popblueL, rounded corners, align=center, font=\small] (s) {Sujet\\ \zh{我}};
\node[draw=popgreen, fill=popgreenL, rounded corners, align=center, font=\small, right=of s] (t) {Temps\\ \zh{每天}};
\node[draw=poporange, fill=poporangeL, rounded corners, align=center, font=\small, right=of t] (l) {Lieu\\ \zh{在网上}};
\node[draw=poppurple, fill=poppurpleL, rounded corners, align=center, font=\small, right=of l] (m) {Avec / instrument\\ \zh{跟朋友} / \zh{用电脑}};
\node[draw=popdark, fill=popcream, rounded corners, align=center, font=\small, right=of m] (v) {Verbe\\ \zh{聊天}};
\node[draw=poppink, fill=poppinkL, rounded corners, align=center, font=\small, right=of v] (c) {Objet\\ (\zh{照片}…)};
\draw[-{Stealth}, popdark] (s) -- (t);
\draw[-{Stealth}, popdark] (t) -- (l);
\draw[-{Stealth}, popdark] (l) -- (m);
\draw[-{Stealth}, popdark] (m) -- (v);
\draw[-{Stealth}, popdark] (v) -- (c);
\end{tikzpicture}
\legende{L'ordre des compléments en chinois : tout ce qui encadre l'action (quand ? où ? avec qui ? avec quoi ?) précède le verbe.}
\end{popfigure}

\begin{exemplebox}[L'instrument avec 用]
\zh{她用电脑上网。} \emph{(Tā yòng diànnǎo shàngwǎng.)} « Elle se connecte avec son ordinateur. »\\
L'instrument introduit par \zh{用} (avec, au moyen de) précède le verbe \zh{上网} (se connecter). En français, « avec son ordinateur » vient après ; en chinois, \zh{用电脑} vient avant.
\end{exemplebox}

\begin{attention}[跟...聊天 : discuter AVEC quelqu'un]
\zh{跟} signifie « avec » et se place \textbf{avant le verbe} : \zh{我跟朋友聊天。} \emph{(Wǒ gēn péngyou liáotiān.)} « Je discute avec mes amis. » Ne dites jamais *\zh{我聊天跟朋友}, calque du français « je discute avec mes amis ».
\end{attention}

\begin{methode}[La méthode du thème en quatre étapes]
\begin{enumerate}
\item \textbf{Analyser la phrase française} : repérer le sujet, le verbe, le complément d'objet et les compléments circonstanciels (temps, lieu, manière).
\item \textbf{Choisir les caractères du thème} : mobiliser le lexique des réseaux sociaux étudié dans cette leçon (\zh{发}, \zh{删}, \zh{评论}, \zh{隐私}, \zh{网友}…) et les classificateurs appropriés (\zh{张} pour les photos, \zh{条} pour les messages).
\item \textbf{Réorganiser à la chinoise} : placer sujet, puis temps, lieu, instrument, puis verbe et objet ; choisir l'aspect (\zh{了}, \zh{过}, \zh{着}) si nécessaire.
\item \textbf{Vérifier les caractères} : relire chaque caractère (radical, traits, caractère proche) et vérifier la ponctuation chinoise 。 ，
\end{enumerate}
\end{methode}

\begin{popfigure}
\begin{tikzpicture}[node distance=0.5cm]
\node[draw=popblue, fill=popblueL, rounded corners, align=center, font=\small, text width=4.6cm] (a) {1. Phrase française\\ « Je publie souvent des photos. »};
\node[draw=popgreen, fill=popgreenL, rounded corners, align=center, font=\small, text width=4.6cm, below=of a] (b) {2. Analyse S / V / C\\ S = je ; V = publier ; C = des photos ; temps = souvent};
\node[draw=poporange, fill=poporangeL, rounded corners, align=center, font=\small, text width=4.6cm, below=of b] (c) {3. Réorganisation chinoise\\ \zh{我 + 经常 + 在社交媒体上 + 发 + 照片}};
\node[draw=poppurple, fill=poppurpleL, rounded corners, align=center, font=\small, text width=4.6cm, below=of c] (d) {4. Vérification\\ caractères, tons, ponctuation 。};
\draw[-{Stealth}, thick, popdark] (a) -- (b);
\draw[-{Stealth}, thick, popdark] (b) -- (c);
\draw[-{Stealth}, thick, popdark] (c) -- (d);
\end{tikzpicture}
\legende{Organigramme du thème : de la phrase française à la phrase chinoise correcte.}
\end{popfigure}

\courssub{Trois thèmes guidés pas à pas}

\textbf{Thème 1 : « Je publie souvent des photos sur les réseaux sociaux. »}\par

Analyse : sujet = je (\zh{我}) ; temps = souvent (\zh{经常}) ; lieu = sur les réseaux sociaux (\zh{在社交媒体上}) ; verbe = publier (\zh{发}) ; objet = des photos (\zh{照片}). En réorganisant : Sujet + temps + lieu + verbe + objet.

\zh{我经常在社交媒体上发照片。} \emph{(Wǒ jīngcháng zài shèjiāo méitǐ shàng fā zhàopiàn.)}

\textbf{Thème 2 : « Il a supprimé ce commentaire parce qu'il était dangereux. »}\par

L'action est accomplie : on ajoute \zh{了} après le verbe \zh{删} (supprimer) — ou le composé \zh{删掉} — pour marquer l'aspect accompli. « Ce commentaire » exige le classificateur \zh{条} (pour les messages, commentaires, informations) : \zh{那条评论}. La cause s'exprime avec \zh{因为} (parce que).

\zh{他删掉了那条评论，因为它很危险。} \emph{(Tā shāndiào le nà tiáo pínglùn, yīnwèi tā hěn wēixiǎn.)}

\textbf{Thème 3 : « Bien que WeChat soit pratique, il faut protéger sa vie privée. »}\par

La concession se traduit par le couple \zh{虽然……但是……} (bien que... mais...). Contrairement au français, le chinois emploie souvent les deux connecteurs ensemble ; à l'écrit, on garde les deux.

\zh{虽然微信很方便，但是我们要保护隐私。} \emph{(Suīrán Wēixìn hěn fāngbiàn, dànshì wǒmen yào bǎohù yǐnsī.)}

\begin{attention}[Le piège des collocations : « perdre du temps »]
Ne traduisez jamais mot à mot ! « Perdre du temps » ne se dit pas *\zh{失时间}. La collocation correcte est \zh{浪费时间} \emph{(làngfèi shíjiān)}, littéralement « gaspiller le temps ». De même, « passer du temps sur les réseaux » se dit \zh{花时间在社交媒体上}. Apprenez les verbes avec leurs compléments habituels : \zh{发照片} (publier des photos), \zh{删评论} (supprimer un commentaire), \zh{保护隐私} (protéger sa vie privée), \zh{浪费时间} (perdre du temps).
\end{attention}

\courssub{La version : du chinois vers le français}

\textbf{Observer d'abord}\par

Lisons ce court texte :

\begin{textebox}[Un message d'élève]
\zh{我的网友每天给我发很多消息，有时候我觉得很累。}\\
\emph{Wǒ de wǎngyǒu měitiān gěi wǒ fā hěn duō xiāoxi, yǒu shíhou wǒ juéde hěn lèi.}\\
« Mon ami en ligne m'envoie beaucoup de messages chaque jour ; parfois je me sens fatigué. »
\end{textebox}

Découpons : \zh{我的网友} (mon ami en ligne) / \zh{每天} (chaque jour) / \zh{给我} (à moi) / \zh{发很多消息} (envoie beaucoup de messages) / \zh{有时候} (parfois) / \zh{我觉得很累} (je me sens très fatigué). Remarquez que \zh{每天} et \zh{给我} précèdent le verbe \zh{发} : en français, « chaque jour » et « m' » se réorganisent autour du verbe « envoie ».

\begin{methode}[La méthode de la version en trois étapes]
\begin{enumerate}
\item \textbf{Souligner d'abord les mots du glossaire} et les mots connus (noms, verbes, connecteurs) ; repérer les mots-outils \zh{的}, \zh{了}, \zh{在}, \zh{跟}, \zh{用}.
\item \textbf{Traduire segment par segment} : découper la phrase en blocs de sens (sujet / circonstances / verbe / objet), traduire chaque bloc séparément.
\item \textbf{Reformuler en français correct et naturel} : ne jamais coller au mot à mot ; replacer les compléments selon la grammaire française et choisir un français idiomatique.
\end{enumerate}
\end{methode}

\textbf{Version guidée, étape par étape}\par

\begin{itemize}
\item Étape 1 — mots connus : \zh{网友} (ami en ligne), \zh{每天} (chaque jour), \zh{发} (envoyer), \zh{消息} (message), \zh{有时候} (parfois), \zh{觉得} (ressentir, trouver), \zh{累} (fatigué).
\item Étape 2 — segments : « mon ami en ligne » / « chaque jour » / « à moi envoie » / « beaucoup de messages » ; « parfois » / « je me sens très fatigué ».
\item Étape 3 — reformulation : « Mon ami en ligne m'envoie beaucoup de messages chaque jour ; parfois je me sens fatigué. »
\end{itemize}

\courssub{Les pièges de traduction : ordre des compléments et aspect}

\begin{propriete}[L'ordre des compléments : 时间 + 地点 + 方式 avant le verbe]
Quand une phrase chinoise cumule plusieurs circonstances, elles se rangent dans l'ordre : \textbf{temps (\zh{时间}) → lieu (\zh{地点}) → manière/instrument (\zh{方式})}, toutes avant le verbe.\\
\zh{我昨天在网上跟朋友聊天。} \emph{(Wǒ zuótiān zài wǎngshàng gēn péngyou liáotiān.)} « Hier, j'ai discuté en ligne avec mes amis. »\\
En version, il faut redistribuer ces éléments selon l'ordre français : « Hier, j'ai discuté avec mes amis en ligne. »
\end{propriete}

\begin{propriete}[了, 过, 着 : trois marques d'aspect à ne pas confondre]
\begin{itemize}
\item \zh{了} : action \textbf{accomplie} → passé composé français. \zh{他删了那条评论。} « Il a supprimé ce commentaire. »
\item \zh{过} : expérience \textbf{vécue au moins une fois} → « avoir déjà ». \zh{我用过微信。} « J'ai déjà utilisé WeChat. »
\item \zh{着} : action ou état \textbf{en cours, duratif} → « en train de ». \zh{他看着手机。} « Il est en train de regarder son téléphone. »
\end{itemize}
En thème, demandez-vous toujours : l'action est-elle accomplie (\zh{了}), vécue (\zh{过}) ou en cours (\zh{着}) ?
\end{propriete}

\begin{center}
\begin{tabularx}{\linewidth}{|X|X|X|X|}
\hline
\rowcolor{popblueL} Caractères & Pinyin & Nature & Traduction \\
\hline
发照片 & fā zhàopiàn & verbe + nom & publier des photos \\
\hline
删评论 & shān pínglùn & verbe + nom & supprimer un commentaire \\
\hline
保护隐私 & bǎohù yǐnsī & verbe + nom & protéger sa vie privée \\
\hline
浪费时间 & làngfèi shíjiān & verbe + nom & perdre (gaspiller) du temps \\
\hline
网友 & wǎngyǒu & nom & ami en ligne \\
\hline
消息 & xiāoxi & nom & message, nouvelle \\
\hline
跟……聊天 & gēn... liáotiān & structure & discuter avec... \\
\hline
用电脑上网 & yòng diànnǎo shàngwǎng & structure & se connecter avec un ordinateur \\
\hline
虽然……但是…… & suīrán... dànshì... & connecteurs & bien que... mais... \\
\hline
因为 & yīnwèi & conjonction & parce que \\
\hline
\end{tabularx}
\end{center}

\begin{exoresolu}[Thème et version au baccalauréat]
Énoncé — 1) Traduisez en chinois : « Hier, elle a publié deux photos sur les réseaux sociaux. » 2) Traduisez en français : \zh{虽然网络很方便，但是我们不应该浪费时间。}
\tcblower
Solution détaillée — 1) Analyse : sujet = elle (\zh{她}) ; temps = hier (\zh{昨天}) ; lieu = sur les réseaux sociaux (\zh{在社交媒体上}) ; verbe = publier (\zh{发}) avec \zh{了} (action accomplie) ; objet = deux photos (\zh{两张照片}, classificateur \zh{张} pour les objets plats). Réorganisation : \zh{她昨天在社交媒体上发了两张照片。} \emph{(Tā zuótiān zài shèjiāo méitǐ shàng fā le liǎng zhāng zhàopiàn.)}\\
2) Segments : \zh{虽然} (bien que) / \zh{网络很方便} (Internet soit pratique) / \zh{但是} (mais) / \zh{我们不应该} (nous ne devrions pas) / \zh{浪费时间} (perdre du temps). Reformulation : « Bien qu'Internet soit pratique, nous ne devrions pas perdre notre temps. »
\end{exoresolu}

\begin{exercice}[À vous de jouer]
1) Thème : traduisez en chinois. a) « Je discute souvent avec mes amis sur WeChat. » b) « Il a supprimé ce message parce qu'il était faux. »\\
2) Version : traduisez en français. \zh{我的网友每天用电脑给我发消息，有时候我觉得很累。}
\end{exercice}

\begin{corrige}[Exercice « À vous de jouer »]
1) a) Sujet + temps + lieu + avec + verbe : \zh{我经常在微信上跟朋友聊天。} \emph{(Wǒ jīngcháng zài Wēixìn shàng gēn péngyou liáotiān.)} b) Action accomplie + cause : \zh{他删掉了那条消息，因为它是假的。} \emph{(Tā shāndiào le nà tiáo xiāoxi, yīnwèi tā shì jiǎ de.)}\\
2) Segments puis reformulation : « Mon ami en ligne m'envoie des messages chaque jour avec son ordinateur ; parfois je me sens fatigué. »
\end{corrige}

\begin{methode}[Grille d'évaluation et entraînement au bac]
\textbf{Critères de notation (sur 10 points pour l'ensemble thème + version) :}
\begin{enumerate}
\item \textbf{Respect de l'ordre des mots chinois (3 pts)} : Sujet -- Temps -- Lieu -- Manière/Instrument -- Verbe -- Objet. Toute inversion (ex: verbe avant lieu) coûte 1 pt.
\item \textbf{Justesse lexicale et collocations (3 pts)} : Vocabulaire du thème utilisé correctement (\zh{发照片}, \zh{删评论}, \zh{保护隐私}, \zh{浪费时间}). Pas de calque (*\zh{失时间}, *\zh{聊天跟}).
\item \textbf{Maîtrise des mots-outils et aspect (2 pts)} : \zh{了} (accompli), \zh{过} (expérience), \zh{着} (duratif) ; \zh{在} (lieu), \zh{跟/用} (compagnon/instrument), \zh{的/得/地}.
\item \textbf{Qualité du français (version) / Caractères (thème) (2 pts)} : Français naturel, idiomatique (pas de mot à mot) ; caractères lisibles, traits dans l'ordre, radicaux respectés, ponctuation pleine chasse 。
\end{enumerate}
\textbf{Conseil d'examen} : Faites d'abord le thème (plus contraignant), relisez vos caractères. Pour la version, écrivez un brouillon mot à mot, puis une version soignée sur la copie.
\end{methode}

\begin{aretenir}[L'essentiel de la section]
\begin{itemize}
\item Thème : analyser (S/V/C), réorganiser (temps + lieu + manière \textbf{avant} le verbe), choisir l'aspect (\zh{了/过/着}), vérifier les caractères et classificateurs (\zh{张}, \zh{条}).
\item Version : souligner les mots connus, traduire segment par segment, reformuler en français naturel.
\item Collocations à mémoriser : \zh{发照片}, \zh{删评论}, \zh{保护隐私}, \zh{浪费时间} — jamais de mot à mot !
\end{itemize}
\end{aretenir}

\coursec{Grille d'évaluation et entraînement à l'examen}

Nous savons désormais traduire dans les deux sens des phrases simples sur les réseaux sociaux. Il reste une dernière étape, décisive : se préparer concrètement à l'épreuve d'expression écrite. Au baccalauréat, la copie n'est pas notée « au feeling » : le correcteur applique une \cle{grille d'évaluation} précise. Connaître cette grille, s'entraîner avec elle, c'est déjà gagner des points. Cette dernière section vous donne les clés de l'examen : le barème, la méthode de relecture, la gestion du temps et la correction d'une copie d'élève.

\courssub{Observer d'abord : une copie d'élève corrigée}

Avant d'énoncer les critères, observons un texte réel d'élève (niveau Tle A) sur le sujet « Mes réseaux sociaux préférés ». Ce texte contient \textbf{5 erreurs typiques}. Cherchez-les avant de lire la correction.

\begin{textebox}[Copie d'élève — avec 5 erreurs]
我用微信每天。微信是很好，我和朋友聊天。\\
我喜观抖音，因为视频很有意思。\\
但是，上网很多时间是不好的。\\
所以，我们应该少用手机，多学习。\\
(Wǒ yòng Wēixìn měitiān. Wēixìn shì hěn hǎo, wǒ hé péngyou liáotiān. Wǒ xǐguān Dǒuyīn, yīnwèi shìpín hěn yǒu yìsi. Dànshì, shàngwǎng hěn duō shíjiān shì bù hǎo de. Suǒyǐ, wǒmen yīnggāi shǎo yòng shǒujī, duō xuéxí.)\\
« J'utilise WeChat chaque jour. WeChat est bien, je discute avec mes amis. J'aime Douyin, parce que les vidéos sont intéressantes. Mais passer beaucoup de temps sur Internet n'est pas bon. Donc nous devons moins utiliser le téléphone et plus étudier. »
\end{textebox}

\textbf{Les 5 erreurs cachées :}
\begin{enumerate}
\item \textbf{Ordre des mots} : ✗ \zh{我用微信每天。} → ✓ \zh{我每天用微信。} (Wǒ měitiān yòng Wēixìn.) — le complément de temps \zh{每天} se place \emph{avant} le verbe, jamais à la fin comme en français.
\item \textbf{Structure 是 + adjectif} : ✗ \zh{微信是很好} → ✓ \zh{微信很好} (Wēixìn hěn hǎo.) — on n'utilise pas \zh{是} devant un adjectif qualificatif : l'adjectif seul (précédé de \zh{很}) suffit comme prédicat.
\item \textbf{Caractère proche} : ✗ \zh{喜观} → ✓ \zh{喜欢} (xǐhuan, aimer) — confusion entre \zh{观} (guān, regarder) et \zh{欢} (huan, joyeux).
\item \textbf{Ordre des mots (complément)} : ✗ \zh{上网很多时间} → ✓ \zh{花很多时间上网} (huā hěn duō shíjiān shàngwǎng) ou \zh{上网时间很长} (shàngwǎng shíjiān hěn cháng) — le français « passer beaucoup de temps sur Internet » ne se calque pas mot à mot : il faut le verbe \zh{花} (consacrer) + durée + activité.
\item \textbf{Connecteur manquant / phrase télégraphique} : la deuxième phrase juxtapose deux idées sans lien ; mieux : \zh{微信很好，因为我和朋友聊天。} (Wēixìn hěn hǎo, yīnwèi wǒ hé péngyou liáotiān.) avec le connecteur \zh{因为}.
\end{enumerate}

\courssub{La grille d'évaluation : critères et barème}

\begin{definition}[Grille d'évaluation]
La \cle{grille d'évaluation} est le tableau que le correcteur utilise pour noter votre production écrite. Elle répartit les 20 points entre plusieurs critères : l'exactitude des caractères, la grammaire, les connecteurs, le contenu et la présentation. Connaître la grille, c'est savoir où sont les points.
\end{definition}

\begin{center}
\begin{tabularx}{\linewidth}{|X|X|X|X|}\hline
\rowcolor{popblueL} Critère & Points & Ce que l'examinateur attend & Erreur qui fait perdre des points \\ \hline
Exactitude des caractères & 4 & Traits complets, radicaux corrects, aucun caractère proche confondu & \zh{观} pour \zh{欢}, trait manquant dans \zh{微} \\ \hline
Grammaire et ordre des mots & 6 & Sujet + temps + verbe ; pas de \zh{是} devant un adjectif ; 了/过 bien employés & \zh{我用微信每天} \\ \hline
Connecteurs et cohérence & 4 & Au moins 3 connecteurs variés : \zh{因为、但是、所以、首先、然后} & phrases juxtaposées sans lien \\ \hline
Contenu et respect du sujet & 4 & Sujet traité, 6 à 10 phrases, idées personnelles & hors-sujet, texte trop court \\ \hline
Présentation & 2 & Écriture lisible, ponctuation chinoise ，。 propre & ratures, ponctuation française \\ \hline
\end{tabularx}
\end{center}

\begin{popfigure}
\begin{tikzpicture}[
  crit/.style={rectangle, draw=popblue, fill=popblueL, rounded corners, align=center, text width=3.1cm, minimum height=1.1cm, font=\small},
  pts/.style={rectangle, draw=poporange, fill=poporangeL, rounded corners, align=center, text width=1.3cm, minimum height=1.1cm, font=\small\bfseries},
  att/.style={rectangle, draw=popgreen, fill=popgreenL, rounded corners, align=center, text width=5.6cm, minimum height=1.1cm, font=\small},
  head/.style={font=\small\bfseries, text=popdark}
]
\node[head] at (1.55,5.6) {Critère};
\node[head] at (5.0,5.6) {Points};
\node[head] at (9.2,5.6) {Ce que l'examinateur attend};
\node[crit] at (1.55,4.5) {Caractères (traits, radicaux)};
\node[pts] at (5.0,4.5) {4 pts};
\node[att] at (9.2,4.5) {chaque caractère complet, aucune confusion de caractères proches};
\node[crit] at (1.55,3.3) {Grammaire et ordre des mots};
\node[pts] at (5.0,3.3) {6 pts};
\node[att] at (9.2,3.3) {temps avant le verbe, adjectif sans 是, 了/过 corrects};
\node[crit] at (1.55,2.1) {Connecteurs et cohérence};
\node[pts] at (5.0,2.1) {4 pts};
\node[att] at (9.2,2.1) {au moins 3 connecteurs variés, paragraphes liés};
\node[crit] at (1.55,0.9) {Contenu et longueur};
\node[pts] at (5.0,0.9) {4 pts};
\node[att] at (9.2,0.9) {sujet respecté, 6 à 10 phrases simples};
\node[crit] at (1.55,-0.3) {Présentation};
\node[pts] at (5.0,-0.3) {2 pts};
\node[att] at (9.2,-0.3) {écriture lisible, ponctuation chinoise};
\end{tikzpicture}
\legende{Grille d'évaluation type de l'épreuve d'expression écrite (sur 20 points)}
\end{popfigure}

\begin{propriete}[La correction prime sur l'ambition]
Au baccalauréat, mieux vaut des phrases \textbf{simples et justes} que des phrases complexes et fautives. Une phrase du type \zh{我每天用微信和朋友聊天。} (Wǒ měitiān yòng Wēixìn hé péngyou liáotiān. — « Chaque jour, j'utilise WeChat pour discuter avec mes amis. ») rapporte tous ses points ; une phrase longue avec trois fautes en perd. La règle d'or : \emph{je n'écris que ce que je sais écrire correctement}.
\end{propriete}

\courssub{Gérer son temps : 5 + 20 + 5}

L'épreuve laisse environ 30 minutes pour la production écrite. La répartition optimale est la suivante : \textbf{5 minutes de plan}, \textbf{20 minutes de rédaction}, \textbf{5 minutes de relecture}. Beaucoup d'élèves perdent des points non par manque de connaissances, mais parce qu'ils rédigent sans plan et rendent sans relire.

\begin{popfigure}
\begin{tikzpicture}[
  phase/.style={rectangle, rounded corners, align=center, text width=3.4cm, minimum height=1.6cm, font=\small},
  fleche/.style={-{Stealth[length=3mm]}, very thick, popdark},
  duree/.style={font=\small\bfseries, text=popdark}
]
\node[phase, draw=popblue, fill=popblueL, align=center] (p1) at (0,0){\textbf{PLAN}\\ noter 3 idées + 3 connecteurs + les caractères difficiles};
\node[phase, draw=poporange, fill=poporangeL, align=center] (p2) at (5.6,0){\textbf{RÉDACTION}\\ phrases simples, sujet + temps + verbe, ponctuation chinoise};
\node[phase, draw=popgreen, fill=popgreenL, align=center] (p3) at (11.2,0){\textbf{RELECTURE}\\ 3 passages : sens, grammaire, caractères};
\draw[fleche] (p1) -- (p2);
\draw[fleche] (p2) -- (p3);
\node[duree] at (0,1.6) {5 min};
\node[duree] at (5.6,1.6) {20 min};
\node[duree] at (11.2,1.6) {5 min};
\end{tikzpicture}
\legende{Frise de gestion du temps : 5 min de plan, 20 min de rédaction, 5 min de relecture}
\end{popfigure}

\begin{methode}[La relecture en trois passages]
Relire une seule fois ne suffit pas : l'œil ne peut pas tout voir à la fois. Relisez \textbf{trois fois}, avec un objectif différent à chaque passage.
\begin{enumerate}
\item \textbf{Premier passage — le sens général} : mon texte répond-il au sujet ? Les idées se suivent-elles ? Y a-t-il un connecteur dans chaque paragraphe (\zh{首先、因为、但是、所以}) ?
\item \textbf{Deuxième passage — la grammaire} : vérifier chaque phrase : le complément de temps est-il \emph{avant} le verbe ? Ai-je mis \zh{是} devant un adjectif (à supprimer) ? Les particules \zh{了} et \zh{过} sont-elles justifiées ?
\item \textbf{Troisième passage — les caractères, trait par trait} : repasser sur chaque caractère du thème (\zh{微、信、聊、络、险、播}) : aucun trait manquant ? aucun caractère proche (\zh{观/欢}, \zh{未/末}, \zh{己/已}) ?
\end{enumerate}
\end{methode}

\begin{attention}[Le texte trop court coûte cher]
Un texte de moins de 5 phrases perd des points de contenu, même s'il est parfait. Visez \textbf{6 à 10 phrases simples mais correctes}. Si vous manquez d'idées, utilisez le plan en trois temps : 1) ce que j'utilise, 2) pourquoi, 3) mon avis (avantages et dangers). Et attention : une phrase juste de plus vaut mieux qu'une phrase longue fautive.
\end{attention}

\courssub{Auto-évaluation : la liste de contrôle avant de rendre}

Avant de rendre votre copie, cochez mentalement chaque point de cette liste. Chaque « non » est un point potentiellement perdu.

\begin{center}
\begin{tabularx}{\linewidth}{|X|X|}\hline
\rowcolor{popblueL} Question de contrôle & Critère visé \\ \hline
Chaque caractère est-il complet (traits, radicaux) ? & Caractères (4 pts) \\ \hline
Ai-je évité les confusions 观/欢, 未/末, 己/已 ? & Caractères (4 pts) \\ \hline
Le complément de temps est-il avant le verbe ? & Grammaire (6 pts) \\ \hline
Ai-je supprimé tout 是 devant un adjectif ? & Grammaire (6 pts) \\ \hline
Y a-t-il un connecteur dans chaque paragraphe ? & Connecteurs (4 pts) \\ \hline
Ai-je écrit entre 6 et 10 phrases sur le sujet ? & Contenu (4 pts) \\ \hline
Ma ponctuation est-elle chinoise （，。） et mon écriture lisible ? & Présentation (2 pts) \\ \hline
\end{tabularx}
\end{center}

\courssub{Exercice résolu : corriger la copie d'élève}

\begin{exoresolu}[Retrouver et corriger les erreurs]
Énoncé — Voici trois phrases extraites de copies d'élèves. Corrigez-les et justifiez chaque correction.\\
1. 我用微信每天。\\
2. 抖音是很有意思。\\
3. 我喜观上网，但是上网很多时间不好。
\tcblower
Solution détaillée —\\
1. ✗ \zh{我用微信每天。} → ✓ \zh{我每天用微信。} (Wǒ měitiān yòng Wēixìn. — « J'utilise WeChat chaque jour. ») Justification : en chinois, le complément de temps \zh{每天} se place entre le sujet et le verbe ; l'ordre français « j'utilise WeChat chaque jour » ne peut pas être calqué.\\
2. ✗ \zh{抖音是很有意思。} → ✓ \zh{抖音很有意思。} (Dǒuyīn hěn yǒu yìsi. — « Douyin est très intéressant. ») Justification : l'adjectif qualificatif fonctionne seul comme prédicat, sans \zh{是} ; \zh{很} sert de lien naturel.\\
3. ✗ \zh{我喜观上网，但是上网很多时间不好。} → ✓ \zh{我喜欢上网，但是花很多时间上网不好。} (Wǒ xǐhuan shàngwǎng, dànshì huā hěn duō shíjiān shàngwǎng bù hǎo. — « J'aime aller sur Internet, mais passer beaucoup de temps en ligne n'est pas bon. ») Justification : confusion de caractères proches (\zh{观} → \zh{欢}) et calque du français « passer beaucoup de temps » : il faut le verbe \zh{花} (consacrer) + durée + activité.
\end{exoresolu}

\courssub{Entraînement à l'examen : sujets probables}

\begin{exercice}[Sujet d'entraînement complet (30 minutes)]
Sujet : « Tes réseaux sociaux préférés » — Rédigez un texte de 6 à 10 phrases en chinois. Consignes : utilisez au moins 3 connecteurs (\zh{首先、因为、但是、所以}), au moins 5 mots du lexique du thème (\zh{微信、聊天、视频、网络、手机}), et respectez la gestion du temps 5 + 20 + 5. Notez-vous ensuite vous-même avec la grille sur 20.
\end{exercice}

\begin{corrige}[Sujet d'entraînement]
Modèle de copie attendue (une possibilité parmi d'autres) :\\
\zh{首先，我最喜欢微信，因为我每天用它和朋友聊天。} (Shǒuxiān, wǒ zuì xǐhuan Wēixìn, yīnwèi wǒ měitiān yòng tā hé péngyou liáotiān. — « D'abord, je préfère WeChat, parce que je l'utilise chaque jour pour discuter avec mes amis. »)\\
\zh{我也喜欢看抖音的视频，很有意思。} (Wǒ yě xǐhuan kàn Dǒuyīn de shìpín, hěn yǒu yìsi. — « J'aime aussi regarder les vidéos de Douyin, c'est très intéressant. »)\\
\zh{但是，花很多时间上网不好，对眼睛和学习都不好。} (Dànshì, huā hěn duō shíjiān shàngwǎng bù hǎo, duì yǎnjing hé xuéxí dōu bù hǎo. — « Mais passer beaucoup de temps en ligne n'est pas bon, ni pour les yeux ni pour les études. »)\\
\zh{所以，我们应该少用手机，多运动、多学习。} (Suǒyǐ, wǒmen yīnggāi shǎo yòng shǒujī, duō yùndòng, duō xuéxí. — « Donc nous devons moins utiliser le téléphone, faire plus de sport et plus étudier. »)\\
Barème indicatif : caractères corrects 4/4 ; ordre des mots et adjectifs sans \zh{是} 6/6 ; connecteurs \zh{首先、因为、但是、所以} 4/4 ; sujet respecté, idées personnelles 4/4 ; ponctuation chinoise 2/2. Pour atteindre les 6 phrases, on peut ajouter par exemple : \zh{周末，我也和家人视频聊天。} (Zhōumò, wǒ yě hé jiārén shìpín liáotiān. — « Le week-end, je fais aussi des appels vidéo avec ma famille. ») et \zh{网络很有用，但是要小心。} (Wǎngluò hěn yǒuyòng, dànshì yào xiǎoxīn. — « Internet est très utile, mais il faut être prudent. »)
\end{corrige}

\begin{exercice}[Sujets probables complémentaires]
Entraînez-vous avec les mêmes consignes sur ces deux sujets : 1) « Les dangers d'Internet » (utilisez \zh{危险、应该、少、但是}) ; 2) « Un message à un ami chinois » (formule d'appel \zh{亲爱的} (qīn'ài de, « cher/chère ») + nom, formule finale \zh{祝你好运！} (Zhù nǐ hǎoyùn! — « Bonne chance ! »)).
\end{exercice}

\begin{savaistu}[Les réseaux sociaux, un thème très probable au Bac]
Au Cameroun, les sujets d'expression écrite de chinois au baccalauréat portent très souvent sur la vie quotidienne des jeunes : l'école, la famille, le sport, le téléphone et les réseaux sociaux. Le thème « réseaux sociaux et Internet » fait partie des sujets les plus probables des prochaines sessions, car il combine vie quotidienne et vocabulaire du module « Mass médias et télécommunication ». En Chine, WeChat (\zh{微信}) compte plus d'un milliard d'utilisateurs et sert à tout : discuter, payer, commander un taxi — bien plus que nos applications occidentales. Préparer ce thème, c'est donc préparer à la fois l'examen et la communication réelle avec des Chinois.
\end{savaistu}

\begin{aretenir}[L'essentiel de la section]
\begin{itemize}
\item La copie est notée sur 20 : caractères 4, grammaire 6, connecteurs 4, contenu 4, présentation 2.
\item Gestion du temps : 5 min de plan, 20 min de rédaction, 5 min de relecture.
\item Relecture en 3 passages : sens général, grammaire, caractères trait par trait.
\item 6 à 10 phrases simples et justes valent mieux que des phrases complexes fautives.
\item Pièges à éliminer : temps avant le verbe, pas de \zh{是} devant l'adjectif, caractères proches \zh{观/欢}.
\end{itemize}
\end{aretenir}

\newpage

\coursec{Activité d'intégration}

\courssub{Objectifs de l'activité}

À l'issue de cette activité d'intégration, tu seras capable de :
\begin{itemize}
\item rédiger un texte simple, structuré et correct de 8 à 10 phrases sur le thème des réseaux sociaux ;
\item écrire sans faute au moins 8 caractères complexes du lexique de la leçon (\zh{社交媒体}, \zh{微信}, \zh{微博}, \zh{点赞}, \zh{评论}, \zh{关注}, \zh{转发}, \zh{隐私}…) ;
\item utiliser au moins 3 connecteurs différents (\zh{首先}, \zh{而且}, \zh{但是}, \zh{所以}, \zh{我觉得}) ;
\item mobiliser les structures étudiées : \zh{用 + plateforme + 来 + verbe}, \zh{虽然……但是……}, \zh{我觉得 + phrase}, \zh{我们应该 + verbe} ;
\item corriger le texte d'un camarade à l'aide d'une grille d'évaluation (caractères, grammaire, connecteurs, contenu) ;
\item traduire des phrases simples dans les deux sens (thème et version).
\end{itemize}

\courssub{Situation}

Le club de chinois de ton lycée de Yaoundé prépare un échange en ligne avec une classe partenaire de Chengdu (成都), en Chine. Les élèves chinois ont envoyé un message sur le groupe de discussion du club : ils veulent connaître les habitudes des jeunes Camerounais sur les réseaux sociaux. Le professeur te demande de répondre par écrit, en chinois, pour le « mur des publications » (\zh{朋友圈}) de la classe.

\begin{textebox}[Message reçu de la classe de Chengdu]
你们好！我们是成都的学生。\\\\
\emph{Nǐmen hǎo! Wǒmen shì Chéngdū de xuésheng.}\\\\
« Bonjour ! Nous sommes des élèves de Chengdu. »\\\\
我们想知道：喀麦隆的学生常用什么社交媒体？\\\\
\emph{Wǒmen xiǎng zhīdao : Kāmàilóng de xuésheng chángyòng shénme shèjiāo méitǐ?}\\\\
« Nous voulons savoir : quels réseaux sociaux les élèves camerounais utilisent-ils souvent ? »\\\\
你们在网上做什么？社交媒体有什么好处和问题？\\\\
\emph{Nǐmen zài wǎngshang zuò shénme? Shèjiāo méitǐ yǒu shénme hǎochù hé wèntí?}\\\\
« Que faites-vous en ligne ? Quels sont les avantages et les problèmes des réseaux sociaux ? »\\\\
请写八到十个句子，告诉我们你的想法！\\\\
\emph{Qǐng xiě bā dào shí ge jùzi, gàosu wǒmen nǐ de xiǎngfa!}\\\\
« Écrivez huit à dix phrases et dites-nous ce que vous en pensez ! »
\end{textebox}

Ta mission : rédiger un texte intitulé \zh{我和社交媒体} (\emph{Wǒ hé shèjiāo méitǐ}, « Moi et les réseaux sociaux ») qui répond aux trois questions des élèves chinois : les plateformes utilisées, les actions habituelles en ligne, un avantage et un inconvénient, avec une conclusion personnelle.

\courssub{Consignes}

\begin{enumerate}
\item \textbf{Lecture et compréhension.} Lis le message des élèves de Chengdu. Souligne les trois questions posées et traduis-les en français à l'écrit. Repère dans le message les caractères du lexique de la leçon (\zh{社交媒体}, \zh{知道}, \zh{问题}).
\item \textbf{Préparation du lexique.} Recopie sur ton brouillon la liste des mots que tu comptes utiliser : au moins 8 mots du thème, écrits en caractères, avec le pinyin. Vérifie l'ordre des traits des caractères complexes (\zh{微}, \zh{博}, \zh{赞}, \zh{隐}) dans le tableau de la leçon.
\item \textbf{Plan du texte.} Construis ton plan en quatre parties : (a) présentation des plateformes que tu utilises ; (b) tes actions habituelles (\zh{发照片}, \zh{点赞}, \zh{评论}, \zh{关注}) ; (c) un avantage et un inconvénient ; (d) conclusion avec \zh{我觉得} ou \zh{我们应该}.
\item \textbf{Rédaction.} Écris ton texte de 8 à 10 phrases, en caractères simplifiés, avec au moins 3 connecteurs différents. Soigne la ponctuation chinoise pleine chasse (，。？).
\item \textbf{Relecture.} Vérifie chaque caractère (traits, radicaux, caractères proches), la place des compléments et l'usage de \zh{了}.
\item \textbf{Correction croisée.} Échange ton texte avec ton binôme et corrige-le à l'aide de la grille d'évaluation (voir Ressources). Note sur 20 et écris un conseil d'amélioration en français.
\item \textbf{Publication.} Les meilleurs textes sont affichés au mur de la classe, rebaptisé \zh{朋友圈} (\emph{péngyouquān}, « cercle d'amis »), et lus à voix haute.
\end{enumerate}

\courssub{Ressources à mobiliser}

\textbf{Lexique du thème.}

\begin{center}\begin{tabularx}{\linewidth}{|X|X|X|X|}\hline
\rowcolor{popblueL} Caractères & Pinyin & Nature & Traduction \\ \hline
社交媒体 & shèjiāo méitǐ & nom & réseaux sociaux \\ \hline
微信 & Wēixìn & nom propre & WeChat \\ \hline
微博 & Wēibó & nom propre & Weibo \\ \hline
发照片 & fā zhàopiàn & verbe + nom & publier des photos \\ \hline
点赞 & diǎnzàn & verbe & liker \\ \hline
评论 & pínglùn & verbe / nom & commenter ; commentaire \\ \hline
关注 & guānzhù & verbe & suivre (un compte) \\ \hline
转发 & zhuǎnfā & verbe & partager, republier \\ \hline
账号 & zhànghào & nom & compte (en ligne) \\ \hline
隐私 & yǐnsī & nom & vie privée \\ \hline
保护 & bǎohù & verbe & protéger \\ \hline
浪费时间 & làngfèi shíjiān & locution & perdre du temps \\ \hline
\end{tabularx}\end{center}

\textbf{Écriture des caractères complexes : radicaux et ordre des traits.}

\begin{itemize}
\item \zh{微} (\emph{wēi}, 13 traits) : clé \zh{彳} (pas, à gauche). On écrit d'abord \zh{彳} (gauche), puis la partie centrale (\zh{山} surmonté d'un trait horizontal, puis \zh{几}), puis \zh{攵} (droite). Ordre général : gauche → centre → droite.
\item \zh{博} (\emph{bó}, 12 traits) : clé \zh{十} (à gauche, forme étroite \zh{忄}-like non : c'est bien \zh{十}). On écrit \zh{十} à gauche, puis la partie droite \zh{甫} avec \zh{寸} en bas : de haut en bas, et le point de \zh{寸} en dernier.
\item \zh{赞} (\emph{zàn}, 16 traits) : caractère haut-bas. On écrit d'abord les deux \zh{先} côte à côte en haut (gauche puis droite), puis \zh{贝} en bas (clé « coquillage / argent »), le dernier point de \zh{贝} en fin.
\item \zh{隐} (\emph{yǐn}, 11 traits) : clé \zh{阝} (colline, à gauche, 2 traits : oreille puis vertical). On écrit \zh{阝} à gauche, puis la partie droite de haut en bas : \zh{⺈}, puis \zh{彐}, puis \zh{心} en bas.
\end{itemize}

\textbf{Caractères proches : ne plus se tromper.}

\begin{center}\begin{tabularx}{\linewidth}{|X|X|X|}\hline
\rowcolor{popblueL} Caractère de la leçon & Caractère proche & Astuce pour les distinguer \\ \hline
博 \emph{bó} (Weibo) & 薄 \emph{báo} (mince) & \zh{薄} a la clé « herbe » \zh{艹} en haut ; \zh{博} commence par \zh{十} à gauche. \\ \hline
微 \emph{wēi} (WeChat) & 徽 \emph{huī} (emblème) & \zh{徽} contient \zh{糸} (soie) au centre ; \zh{微} contient \zh{山} et \zh{几}. \\ \hline
赞 \emph{zàn} (liker) & 替 \emph{tì} (remplacer) & \zh{赞} a \zh{贝} (argent) en bas ; \zh{替} a \zh{曰} en bas. \\ \hline
隐 \emph{yǐn} (caché) & 稳 \emph{wěn} (stable) & \zh{隐} a la clé \zh{阝} à gauche ; \zh{稳} a la clé \zh{禾} (céréale). \\ \hline
\end{tabularx}\end{center}

\textbf{Structures utiles.}
\begin{itemize}
\item \zh{用 + plateforme + 来 + verbe} : \zh{我用微信来跟朋友聊天。} (\emph{Wǒ yòng Wēixìn lái gēn péngyou liáotiān.}) « J'utilise WeChat pour discuter avec mes amis. »
\item \zh{虽然……但是……} : \zh{虽然社交媒体很方便，但是有时候会浪费时间。} (\emph{Suīrán shèjiāo méitǐ hěn fāngbiàn, dànshì yǒu shíhou huì làngfèi shíjiān.}) « Bien que les réseaux sociaux soient pratiques, ils font parfois perdre du temps. »
\item \zh{我觉得 + phrase} : \zh{我觉得我们应该保护隐私。} (\emph{Wǒ juéde wǒmen yīnggāi bǎohù yǐnsī.}) « Je pense que nous devons protéger notre vie privée. »
\end{itemize}

\textbf{Connecteurs} : \zh{首先} (d'abord), \zh{然后} (ensuite), \zh{而且} (de plus), \zh{但是} (mais), \zh{所以} (donc), \zh{最后} (enfin).

\textbf{Grille d'évaluation (sur 20).} Caractères corrects : 5 points ; grammaire et ordre des mots : 5 points ; au moins 3 connecteurs : 3 points ; contenu complet (4 parties du plan) : 4 points ; ponctuation et présentation : 3 points.

\begin{attention}[Pièges à éviter]
N'oublie pas la partie centrale de \zh{微} (\zh{山} + \zh{几}) entre \zh{彳} et \zh{攵} ; ne confonds pas \zh{博} (\emph{bó}, Weibo) et \zh{薄} (\emph{báo}, mince) ; \zh{赞} s'écrit avec deux \zh{先} au-dessus de \zh{贝}. Attention aussi à la place du complément de temps : on dit \zh{我每天用微信} et non \zh{我用微信每天}. Enfin, dans \zh{很快地得到新闻}, c'est bien \zh{地} (complément de manière devant le verbe) et non \zh{得}.
\end{attention}

\courssub{Proposition de production et corrigé commenté}

\begin{exoresolu}[Modèle de texte : 我和社交媒体]
Rédige un texte de 8 à 10 phrases répondant au message de la classe de Chengdu, avec au moins 3 connecteurs et 8 mots du lexique de la leçon.
\tcblower
\textbf{Production modèle}\\\\
我和社交媒体\\\\
\emph{Wǒ hé shèjiāo méitǐ}\\\\
« Moi et les réseaux sociaux »\\\\
首先，我每天都用社交媒体，比如微信和微博。\\\\
\emph{Shǒuxiān, wǒ měitiān dōu yòng shèjiāo méitǐ, bǐrú Wēixìn hé Wēibó.}\\\\
« D'abord, j'utilise chaque jour les réseaux sociaux, par exemple WeChat et Weibo. »\\\\
我用微信来跟家人和朋友聊天。\\\\
\emph{Wǒ yòng Wēixìn lái gēn jiārén hé péngyou liáotiān.}\\\\
« J'utilise WeChat pour discuter avec ma famille et mes amis. »\\\\
我也喜欢在网上发照片，比如雅温得的风景。\\\\
\emph{Wǒ yě xǐhuan zài wǎngshang fā zhàopiàn, bǐrú Yǎwēndé de fēngjǐng.}\\\\
« J'aime aussi publier des photos en ligne, par exemple des paysages de Yaoundé. »\\\\
看到朋友的照片，我常常点赞和评论。\\\\
\emph{Kàndào péngyou de zhàopiàn, wǒ chángcháng diǎnzàn hé pínglùn.}\\\\
« Quand je vois les photos de mes amis, je like et je commente souvent. »\\\\
而且，我关注了几个学习中文的账号。\\\\
\emph{Érqiě, wǒ guānzhù le jǐ ge xuéxí Zhōngwén de zhànghào.}\\\\
« De plus, je suis quelques comptes pour apprendre le chinois. »\\\\
社交媒体有很多好处：我们可以很快地得到新闻。\\\\
\emph{Shèjiāo méitǐ yǒu hěn duō hǎochù : wǒmen kěyǐ hěn kuài de dédào xīnwén.}\\\\
« Les réseaux sociaux ont beaucoup d'avantages : nous pouvons obtenir les nouvelles très vite. »\\\\
但是，有时候它们让我们浪费时间。\\\\
\emph{Dànshì, yǒu shíhou tāmen ràng wǒmen làngfèi shíjiān.}\\\\
« Mais parfois, ils nous font perdre du temps. »\\\\
所以，我觉得我们应该少用手机，多学习。\\\\
\emph{Suǒyǐ, wǒ juéde wǒmen yīnggāi shǎo yòng shǒujī, duō xuéxí.}\\\\
« Donc, je pense que nous devons moins utiliser le téléphone portable et plus étudier. »\\\\
最后，我们也要保护自己的隐私。\\\\
\emph{Zuìhòu, wǒmen yě yào bǎohù zìjǐ de yǐnsī.}\\\\
« Enfin, nous devons aussi protéger notre vie privée. »\\\\
\textbf{Commentaire des choix de langue.} Le texte compte 9 phrases et respecte le plan en quatre parties. Il contient 5 connecteurs (\zh{首先}, \zh{而且}, \zh{但是}, \zh{所以}, \zh{最后}), plus de 10 mots du lexique (\zh{社交媒体}, \zh{微信}, \zh{微博}, \zh{发照片}, \zh{点赞}, \zh{评论}, \zh{关注}, \zh{账号}, \zh{浪费时间}, \zh{隐私}), la structure \zh{用……来……} et la conclusion personnelle avec \zh{我觉得我们应该}. Le complément de temps \zh{每天} est bien placé avant le verbe, et \zh{了} dans \zh{关注了几个账号} marque une action accomplie.
\end{exoresolu}

\courssub{Entraînement à la traduction : thème et version}

\begin{exoresolu}[Thème : français → chinois]
Traduis en chinois : « J'utilise Weibo pour suivre les nouvelles, mais je pense que nous devons protéger notre vie privée. »
\tcblower
\zh{我用微博来关注新闻，但是我觉得我们应该保护自己的隐私。}\\\\
\emph{Wǒ yòng Wēibó lái guānzhù xīnwén, dànshì wǒ juéde wǒmen yīnggāi bǎohù zìjǐ de yǐnsī.}\\\\
Points de méthode : structure \zh{用……来……} pour « utiliser … pour … » ; le connecteur \zh{但是} relie les deux idées opposées ; \zh{自己的} renforce « notre (propre) vie privée ».
\end{exoresolu}

\begin{exoresolu}[Version : chinois → français]
Traduis en français : \zh{虽然社交媒体很方便，但是有时候它让我们浪费时间。}\\\\
\emph{Suīrán shèjiāo méitǐ hěn fāngbiàn, dànshì yǒu shíhou tā ràng wǒmen làngfèi shíjiān.}
\tcblower
« Bien que les réseaux sociaux soient pratiques, ils nous font parfois perdre du temps. »\\\\
Points de méthode : \zh{虽然……但是……} se rend par « bien que … (mais) … » — en français, on ne répète pas « mais » ; \zh{让 + quelqu'un + verbe} exprime « faire faire quelque chose à quelqu'un ».
\end{exoresolu}

\begin{exercice}[À toi de traduire]
1. Thème : traduis en chinois : « Mes amis aiment publier des photos en ligne, et je like souvent leurs photos. »\\
2. Version : traduis en français : \zh{我每天都用微信跟家人聊天。} (\emph{Wǒ měitiān dōu yòng Wēixìn gēn jiārén liáotiān.})
\end{exercice}

\begin{corrige}[Exercice « À toi de traduire »]
1. \zh{我的朋友们喜欢在网上发照片，我常常给他们的照片点赞。} (\emph{Wǒ de péngyoumen xǐhuan zài wǎngshang fā zhàopiàn, wǒ chángcháng gěi tāmen de zhàopiàn diǎnzàn.}) — On utilise \zh{给……点赞} (« liker les photos de quelqu'un ») ; le complément de lieu \zh{在网上} se place avant le verbe.\\
2. « Chaque jour, j'utilise WeChat pour discuter avec ma famille. » — \zh{每天} est placé avant le verbe en chinois, mais se traduit naturellement en tête ou en fin de phrase française.
\end{corrige}

\courssub{Bilan de l'activité}

Cette activité t'a permis de mobiliser en situation réelle tous les savoir-faire de la leçon : comprendre un message authentique en chinois, choisir et écrire correctement les caractères complexes du thème des réseaux sociaux, organiser un texte en quatre parties grâce aux connecteurs, et produire une opinion personnelle avec \zh{我觉得} et \zh{我们应该}. La correction croisée t'a aussi entraîné à repérer les erreurs typiques des francophones : ordre des mots, place du complément de temps, confusion entre \zh{地} et \zh{得}, caractères proches (\zh{博} / \zh{薄}, \zh{微} / \zh{徽}, \zh{赞} / \zh{替}). Pour t'autoévaluer, vérifie que tu sais : écrire de mémoire les 12 mots du tableau, rédiger sans modèle un texte de 8 phrases, et traduire dans les deux sens une phrase comme \zh{我觉得我们应该保护隐私。} « Je pense que nous devons protéger notre vie privée. » Ces compétences seront réinvesties à l'examen, dans l'épreuve d'expression écrite et de traduction.

\newpage

\coursec{Méthodes et exercices résolus}

\courssub{Savoir-faire 1 : Écrire correctement les caractères complexes du thème}

\begin{methode}[Écrire un caractère complexe sans faute]
Pour écrire correctement les caractères du vocabulaire des réseaux sociaux (网, 络, 媒, 博, 微, 信, 赞, 评论…), applique la démarche suivante :
\begin{enumerate}
\item \cle{Identifier le radical} (clé) : il donne le sens et se place souvent à gauche ou en haut. Exemples : 纟(soie) dans 网络, 女 (femme) dans 媒, 讠(parole) dans 评论.
\item \cle{Décomposer le caractère} en composants : 媒 = 女 + 某 ; 博 = 十 + 尃 ; 赞 = 兟 + 贝.
\item \cle{Respecter l'ordre des traits} : de haut en bas, de gauche à droite, l'horizontal avant le vertical, l'extérieur avant l'intérieur, on ferme la boîte en dernier (comme dans 网 : on écrit le cadre, puis l'intérieur, puis on ferme).
\item \cle{Vérifier les proportions} : dans un caractère gauche-droite comme 信 ou 评, la partie gauche (亻, 讠) est étroite, la partie droite plus large.
\item \cle{Réécrire de mémoire} trois fois, en prononçant le pinyin à voix haute : mémoire de la main + mémoire de l'oreille.
\end{enumerate}
Réflexe d'examen : relis chaque caractère écrit et demande-toi « le radical est-il le bon ? le composant phonétique est-il complet ? ».
\end{methode}

\begin{exoresolu}[Écriture guidée des caractères]
Énoncé. Écris en caractères chinois les mots suivants, en indiquant le radical de chacun : réseau (wǎng), média (méi), microblog (wēibó), commentaire (pínglùn).
\tcblower
Solution détaillée.
\begin{enumerate}
\item \zh{网} (wǎng, « filet, réseau ») : radical 冂 (cadre). On trace d'abord le cadre extérieur, puis les deux croix intérieures, puis on ferme en bas. Attention à ne pas écrire 冈 (colline) qui n'a qu'une seule croix.
\item \zh{媒} (méi, « intermédiaire, média ») : radical 女 (femme) à gauche, étroit ; à droite 某 (甘 + 木). Ordre : 女 d'abord, puis 某 de haut en bas.
\item \zh{微博} (wēibó, « microblog ») : 微 a le radical 彳 (pas, à gauche) ; 博 a le radical 十, avec 尃 à droite — ne pas oublier le petit trait (点) en bas à droite de 寸.
\item \zh{评论} (pínglùn, « commenter ») : les deux caractères portent le radical de la parole 讠 à gauche (forme simplifiée de 言) ; à droite 平 et 仑. Le radical 讠 s'écrit en deux traits : le point, puis le pli horizontal-vertical.
\end{enumerate}
Conclusion : un caractère juste = bon radical + bon composant + ordre des traits respecté. Au Bac, un caractère mal tracé (trait manquant, radical confondu) est compté faux.
\end{exoresolu}

\courssub{Savoir-faire 2 : Distinguer les caractères proches}

\begin{methode}[Ne plus confondre les caractères proches]
\begin{enumerate}
\item \cle{Repérer la différence visuelle} : un trait de plus ou de moins change le sens. Exemples : 网 / 冈 ; 博 / 搏 (搏 = se battre, radical 扌main) ; 评 / 苹 (苹 de 苹果, radical 艹 herbe).
\item \cle{Utiliser le radical comme indice de sens} : 讠→ parole (评论), 扌 → action de la main (拍 photo, 点赞), 纟 → fil, réseau (网络).
\item \cle{Créer une phrase-repère} qui contient le mot correct : \zh{我在网上评论。} (Wǒ zài wǎngshang pínglùn. « Je commente sur Internet. »)
\item \cle{Faire une dictée croisée} : un camarade dicte le pinyin, tu écris le caractère, puis tu vérifies dans le tableau.
\end{enumerate}
\end{methode}

\begin{exoresolu}[Choisir le bon caractère]
Énoncé. Complète avec le bon caractère : (a) 网 / 冈 : 上\_\_\_ (se connecter) ; (b) 博 / 搏 : 微\_\_\_ (microblog) ; (c) 评 / 苹 : \_\_\_论 (commentaire) ; (d) 信 / 言 : 微\_\_\_ (WeChat).
\tcblower
Solution détaillée.
\begin{enumerate}
\item (a) \zh{上网} (shàngwǎng, « aller sur Internet ») : 网 = filet/réseau ; 冈 (gāng) = colline, sens impossible ici.
\item (b) \zh{微博} (wēibó) : 博 = vaste, abondant ; 搏 (bó) = lutter, avec le radical 扌 de la main — un microblog ne « se bat » pas.
\item (c) \zh{评论} (pínglùn) : 评 porte le radical 讠 (parole), cohérent avec « commenter » ; 苹 n'existe que dans 苹果 (píngguǒ, « pomme »).
\item (d) \zh{微信} (Wēixìn, « WeChat ») : 信 = message, confiance (亻+ 言) ; 言 seul signifie « parole » mais ne forme pas ce mot.
\end{enumerate}
Conclusion : la logique du radical (parole, main, soie) est ta meilleure alliée pour choisir le bon caractère.
\end{exoresolu}

\courssub{Savoir-faire 3 : Produire un texte simple sur les réseaux sociaux}

\begin{methode}[Rédiger un court texte sur les réseaux sociaux]
\begin{enumerate}
\item \cle{Planifier en 3 parties} : présentation (quel réseau, qui l'utilise) → activités (ce qu'on y fait) → avis personnel (avantages / limites).
\item \cle{Utiliser les connecteurs} : \zh{首先} (shǒuxiān, « d'abord »), \zh{然后} (ránhòu, « ensuite »), \zh{但是} (dànshì, « mais »), \zh{因为} (yīnwèi, « parce que »), \zh{所以} (suǒyǐ, « donc »).
\item \cle{Employer les structures du thème} : Sujet + 用 + réseau + 做 + action (\zh{我用微信和朋友聊天。}) ; Sujet + 在 + réseau + 上 + verbe (\zh{他在微博上发照片。}).
\item \cle{Donner un avis} avec \zh{我觉得} (wǒ juéde, « je pense que ») + phrase, ou \zh{对……有好处/坏处} (duì … yǒu hǎochù / huàichù, « être bon / mauvais pour »).
\item \cle{Relire} : ordre Sujet–Temps–Lieu–Verbe, un connecteur par phrase, caractères vérifiés.
\end{enumerate}
\end{methode}

\begin{exoresolu}[Rédaction guidée]
Énoncé. Rédige un texte de 4 à 6 phrases en chinois : « Tu parles de ton utilisation des réseaux sociaux (réseau utilisé, activités, avis) ».
\tcblower
Solution détaillée. Plan : présentation → activités → avis. Structures mobilisées : 用 + réseau ; 在 … 上 + verbe ; 我觉得 ; connecteurs 每天, 但是, 因为.

Texte rédigé :

\zh{我每天用微信和朋友聊天。}(Wǒ měitiān yòng Wēixìn hé péngyou liáotiān. « Chaque jour, j'utilise WeChat pour discuter avec mes amis. »)

\zh{我也在微博上发照片、看新闻。}(Wǒ yě zài wēibó shang fā zhàopiàn, kàn xīnwén. « Je publie aussi des photos et je lis les nouvelles sur Weibo. »)

\zh{但是我不花太多时间上网。}(Dànshì wǒ bù huā tài duō shíjiān shàngwǎng. « Mais je ne passe pas trop de temps sur Internet. »)

\zh{我觉得网络很有用，因为它帮助我学习汉语。}(Wǒ juéde wǎngluò hěn yǒuyòng, yīnwèi tā bāngzhù wǒ xuéxí Hànyǔ. « Je pense qu'Internet est très utile, parce qu'il m'aide à apprendre le chinois. »)

Justifications : le complément de temps 每天 se place avant le verbe ; 在 … 上 encadre le nom du réseau ; 因为 introduit la cause en fin de phrase. Conclusion : 4 phrases, 3 connecteurs, zéro faute de caractère = copie réussie.
\end{exoresolu}

\courssub{Savoir-faire 4 : Traduire — thème (français → chinois)}

\begin{methode}[Réussir le thème]
\begin{enumerate}
\item \cle{Analyser la phrase française} : qui ? quand ? où ? fait quoi ?
\item \cle{Poser le squelette chinois} : Sujet + (temps) + (lieu avec 在) + verbe + complément.
\item \cle{Choisir le mot-outil exact} : 在 devant le lieu, 用 pour l'outil (« avec »), 和 pour « avec » une personne, 了 pour une action accomplie.
\item \cle{Traduire mot à mot le vocabulaire du thème} : réseaux sociaux = \zh{社交网络} (shèjiāo wǎngluò), publier = \zh{发} (fā), aimer/liker = \zh{点赞} (diǎnzàn), partager = \zh{分享} (fēnxiǎng).
\item \cle{Relire en chinois uniquement} : la phrase doit être naturelle pour un Chinois, pas un calque du français.
\end{enumerate}
\end{methode}

\begin{exoresolu}[Thème : phrases à traduire]
Énoncé. Traduis en chinois : (a) « Ma sœur publie des photos sur les réseaux sociaux. » (b) « Hier, j'ai partagé une vidéo avec mes camarades de classe. »
\tcblower
Solution détaillée.

(a) Sujet : 我姐姐 (wǒ jiějie, « ma grande sœur »). Lieu : 在社交网络上 (zài shèjiāo wǎngluò shang, « sur les réseaux sociaux ») placé avant le verbe. Verbe : 发照片 (fā zhàopiàn, « publier des photos »).

Réponse : \zh{我姐姐在社交网络上发照片。}(Wǒ jiějie zài shèjiāo wǎngluò shang fā zhàopiàn.)

Piège évité : en français le lieu est après le verbe ; en chinois, 在 + lieu se place AVANT le verbe.

(b) Temps : 昨天 (zuótiān, « hier ») après le sujet. Action accomplie : 了 après le verbe 分享. « Avec mes camarades » : 和同学们, placé avant le verbe ou après 了 selon le sens ; ici : 和同学们一起 (« ensemble avec »).

Réponse : \zh{昨天我和同学们一起分享了一个视频。}(Zuótiān wǒ hé tóngxuémen yìqǐ fēnxiǎng le yí ge shìpín.)

Justifications : 了 marque l'accompli (hier) ; le classificateur de 视频 est 个 ; 一起 renforce l'idée de « avec ». Conclusion : squelette Sujet + Temps + 和 + personne + Verbe + 了 + complément.
\end{exoresolu}

\courssub{Savoir-faire 5 : Traduire — version (chinois → français)}

\begin{methode}[Réussir la version]
\begin{enumerate}
\item \cle{Lire la phrase entière} avant de traduire ; repérer sujet, verbe, compléments.
\item \cle{Découper en blocs} : Sujet | temps/lieu | verbe | objet.
\item \cle{Traduire les mots-outils correctement} : 在 = « sur/dans/à », 用 = « avec (outil) », 对……感兴趣 = « s'intéresser à », 越来越 = « de plus en plus ».
\item \cle{Rédiger un français correct et naturel}, pas un mot-à-mot : \zh{上网} ne se traduit pas « monter le filet » mais « aller sur Internet ».
\item \cle{Vérifier la cohérence} : temps du verbe français (了 → passé), nombres, négation (不 présent/futur, 没 passé).
\end{enumerate}
\end{methode}

\begin{exoresolu}[Version : texte à traduire]
Énoncé. Traduis en français : \zh{现在越来越多的年轻人喜欢用手机上网。他们在微信上和朋友们聊天，在微博上看新闻。}(Xiànzài yuè lái yuè duō de niánqīngrén xǐhuan yòng shǒujī shàngwǎng. Tāmen zài Wēixìn shang hé péngyoumen liáotiān, zài wēibó shang kàn xīnwén.)
\tcblower
Solution détaillée. Découpage en blocs :

Bloc 1 : 现在 (maintenant) | 越来越多的年轻人 (de plus en plus de jeunes) | 喜欢 (aiment) | 用手机上网 (utiliser le téléphone pour aller sur Internet).

Bloc 2 : 他们 (ils) | 在微信上 (sur WeChat) | 和朋友们聊天 (discutent avec leurs amis) ; 在微博上 (sur Weibo) | 看新闻 (lisent les nouvelles).

Traduction rédigée : « Aujourd'hui, de plus en plus de jeunes aiment utiliser leur téléphone portable pour aller sur Internet. Ils discutent avec leurs amis sur WeChat et lisent les nouvelles sur Weibo. »

Justifications : 越来越 + adjectif = « de plus en plus » ; 用 + outil = « avec / au moyen de » ; 在 + plateforme + 上 = « sur ». Le verbe 喜欢 + verbe se traduit « aimer + infinitif ». Conclusion : une version réussie est un français fluide qui respecte exactement le sens du chinois.
\end{exoresolu}

\begin{attention}[Les 5 erreurs qui coûtent des points au Bac]
\begin{enumerate}
\item \cle{Tons négligés} : écrire \emph{wang} sans ton ou confondre wǎng (网, réseau) et wàng (忘, oublier) fausse le sens. Le pinyin avec tons est exigé dans les réponses en transcription.
\item \cle{Caractères confondus} : 博 / 搏, 评 / 苹, 网 / 冈, 信 / 言. Un trait ou un radical faux = caractère faux = point perdu. Vérifie toujours le radical (讠 parole, 扌 main, 纟 soie).
\item \cle{Ordre des mots calqué sur le français} : ne dis jamais \zh{我发照片在社交网络上} ; le complément de lieu avec 在 se place AVANT le verbe : \zh{我在社交网络上发照片。}
\item \cle{Mots-outils oubliés ou mal choisis} : oublier 了 pour une action passée (\zh{昨天我分享了} et non \zh{昨天我分享}), confondre 和 (avec une personne) et 用 (avec un outil : \zh{用手机上网}).
\item \cle{Classificateur absent} : on écrit \zh{一个视频} (yí ge shìpín), \zh{一张照片} (yì zhāng zhàopiàn, 张 pour les objets plats), jamais le nom seul après un nombre.
\end{enumerate}
\end{attention}

\newpage

\coursec{Exercices}

\courssub{Je vérifie mes connaissances}

Avant de commencer les exercices, rappelle-toi des quatre paires de caractères proches qui reviennent souvent dans les épreuves. Observe bien les clés (radicaux) : c'est souvent la clé qui donne le sens.

\begin{aretenir}[Caractères proches : les quatre paires à ne pas confondre]
\begin{center}
\begin{tabularx}{\linewidth}{|X|X|X|X|}
\hline
\rowcolor{popblueL} Caractère & Pinyin & Clé & Sens \\
\hline
微 & wēi & 彳 (pas, cheminer) & petit, micro → 微信 WeChat \\
\hline
徽 & huī & 彳 & insigne, emblème (安徽 Ānhuī) \\
\hline
博 & bó & 十 & vaste, érudit → 博客 blog \\
\hline
搏 & bó & 扌(main) & combattre → 拼搏 lutter \\
\hline
络 & luò & 纟(soie) & relier, filet → 网络 réseau \\
\hline
路 & lù & 足 (pied) & route, chemin \\
\hline
账 & zhàng & 贝 (argent) & compte → 账号 numéro de compte \\
\hline
帐 & zhàng & 巾 (tissu) & tente, rideau → 帐篷 tente \\
\hline
\end{tabularx}
\end{center}
\end{aretenir}

\begin{exercice}[QCM : le bon caractère]
Pour chaque phrase, choisis le caractère qui convient parmi les deux proposés, puis recopie la phrase complète.
\begin{enumerate}
\item 我每天都在＿＿信上聊天。a) 微 \quad b) 徽
\item 他喜欢写＿＿客。a) 博 \quad b) 搏
\item 这个网＿＿很快。a) 络 \quad b) 路
\item 请告诉我你的＿＿号。a) 账 \quad b) 帐
\end{enumerate}
\end{exercice}

\begin{exercice}[Vrai ou faux — justifie ta réponse]
Dis si chaque affirmation est vraie ou fausse, et justifie en une phrase en français.
\begin{enumerate}
\item Le caractère \zh{网} (wǎng, « filet, réseau ») contient la clé 门 (porte).
\item Dans la phrase \zh{我每天用微信跟朋友聊天。}, le complément de moyen \zh{用微信} se place avant le verbe.
\item \zh{赞} (zàn) signifie « supprimer ».
\item Pour exprimer une opposition, on peut utiliser la structure \zh{虽然……但是……} (suīrán... dànshì...).
\end{enumerate}
\end{exercice}

\begin{exercice}[Vocabulaire : complète le tableau]
Recopie et complète le tableau suivant avec la nature du mot et sa traduction française.
\begin{center}
\begin{tabularx}{\linewidth}{|X|X|X|X|}
\hline
\rowcolor{popblueL} Caractères & Pinyin & Nature & Traduction \\
\hline
网络 & wǎngluò & \ldots & \ldots \\
\hline
微博 & Wēibó & \ldots & \ldots \\
\hline
粉丝 & fěnsī & \ldots & \ldots \\
\hline
删除 & shānchú & \ldots & \ldots \\
\hline
密码 & mìmǎ & \ldots & \ldots \\
\hline
视频 & shìpín & \ldots & \ldots \\
\hline
\end{tabularx}
\end{center}
\end{exercice}

\begin{exercice}[Associe caractères et pinyin]
Associe chaque caractère à son pinyin, puis à sa traduction.
\begin{center}
\begin{tabularx}{\linewidth}{|X|X|X|}
\hline
\rowcolor{popblueL} Caractère & Pinyin & Traduction \\
\hline
1. 网 & a. zhàng & A. supprimer \\
\hline
2. 博 & b. zàn & B. mot de passe \\
\hline
3. 赞 & c. wǎng & C. réseau, filet \\
\hline
4. 删 & d. mì & D. blog \\
\hline
5. 账 & e. bó & E. aimer, « liker » \\
\hline
6. 密 & f. shān & F. compte \\
\hline
\end{tabularx}
\end{center}
\end{exercice}

\courssub{Je m'entraîne}

\begin{exercice}[Remets les mots dans l'ordre]
Remets les mots dans le bon ordre pour former des phrases correctes, puis traduis-les en français.
\begin{enumerate}
\item 每天 / 我 / 微信 / 跟 / 聊天 / 朋友 / 用
\item 微博上 / 发 / 照片 / 她 / 在 / 喜欢
\item 虽然 / 很多 / 粉丝 / 但是 / 他 / 不 / 视频 / 的 / 有意思 / 很
\item 密码 / 我 / 忘了 / 账号 / 的
\end{enumerate}
\end{exercice}

\begin{textebox}[Mon compte Weibo — texte à compléter]
\zh{＿＿＿＿我很喜欢网络，＿＿＿＿我不花太多时间在网上。＿＿＿＿，我每天早上看微博上的新闻。＿＿＿＿，我跟朋友在微信上聊天。＿＿＿＿我要学习，＿＿＿＿晚上我不用手机。}\par
\emph{＿＿＿＿ wǒ hěn xǐhuan wǎngluò, ＿＿＿＿ wǒ bù huā tài duō shíjiān zài wǎng shàng. ＿＿＿＿, wǒ měitiān zǎoshang kàn Wēibó shàng de xīnwén. ＿＿＿＿, wǒ gēn péngyou zài Wēixìn shàng liáotiān. ＿＿＿＿ wǒ yào xuéxí, ＿＿＿＿ wǎnshang wǒ bú yòng shǒujī.}\par
« Bien que j'aime beaucoup Internet, je ne passe pas trop de temps en ligne. D'abord, chaque matin je lis les nouvelles sur Weibo. Ensuite, je discute avec mes amis sur WeChat. Comme je dois étudier, le soir je n'utilise pas mon téléphone. »
\end{textebox}

\begin{exercice}[Texte à trous : les connecteurs]
Complète le texte « Mon compte Weibo » ci-dessus avec les connecteurs de la boîte : \zh{首先}、\zh{然后}、\zh{因为……所以……}、\zh{虽然……但是……}（chaque connecteur n'est utilisé qu'une seule fois）. Aide-toi de la traduction française pour choisir le bon connecteur.
\end{exercice}

\begin{exercice}[Traduction : thème]
Traduis les phrases suivantes en chinois (caractères + pinyin).
\begin{enumerate}
\item J'utilise WeChat tous les jours pour discuter avec mes amis.
\item Elle a beaucoup de fans sur Weibo.
\item Ne donne ton mot de passe à personne.
\item Bien que cette vidéo soit intéressante, je n'ai pas le temps de la regarder.
\end{enumerate}
\end{exercice}

\begin{textebox}[Texte d'élève à corriger]
\zh{我很喜欢用徽信。每天我聊天跟朋友在网上。我的粉丝很多，所以我不喜欢分享照片。虽然网络很有意思，我不学习。昨天我发了一个视频，很多朋友都赞了它。}\par
\emph{Wǒ hěn xǐhuan yòng Huīxìn. Měitiān wǒ liáotiān gēn péngyou zài wǎng shàng. Wǒ de fěnsī hěn duō, suǒyǐ wǒ bù xǐhuan fēnxiǎng zhàopiàn. Suīrán wǎngluò hěn yǒu yìsi, wǒ bù xuéxí. Zuótiān wǒ fā le yí ge shìpín, hěn duō péngyou dōu zàn le tā.}
\end{textebox}

\begin{exercice}[Corrige les erreurs]
Le texte ci-dessus, écrit par un élève, contient \textbf{cinq erreurs} : un caractère faux, un problème d'ordre des mots, un connecteur mal choisi, un connecteur manquant et une faute de structure. Relève chaque erreur et écris la correction.
\end{exercice}

\courssub{Je me prépare au Bac}

\begin{textebox}[Les réseaux sociaux au lycée]
\zh{现在，很多喀麦隆中学生每天都用网络。他们用手机上微信、微博，看朋友发的照片和视频。虽然网络很方便，但是也有一些问题：有的学生花太多时间在网上，不学习；有的学生把密码告诉别人，账号就不安全了。老师说：“网络是好东西，但是我们要学会正确地用网络。”}\par
\emph{Xiànzài, hěn duō Kāmàilóng zhōngxuéshēng měitiān dōu yòng wǎngluò. Tāmen yòng shǒujī shàng Wēixìn, Wēibó, kàn péngyou fā de zhàopiàn hé shìpín. Suīrán wǎngluò hěn fāngbiàn, dànshì yě yǒu yìxiē wèntí : yǒu de xuéshēng huā tài duō shíjiān zài wǎng shàng, bù xuéxí ; yǒu de xuéshēng bǎ mìmǎ gàosu biérén, zhànghào jiù bù ānquán le. Lǎoshī shuō : « Wǎngluò shì hǎo dōngxi, dànshì wǒmen yào xuéhuì zhèngquè de yòng wǎngluò. »}\par
« Aujourd'hui, beaucoup de collégiens et lycéens camerounais utilisent Internet tous les jours. Avec leur téléphone, ils vont sur WeChat et Weibo, et regardent les photos et les vidéos publiées par leurs amis. Bien qu'Internet soit pratique, il y a aussi des problèmes : certains élèves passent trop de temps en ligne et n'étudient pas ; d'autres donnent leur mot de passe à d'autres personnes, et leur compte n'est plus en sécurité. Le professeur dit : “Internet est une bonne chose, mais nous devons apprendre à l'utiliser correctement.” »
\end{textebox}

\begin{exercice}[Compréhension de l'écrit (type Bac)]
Lis le texte « Les réseaux sociaux au lycée » ci-dessus, puis réponds aux questions \textbf{en français}.
\textbf{Questions de compréhension :}
\begin{enumerate}
\item Que font les élèves camerounais sur leur téléphone ?
\item Cite deux problèmes liés à l'utilisation des réseaux sociaux mentionnés dans le texte.
\item Quel conseil le professeur donne-t-il aux élèves ?
\end{enumerate}
\textbf{Questions de langue :}
\begin{enumerate}
\item Releve dans le texte une structure d'opposition et traduis-la.
\item Recopie les caractères \zh{网}、\zh{络}、\zh{账}、\zh{密} en respectant l'ordre des traits, et donne la clé (radical) de chacun.
\end{enumerate}
\end{exercice}

\begin{exercice}[Thème et version (type Bac)]
\textbf{A. Thème.} Traduis en chinois (caractères + pinyin) :
\begin{enumerate}
\item Bien que les réseaux sociaux soient pratiques, ils font perdre beaucoup de temps aux jeunes.
\item Mon frère aîné a ouvert un compte Weibo, mais il a oublié son mot de passe.
\end{enumerate}
\textbf{B. Version.} Traduis en français la phrase suivante :\\
\zh{我的粉丝不多，但是我每天还是在微博上发照片，因为我很喜欢分享我的生活。}\\
\emph{Wǒ de fěnsī bù duō, dànshì wǒ měitiān háishì zài Wēibó shàng fā zhàopiàn, yīnwèi wǒ hěn xǐhuan fēnxiǎng wǒ de shēnghuó.}
\end{exercice}

\begin{exercice}[Expression écrite (type Bac)]
\textbf{Sujet :} Raconte comment tu utilises les réseaux sociaux dans ta vie quotidienne (quels réseaux, pour quoi faire, combien de temps, les avantages et les dangers).
\begin{itemize}
\item Longueur attendue : \textbf{6 à 10 phrases}, soit environ \textbf{80 à 120 caractères chinois}.
\item Utilise \textbf{au moins 3 connecteurs} parmi : \zh{首先}、\zh{然后}、\zh{但是}、\zh{因为……所以……}、\zh{虽然……但是……}.
\item Utilise au moins 5 mots du vocabulaire du thème : \zh{网络}、\zh{微信}、\zh{微博}、\zh{粉丝}、\zh{视频}、\zh{赞}、\zh{删除}、\zh{账号}、\zh{密码}、\zh{分享}.
\item Soigne l'écriture des caractères : traits dans le bon ordre, caractères proches non confondus (\zh{微/徽}、\zh{博/搏}、\zh{络/路}、\zh{账/帐}).
\end{itemize}
\end{exercice}

\begin{methode}[Réussir ta production écrite en 5 étapes]
\begin{enumerate}
\item \textbf{Planifie} : note en français 3 ou 4 idées (quels réseaux ? quelle utilisation ? avantages ? dangers ?).
\item \textbf{Structure} : une phrase d'introduction, deux ou trois phrases de développement avec \zh{首先} / \zh{然后}, une phrase d'opposition avec \zh{虽然……但是……}, une conclusion avec \zh{因为……所以……}.
\item \textbf{Rédige} en phrases courtes : sujet + (temps / lieu / moyen) + verbe + complément. N'écris que ce que tu sais écrire correctement.
\item \textbf{Vérifie} les points sensibles : place des compléments avant le verbe, présence de \zh{但是} après \zh{虽然}, caractères proches (\zh{微} et non \zh{徽}).
\item \textbf{Compte} tes caractères (80 à 120) et tes connecteurs (au moins 3).
\end{enumerate}
\end{methode}

\begin{textebox}[Modèle de production (corrigé type)]
\zh{我每天都用网络。首先，我早上在微信上跟朋友聊天，看他们发的照片。然后，我在微博上发视频。虽然我的粉丝不多，但是很多朋友喜欢我的视频，给我点赞。因为我要学习，所以晚上我不用手机。网络很方便，但是我们要小心：不要把密码告诉别人。}\par
\emph{Wǒ měitiān dōu yòng wǎngluò. Shǒuxiān, wǒ zǎoshang zài Wēixìn shàng gēn péngyou liáotiān, kàn tāmen fā de zhàopiàn. Ránhòu, wǒ zài Wēibó shàng fā shìpín. Suīrán wǒ de fěnsī bù duō, dànshì hěn duō péngyou xǐhuan wǒ de shìpín, gěi wǒ diǎnzàn. Yīnwèi wǒ yào xuéxí, suǒyǐ wǎnshang wǒ bú yòng shǒujī. Wǎngluò hěn fāngbiàn, dànshì wǒmen yào xiǎoxīn : bú yào bǎ mìmǎ gàosu biérén.}\par
« J'utilise Internet tous les jours. D'abord, le matin, je discute avec mes amis sur WeChat et je regarde les photos qu'ils publient. Ensuite, je publie des vidéos sur Weibo. Bien que je n'aie pas beaucoup de fans, beaucoup d'amis aiment mes vidéos et les likent. Comme je dois étudier, le soir je n'utilise pas mon téléphone. Internet est pratique, mais nous devons être prudents : il ne faut donner son mot de passe à personne. »
\end{textebox}

\begin{attention}[Grille d'évaluation : ce que l'examinateur regarde]
\begin{itemize}
\item \textbf{Respect du sujet et de la consigne} : longueur (80--120 caractères), 3 connecteurs minimum, 5 mots du thème.
\item \textbf{Exactitude des caractères} : une confusion \zh{微/徽} ou \zh{账/帐} coûte des points ; révise les clés.
\item \textbf{Ordre des mots} : compléments de temps, de lieu, de moyen et de compagnie \emph{avant} le verbe.
\item \textbf{Connecteurs complets} : \zh{虽然} sans \zh{但是} est une faute ; \zh{因为} appelle \zh{所以}.
\item \textbf{Ponctuation chinoise} : utilise ，。？！：et non la ponctuation française dans le texte chinois.
\end{itemize}
\end{attention}

\newpage

\coursec{Corrigés des exercices}

\begin{corrige}[Exercice 1 — QCM : le bon caractère]
\begin{enumerate}
\item a) \zh{微} — \zh{我每天都在微信上聊天。} \emph{Wǒ měitiān dōu zài Wēixìn shàng liáotiān.} « Je discute tous les jours sur WeChat. » \zh{徽} (huī) signifie « insigne, emblème » (\zh{安徽} Ānhuī) : sens impossible ici.
\item a) \zh{博} — \zh{他喜欢写博客。} \emph{Tā xǐhuan xiě bókè.} « Il aime écrire un blog. » \zh{搏} (bó, clé de la main 扌) = « combattre » (\zh{拼搏} pīnbó).
\item a) \zh{络} — \zh{这个网络很快。} \emph{Zhège wǎngluò hěn kuài.} « Ce réseau est rapide. » \zh{路} (lù, clé du pied 足) = « route ».
\item a) \zh{账} — \zh{请告诉我你的账号。} \emph{Qǐng gàosu wǒ nǐ de zhànghào.} « Donne-moi ton numéro de compte, s'il te plaît. » La clé 贝 (argent) confirme le sens financier ; \zh{帐} (巾, tissu) = « tente ».
\end{enumerate}
\textbf{Méthode :} quand deux caractères se ressemblent, identifie d'abord la clé : elle donne le champ sémantique (soie 纟→ relier ; argent 贝 → compte ; main 扌→ action).
\end{corrige}

\begin{corrige}[Exercice 2 — Vrai ou faux]
\begin{enumerate}
\item \textbf{Faux.} \zh{网} (wǎng) est sa propre clé (on le classe aussi sous 冂) ; la clé 门 (mén, « porte ») n'apparaît pas dans ce caractère.
\item \textbf{Vrai.} Le complément de moyen \zh{用微信} se place avant le verbe : Sujet + \zh{用} + moyen + verbe → \zh{我用微信聊天。}
\item \textbf{Faux.} \zh{赞} (zàn) = « aimer, louer, liker » ; « supprimer » se dit \zh{删除} (shānchú).
\item \textbf{Vrai.} \zh{虽然……但是……} exprime la concession ; en chinois, contrairement au français, les deux connecteurs sont normalement présents dans la même phrase.
\end{enumerate}
\end{corrige}

\begin{corrige}[Exercice 3 — Vocabulaire : complète le tableau]
\begin{center}
\begin{tabularx}{\linewidth}{|X|X|X|X|}
\hline
\rowcolor{popblueL} Caractères & Pinyin & Nature & Traduction \\
\hline
网络 & wǎngluò & nom & réseau, Internet \\
\hline
微博 & Wēibó & nom propre & Weibo (réseau social chinois) \\
\hline
粉丝 & fěnsī & nom & fans, abonnés \\
\hline
删除 & shānchú & verbe & supprimer, effacer \\
\hline
密码 & mìmǎ & nom & mot de passe \\
\hline
视频 & shìpín & nom & vidéo \\
\hline
\end{tabularx}
\end{center}
\textbf{Remarque :} \zh{粉丝} (fěnsī) est un emprunt phonétique à l'anglais \emph{fans} ; au sens premier, ce mot désigne les « vermicelles de riz ».
\end{corrige}

\begin{corrige}[Exercice 4 — Associe caractères et pinyin]
1 -- c -- C : \zh{网} wǎng, « réseau, filet ».\\
2 -- e -- D : \zh{博} bó, « blog » (\zh{博客}).\\
3 -- b -- E : \zh{赞} zàn, « aimer, liker ».\\
4 -- f -- A : \zh{删} shān, « supprimer » (\zh{删除}).\\
5 -- a -- F : \zh{账} zhàng, « compte » (\zh{账号}).\\
6 -- d -- B : \zh{密} mì, « secret » → \zh{密码} « mot de passe ».
\end{corrige}

\begin{corrige}[Exercice 5 — Remets les mots dans l'ordre]
\begin{enumerate}
\item \zh{我每天用微信跟朋友聊天。} \emph{Wǒ měitiān yòng Wēixìn gēn péngyou liáotiān.} « Chaque jour, je discute avec mes amis sur WeChat. » Ordre : sujet + temps + moyen (\zh{用微信}) + compagnon (\zh{跟朋友}) + verbe.
\item \zh{她喜欢在微博上发照片。} \emph{Tā xǐhuan zài Wēibó shàng fā zhàopiàn.} « Elle aime publier des photos sur Weibo. » Le lieu \zh{在微博上} précède le verbe \zh{发}.
\item \zh{虽然他的视频很有意思，但是粉丝不很多。} \emph{Suīrán tā de shìpín hěn yǒu yìsi, dànshì fěnsī bù hěn duō.} « Bien que ses vidéos soient très intéressantes, il n'a pas beaucoup de fans. »
\item \zh{我忘了账号的密码。} \emph{Wǒ wàng le zhànghào de mìmǎ.} « J'ai oublié le mot de passe du compte. » \zh{的} marque la possession : le mot de passe \emph{du} compte.
\end{enumerate}
\textbf{Point de vigilance :} en chinois, les compléments de temps, de lieu, de moyen et de compagnie se placent \emph{avant} le verbe, jamais après comme en français.
\end{corrige}

\begin{corrige}[Exercice 6 — Texte à trous : les connecteurs]
\zh{虽然我很喜欢网络，但是我不花太多时间在网上。首先，我每天早上看微博上的新闻。然后，我跟朋友在微信上聊天。因为我要学习，所以晚上我不用手机。}\par
\emph{Suīrán wǒ hěn xǐhuan wǎngluò, dànshì wǒ bù huā tài duō shíjiān zài wǎng shàng. Shǒuxiān, wǒ měitiān zǎoshang kàn Wēibó shàng de xīnwén. Ránhòu, wǒ gēn péngyou zài Wēixìn shàng liáotiān. Yīnwèi wǒ yào xuéxí, suǒyǐ wǎnshang wǒ bú yòng shǒujī.}\par
\textbf{Justification :} « Bien que... » impose la paire \zh{虽然……但是……} ; « D'abord » = \zh{首先} ; « Ensuite » = \zh{然后} ; « Comme... » = \zh{因为……所以……}. Chaque paire de connecteurs doit être complète : \zh{因为} appelle \zh{所以}, \zh{虽然} appelle \zh{但是}.
\end{corrige}

\begin{corrige}[Exercice 7 — Traduction : thème]
\begin{enumerate}
\item \zh{我每天用微信跟朋友聊天。} \emph{Wǒ měitiān yòng Wēixìn gēn péngyou liáotiān.} (« pour discuter » est rendu directement par le verbe \zh{聊天} après le complément de moyen.)
\item \zh{她在微博上有很多粉丝。} \emph{Tā zài Wēibó shàng yǒu hěn duō fěnsī.} (Lieu \zh{在微博上} avant le verbe \zh{有}.)
\item \zh{不要把你的密码告诉别人。} \emph{Bú yào bǎ nǐ de mìmǎ gàosu biérén.} Impératif négatif : \zh{不要} + verbe ; la structure en \zh{把} (Sujet + 把 + objet + verbe + complément) met l'objet \zh{密码} en avant.
\item \zh{虽然这个视频很有意思，但是我没有时间看。} \emph{Suīrán zhège shìpín hěn yǒu yìsi, dànshì wǒ méiyǒu shíjiān kàn.} (« ne pas avoir le temps de » = \zh{没有时间} + verbe.)
\end{enumerate}
\end{corrige}

\begin{corrige}[Exercice 8 — Corrige les erreurs]
\begin{enumerate}
\item \textbf{Caractère faux :} \zh{徽信} → \zh{微信} (Wēixìn). C'est \zh{微} (wēi, « petit ») qu'il faut ; \zh{徽} (huī) = « insigne ».
\item \textbf{Ordre des mots :} \zh{每天我聊天跟朋友在网上} → \zh{每天我跟朋友在网上聊天。} Le compagnon \zh{跟朋友} et le lieu \zh{在网上} se placent avant le verbe \zh{聊天}.
\item \textbf{Connecteur mal choisi :} \zh{所以} → \zh{但是} : \zh{我的粉丝很多，但是我不喜欢分享照片。} La relation est une opposition, pas une conséquence.
\item \textbf{Connecteur manquant :} \zh{虽然网络很有意思，我不学习} → \zh{虽然网络很有意思，但是我不学习。} \zh{虽然} exige \zh{但是} dans la seconde proposition.
\item \textbf{Faute de structure :} \zh{很多朋友都赞了它} → \zh{很多朋友都给我点赞。} \zh{赞} ne prend pas de complément d'objet direct ; on dit \zh{给} + personne + \zh{点赞}.
\end{enumerate}
\textbf{Texte corrigé complet :} \zh{我很喜欢用微信。每天我跟朋友在网上聊天。我的粉丝很多，但是我不喜欢分享照片。虽然网络很有意思，但是我不学习。昨天我发了一个视频，很多朋友都给我点赞。}\par
\emph{Wǒ hěn xǐhuan yòng Wēixìn. Měitiān wǒ gēn péngyou zài wǎng shàng liáotiān. Wǒ de fěnsī hěn duō, dànshì wǒ bù xǐhuan fēnxiǎng zhàopiàn. Suīrán wǎngluò hěn yǒu yìsi, dànshì wǒ bù xuéxí. Zuótiān wǒ fā le yí ge shìpín, hěn duō péngyou dōu gěi wǒ diǎnzàn.}
\end{corrige}

\begin{corrige}[Exercice 9 — Compréhension de l'écrit (type Bac)]
\textbf{Questions de compréhension :}
\begin{enumerate}
\item Avec leur téléphone, ils vont sur WeChat et Weibo et regardent les photos et les vidéos publiées par leurs amis.
\item Deux problèmes : certains élèves passent trop de temps en ligne et n'étudient pas ; d'autres donnent leur mot de passe à d'autres personnes, et leur compte n'est plus en sécurité.
\item Le professeur conseille d'apprendre à utiliser Internet correctement, car Internet est une bonne chose.
\end{enumerate}
\textbf{Questions de langue :}
\begin{enumerate}
\item Structure d'opposition : \zh{虽然网络很方便，但是也有一些问题。} \emph{Suīrán wǎngluò hěn fāngbiàn, dànshì yě yǒu yìxiē wèntí.} « Bien qu'Internet soit pratique, il y a aussi des problèmes. » (Concession : \zh{虽然……但是……}.)
\item Clés et traits : \zh{网} → clé 冂 (le caractère est aussi sa propre clé) : on trace le cadre extérieur, l'intérieur, puis on ferme ; \zh{络} → clé 纟(soie), la clé à gauche avant la partie \zh{各} à droite ; \zh{账} → clé 贝 (coquillage, argent), 贝 à gauche puis \zh{长} à droite ; \zh{密} → clé 宀 (toit), de haut en bas : 宀 puis \zh{必} puis \zh{山}. Règle générale : de haut en bas, de gauche à droite, l'extérieur avant l'intérieur, on ferme les cadres en dernier.
\end{enumerate}
\textbf{Barème indicatif (sur 10) :} compréhension 3 × 2 pts = 6 pts ; question de langue 1 : 2 pts ; question de langue 2 : 2 pts (0,5 pt par clé correcte).
\end{corrige}

\begin{corrige}[Exercice 10 — Thème et version (type Bac)]
\textbf{A. Thème.}
\begin{enumerate}
\item \zh{虽然社交网络很方便，但是它们让年轻人浪费了很多时间。} \emph{Suīrán shèjiāo wǎngluò hěn fāngbiàn, dànshì tāmen ràng niánqīngrén làngfèi le hěn duō shíjiān.} (\zh{让} + personne + verbe = « faire faire à » ; \zh{浪费} = « gaspiller, perdre ».)
\item \zh{我哥哥开了一个微博账号，但是他忘了密码。} \emph{Wǒ gēge kāi le yí ge Wēibó zhànghào, dànshì tā wàng le mìmǎ.} (\zh{开} = « ouvrir » un compte ; classificateur \zh{个} devant \zh{账号}.)
\end{enumerate}
\textbf{B. Version.} \zh{我的粉丝不多，但是我每天还是在微博上发照片，因为我很喜欢分享我的生活。} → « Je n'ai pas beaucoup de fans, mais je publie quand même des photos sur Weibo tous les jours, parce que j'aime beaucoup partager ma vie. » \zh{还是} (háishi) = « quand même, malgré tout » ; \zh{分享} (fēnxiǎng) = « partager ».\par
\textbf{Barème indicatif (sur 10) :} thème 2 × 3 pts = 6 pts (1 pt vocabulaire, 1 pt grammaire, 1 pt caractères) ; version 4 pts (sens global 2 pts, précision des connecteurs 2 pts).\par
\textbf{Points de vigilance :} ne pas oublier \zh{但是} après \zh{虽然} ; bien écrire \zh{账} (clé 贝) et non \zh{帐} ; place du lieu \zh{在微博上} avant le verbe.
\end{corrige}

\begin{corrige}[Exercice 11 — Expression écrite (type Bac)]
\textbf{Proposition de rédaction modèle :}\par
\zh{我每天都用网络。首先，我早上在微信上跟朋友聊天，看他们发的照片。然后，我在微博上发视频。虽然我的粉丝不多，但是很多朋友喜欢我的视频，给我点赞。因为我要学习，所以晚上我不用手机。网络很方便，但是我们要小心：不要把密码告诉别人。}\par
\emph{Wǒ měitiān dōu yòng wǎngluò. Shǒuxiān, wǒ zǎoshang zài Wēixìn shàng gēn péngyou liáotiān, kàn tāmen fā de zhàopiàn. Ránhòu, wǒ zài Wēibó shàng fā shìpín. Suīrán wǒ de fěnsī bù duō, dànshì hěn duō péngyou xǐhuan wǒ de shìpín, gěi wǒ diǎnzàn. Yīnwèi wǒ yào xuéxí, suǒyǐ wǎnshang wǒ bú yòng shǒujī. Wǎngluò hěn fāngbiàn, dànshì wǒmen yào xiǎoxīn : bú yào bǎ mìmǎ gàosu biérén.}\par
« J'utilise Internet tous les jours. D'abord, le matin, je discute avec mes amis sur WeChat et je regarde les photos qu'ils publient. Ensuite, je publie des vidéos sur Weibo. Bien que je n'aie pas beaucoup de fans, beaucoup d'amis aiment mes vidéos et les likent. Comme je dois étudier, le soir je n'utilise pas mon téléphone. Internet est pratique, mais nous devons être prudents : il ne faut donner son mot de passe à personne. »\par
Ce modèle compte environ 100 caractères, 7 phrases, 4 connecteurs (\zh{首先}、\zh{然后}、\zh{虽然……但是……}、\zh{因为……所以……}) et 8 mots du thème (\zh{网络}、\zh{微信}、\zh{微博}、\zh{粉丝}、\zh{视频}、\zh{点赞}、\zh{密码}、\zh{手机}).\par
\textbf{Barème indicatif (sur 20) :} respect du sujet et des consignes (longueur, connecteurs, vocabulaire) : 5 pts ; exactitude des caractères : 5 pts ; grammaire et ordre des mots : 5 pts ; cohérence et ponctuation chinoise : 3 pts ; présentation générale : 2 pts.\par
\textbf{Points de vigilance :} \zh{微} et non \zh{徽} ; \zh{账} et non \zh{帐} ; compléments de temps, lieu, moyen et compagnie avant le verbe ; \zh{虽然} toujours suivi de \zh{但是} ; ponctuation chinoise ，。：uniquement.
\end{corrige}

\newpage

\coursec{Fiche bilan}

\begin{popfigure}
\begin{tikzpicture}[
  mindmap,
  grow cyclic,
  every node/.style={concept, execute at begin node=\hskip0pt},
  concept color=popblue,
  level 1/.append style={level distance=4.2cm, sibling angle=60},
  level 2/.append style={level distance=3cm, sibling angle=45},
  text=white,
  font=\small\bfseries,
  align=center
]
\node[concept, align=center]{Réseaux sociaux\\社交媒体\\(Shèjiāo méitǐ)}
  child[concept color=poporange] { node[concept, align=center]{Vocabulaire clé\\发布 点赞 转发\\评论 粉丝 直播} }
  child[concept color=popgreen] { node[concept, align=center]{Caractères complexes\\Radicales 扌 讠 纟 页\\Ordre des traits} }
  child[concept color=poppurple] { node[concept, align=center]{Caractères proches\\发/友 评/平 号/令\\辨析与书写} }
  child[concept color=poppink] { node[concept, align=center]{Modèle de texte\\Structure: 标题-正文-标签\\Connecteurs logiques} }
  child[concept color=popgold] { node[concept, align=center]{Traduction\\Thème / Version\\把 被 使 了 过} }
  child[concept color=popdark] { node[concept, align=center]{Évaluation Bac\\Grille: Contenu/Langue/Écriture\\Contexte Cameroun/Chine} };
\end{tikzpicture}
\legende{Carte mentale de la leçon ZH36 : Gestion des réseaux sociaux}
\end{popfigure}

\begin{bilan}

\textbf{1. Lexique clé (HSK 3-4) — Réseaux sociaux \& Gestion de compte}\par
\begin{center}
\begin{tabularx}{\linewidth}{|X|X|X|X|}
\hline
\rowcolor{popblueL} \textbf{Caractères} & \textbf{Pinyin} & \textbf{Nature} & \textbf{Traduction} \\
\hline
发布 & fābù & v. & publier, poster (contenu) \\
\hline
点赞 & diǎnzàn & v. & liker, « aimer » \\
\hline
转发 & zhuǎnfā & v. & partager, retweeter, transférer \\
\hline
评论 & pínglùn & v./n. & commenter / commentaire \\
\hline
粉丝 & fěnsī & n. & fans, abonnés (translit. « fans ») \\
\hline
关注 & guānzhù & v. & suivre, s'abonner, prêter attention \\
\hline
直播 & zhíbō & v./n. & diffuser en direct / live \\
\hline
私信 & sīxìn & v./n. & message privé (MP), DM \\
\hline
标签 & biāoqiān & n. & hashtag, étiquette, tag \\
\hline
热搜 & rèsōu & n. & sujet tendance, « hot search » \\
\hline
流量 & liúliàng & n. & trafic, audience, visibilité \\
\hline
算法 & suànfǎ & n. & algorithme \\
\hline
隐私 & yǐnsī & n. & vie privée, confidentialité \\
\hline
设置 & shèzhì & v./n. & paramétrer / paramètres \\
\hline
封号 & fēnghào & v. & bannir/suspendre un compte \\
\hline
\end{tabularx}
\end{center}

\textbf{2. Règles de grammaire à connaître par cœur (Structures HSK 3-4)}\par
\begin{center}
\begin{tabularx}{\linewidth}{|X|X|X|}
\hline
\rowcolor{popblueL} \textbf{Structure / Règle} & \textbf{Exemple (Caractères + Pinyin)} & \textbf{Traduction} \\
\hline
\textbf{把  bǎ  (objet connu + action + résultat)} & 我\zh{把}视频\zh{发布}了。\newline (Wǒ \zh{bǎ} shìpín \zh{fābù} le.) & J'ai publié la vidéo. (Action accomplie sur objet précis) \\
\hline
\textbf{被  bèi  (passif formel / malheur)} & 他的账号\zh{被}封号了。\newline (Tā de zhànghào \zh{bèi} fēnghào le.) & Son compte a été banni. \\
\hline
\textbf{使/让/叫  shǐ/ràng/jiào  (causatif)} & 这条热搜\zh{让}我\zh{关注}了他。\newline (Zhè tiáo rèsōu \zh{ràng} wǒ \zh{guānzhù} le tā.) & Ce sujet tendance m'a fait le suivre. \\
\hline
\textbf{连...都/也...  (insistance / surprise)} & 他\zh{连}私信\zh{都}不回。\newline (Tā \zh{lián} sīxìn \zh{dōu} bù huí.) & Il ne répond même pas aux MP. \\
\hline
\textbf{由于...因此/所以...  (cause-conséquence)} & \zh{由于}流量大,\zh{因此}直播卡顿。\newline (\zh{Yóuyú} liúliàng dà, \zh{yīncǐ} zhíbō kǎdùn.) & En raison du fort trafic, le direct lag. \\
\hline
\textbf{Aspect 了 (terminé) vs 过 (expérience)} & 我\zh{发过}这条动态。 vs 我\zh{发}了这条动态。\newline (Wǒ \zh{fāguo}... vs Wǒ \zh{fā} le...) & J'ai déjà posté ce statut (expérience). vs J'ai posté ce statut (action close). \\
\hline
\end{tabularx}
\end{center}

\textbf{3. Méthodes express}
\begin{itemize}
  \item \textbf{Rédiger un post / 微博/朋友圈文案} : 1) \textbf{Accroche} (标题/首句 : 疑问句, chiffre, émotion) → 2) \textbf{Contenu} (事实 + 连接词 : 首先, 此外, 最重要的是) → 3) \textbf{Appel à l'action} (欢迎评论, 点赞转发) → 4) \textbf{Hashtags} (\#话题\# @好友).
  \item \textbf{Thème (Fr $\rightarrow$ Zh)} : Identifier le COD $\rightarrow$ Choisir \zh{把} (action concrète) ou \zh{被} (subi) $\rightarrow$ Placer les compléments (Lieu/Temps/Manière) AVANT le verbe $\rightarrow$ Vérifier l'aspect (\zh{了}/\zh{过}) $\rightarrow$ Ajouter connecteurs.
  \item \textbf{Version (Zh $\rightarrow$ Fr)} : Segmenter (Sujet / Prédicat / COD) $\rightarrow$ Repérer mots-outils (\zh{了, 过, 把, 被, 连...都}) $\rightarrow$ Traduire le sens global (pas mot à mot) $\rightarrow$ Adapter le registre (familier « story » / formel « communiqué »).
  \item \textbf{Écrire les caractères complexes} : 1) Identifier le \textbf{radical} (sémantique) $\rightarrow$ 2) Compter les \textbf{traits} du corps $\rightarrow$ 3) Appliquer l'ordre : \emph{horizontal avant vertical, gauche avant droite, haut avant bas, centre avant ailes, fermer en dernier} $\rightarrow$ 4) Comparer aux \textbf{caractères proches} (ex: \zh{发} vs \zh{友}, \zh{评} vs \zh{平}).
\end{itemize}

\end{bilan}

\begin{aretenir}[Checklist avant le Bac]
$\square$ Je connais par cœur les 15 mots du tableau de vocabulaire (caractères, pinyin, sens, classe).\\
$\square$ J'écris correctement les 10 caractères complexes ciblés (ordre des traits, radicaux : \zh{发布 评论 转发 直播 算法 隐私 设置 封号 粉丝 标签}).\\
$\square$ Je distingue les paires de caractères proches : \zh{发/友}, \zh{评/平}, \zh{号/令}, \zh{设/订}, \zh{算/管}.\\
$\square$ Je maîtrise la structure \zh{把} (objet connu + résultat) pour décrire une action de gestion (poster, supprimer, paramétrer).\\
$\square$ Je maîtrise la structure \zh{被} pour exprimer le passif (compte banni, message supprimé, vie privée violée).\\
$\square$ J'utilise au moins 3 connecteurs logiques (\zh{此外, 因此, 不过, 总之, 首先}) pour structurer un texte de 80-100 caractères.\\
$\square$ Je sais traduire un thème « gestion de compte » (Fr $\rightarrow$ Zh) en respectant l'ordre des compléments (Temps/Lieu/Manière avant verbe).\\
$\square$ Je sais faire une version (Zh $\rightarrow$ Fr) fluide, en rendant les aspects \zh{了/过} et les structures \zh{把/被} par des temps/modes français corrects.\\
$\square$ Je connais 2 différences socioculturelles majeures (ex: \zh{微信} écosystème tout-en-un vs apps FR séparées ; paiement QR code ubiquitaire en Chine).\\
$\square$ Je gère mon temps : 15 min analyse/plan, 45 min rédaction, 15 min relecture (caractères, tons, ponctuation \zh{，。！？}).\\
$\square$ Je relis spécifiquement : radicaux oubliés, traits en trop/manquants, confusion \zh{的/得/地}, accord \zh{了/过}.
\end{aretenir}

\findecours

\end{document}
